\documentclass[11pt,reqno]{amsart}
\usepackage[T1]{fontenc}
\usepackage{lmodern,amsmath,amscd,amssymb,amsthm,mathtools,mathrsfs,microtype,booktabs,longtable}
\usepackage[margin=1.08in]{geometry}
\usepackage{xr-hyper}
\usepackage[hidelinks]{hyperref}
\numberwithin{equation}{section}
\newtheorem{theorem}{Theorem}[section]
\newtheorem{proposition}[theorem]{Proposition}
\newtheorem{lemma}[theorem]{Lemma}
\newtheorem{corollary}[theorem]{Corollary}
\theoremstyle{definition}\newtheorem{definition}[theorem]{Definition}
\newtheorem{example}[theorem]{Example}
\theoremstyle{remark}\newtheorem{remark}[theorem]{Remark}
\newcommand{\C}{\mathbb C}\newcommand{\Q}{\mathbb Q}
\newcommand{\Pj}{\mathbb P}\newcommand{\A}{\mathbb A}
\newcommand{\OO}{\mathcal O}\newcommand{\II}{\mathcal I}
\newcommand{\NN}{\mathcal N}\newcommand{\LL}{\mathcal L}
\newcommand{\ZZ}{\mathcal Z}\newcommand{\RR}{\mathcal R}
\DeclareMathOperator{\Sym}{Sym}\DeclareMathOperator{\Gr}{Gr}
\DeclareMathOperator{\Hom}{Hom}\DeclareMathOperator{\Tot}{Tot}
\DeclareMathOperator{\Fitt}{Fitt}\DeclareMathOperator{\Ann}{Ann}
\DeclareMathOperator{\rank}{rank}\DeclareMathOperator{\Proj}{Proj}
\DeclareMathOperator{\Spec}{Spec}\DeclareMathOperator{\PGL}{PGL}
\DeclareMathOperator{\Pic}{Pic}\DeclareMathOperator{\diag}{diag}
\DeclareMathOperator{\Tor}{Tor}\DeclareMathOperator{\im}{im}
\DeclareMathOperator{\coker}{coker}\DeclareMathOperator{\codim}{codim}
\DeclareMathOperator{\End}{End}
\DeclareMathOperator{\disc}{disc}\DeclareMathOperator{\Disc}{Disc}
\DeclareMathOperator{\Supp}{Supp}\DeclareMathOperator{\Hilb}{Hilb}
\DeclareMathOperator{\rad}{rad}\DeclareMathOperator{\Ass}{Ass}
\DeclareMathOperator{\gr}{gr}\DeclareMathOperator{\Gal}{Gal}
\DeclareMathOperator{\length}{length}
\allowdisplaybreaks[2]
\setlength{\emergencystretch}{3em}

\begin{document}
\title[Finite failure schemes and quadratic pencils]{Finite failure schemes and the reconstruction of quadratic pencils}
\author{Qian Qi}
\date{September 25, 2026}
\subjclass[2020]{14M12, 14B05, 14D20, 13C40}
\keywords{Quadratic pencil, finite failure algebra, intrinsic reconstruction, common divisor, conductor, normalization fibre}
\hypersetup{pdftitle={Finite failure schemes and the reconstruction of quadratic pencils, revision 154},pdfauthor={Qian Qi}}
\begin{abstract}
Every complex quadratic pencil is recovered from its unmarked,
ungraded finite multiplication-failure algebra at a sharp order.
We identify the intrinsic Pluecker polarization of this inverse and
its projective GIT completion. On the regular-pencil stack, a tower
of spectral Fitting divisors is realized as binary normalization
fibres; their successive local lengths recover the full Segre
symbol. We then study the common-divisor geometry of first relations.
The binary normalization is finite flat on each exact-gcd stratum,
with all collision fibres explicitly determined. On the squarefree
locus we compute the conductor, branch intersections and graded
Betti numbers. Arbitrary collisions have a completed incidence
presentation. In every embedding dimension the multivariate
pure-power fibres admit a determinant-apolar quotient, giving a
uniform length bound. The first multivariate two-factor fibre has
length 22, type 7 and an explicit weighted Betti table. These results
distinguish native pencil degeneration from the common extreme
limit under all cotangent coordinate changes.
\end{abstract}
\maketitle

\section{Introduction}
\label{sec:introduction-v153}
The nonreduced structure of a multiplication-failure scheme can
retain a multiplication law after its coordinates, grading and
tensor marking have been forgotten. For quadratic pencils this
retention is exact at one sharp infinitesimal order. The question
pursued here is not only whether a pencil can be recovered, but
which geometric constructions survive this recovery in families.
The first relation gives both a coefficient line and a divisor
incidence. We use the former to recover a polarized moduli problem
and the latter to describe its intrinsically induced spectral
fibres and its larger coordinate-change boundary.

Let $\dim V=n\geq3$, $N=\binom{n+1}{2}$, and
$R\in P=\Gr(2,\Sym^2V)$. Put $p=N-2$ and
$D=n+2p=n^2+2n-4$. If $T$ is the universal matrix and $\gamma_R$
the quotient by $R$, the relevant finite algebra is
\[
 \mathfrak A_R^{[D]}=
 \C[t_{ij}]/\bigl((\det T)I_p(\gamma_R\Sym^2T)
                                      +\mathfrak m^{D+1}\bigr).
\]
Theorems~\ref{thm:artin-local-inverse-v146} and
\ref{thm:sharp-finite-pencil} identify its ungraded, unmarked
isomorphism class with the congruence class of $R$. Every smaller
truncation is independent of $R$. The statement includes singular
pencils. The relative results recover moving pencil subbundles and
their actual source bundles, rather than merely fibrewise classes.

\subsection{A polarized inverse and its spectral boundary}
Write $M=\binom N2$. The coefficient line extracted from the first
relation is the Pluecker line of the pencil. Tensoring it with a
fixed $M$-space gives the closed first-relation parameterization
$\Phi$ of Theorem~\ref{thm:finite-parameter-immersion}.
Theorem~\ref{thm:polarized-completion-v154} identifies its
Pluecker polarization with $\OO_P(M)$, equivariantly for
$\operatorname{SL}(V)$. The intrinsic invariant therefore determines
the usual normal projective pencil quotient and a moduli morphism
from the effective semistable failure stack. This is a compatibility
and extension theorem for the failure invariant, not a new
construction of classical GIT or a properness assertion for the
whole stack.

There is a more discriminating invariant before passing to that
coarse quotient. On the recovered regular pencil line let $D_j$
be the divisor defined by the $j$th Fitting ideal of the spectral
torsion sheaf. Theorem~\ref{thm:spectral-fibres-v154} realizes $D_j$
as the full degree-one normalization fibre of the binary space
$G_jH^0(\OO(1))$, where $G_j$ cuts out $D_j$. If $b_{x,j}$ is its
local length at a spectral point, then the individual elementary
divisors are
\[
                    \lambda_{x,j+1}=b_{x,j}-b_{x,j+1}.
\]
Thus the fibre tower, with its common spectral support, distinguishes
Segre types having the same discriminant. On family strata where
the Fitting divisors are flat Cartier divisors, the construction
commutes with arbitrary base change. This supplies a precise link
between pencil degeneration data and divisor ramification.

\subsection{Normalization and the commutative algebra of its fibres}
For a vector space $E$ let $Y_g\subset\Gr(r,\Sym^D E)$ be the
image of $(f,J)\mapsto fJ$ from
$\Pj(\Sym^g E)\times\Gr(r,\Sym^{D-g}E)$. This finite map is its
normalization. For binary forms its restriction over the exact
common-divisor-degree-$s$ stratum is finite flat of degree
$\binom{s}{g}$. Theorem~\ref{thm:binary-fibres-v153} determines
all its scheme fibres at arbitrary root partitions as products
of reciprocal complete intersections.

On the squarefree-common-divisor open $\mathcal U$,
Theorem~\ref{thm:squarefree-v153} identifies the completed image
with a block coordinate-subspace ring and proves
\[
            \mathfrak c_{Y_g}|_{\mathcal U}
                  =\II_{Y_{g+1}/Y_g}|_{\mathcal U}.
\]
It also computes all multiple branch intersections, seminormality,
depth and multiplicity. Theorem~\ref{thm:block-betti-v154} gives
the entire graded Betti table of this model. The collision atlas
of Theorem~\ref{thm:atlas-v153} is an exact completed incidence
presentation; it is not advertised as a classification of all its
analytic singularities. The all-rank double-root singular Fitting
ideal is evaluated in Theorem~\ref{thm:fitting-v153}.

For $E=\C a\oplus W$, pure-power fibres have the reciprocal
presentation $F_{g,k}(W)$. Their full minimal relation representation
and complete-intersection criterion are given in
Theorem~\ref{thm:reciprocal-structure-v153}. We add a functorial
quotient by the symmetric determinant's apolar algebra, giving the
Catalan lower bound of Theorem~\ref{thm:apolar-quotient-v154} in
all dimensions when $k\geq2$. Proposition~\ref{prop:F22-v154}
gives the length, both grading-dependent Hilbert functions, type
and every weighted Betti number of the whole first multivariate
two-factor fibre.

\subsection{The two group actions and the comparisons}
The native action just used is $\PGL(V)$ congruence. By contrast,
Theorem~\ref{thm:universal-boundary-v152} places one explicit
ramified relation algebra in every full $\operatorname{GL}(E)$
first-relation orbit closure. Its construction uses two successive
specializations, not one asserted universal one-parameter subgroup.
The native first relation retains exact common factor
$(\det T)^{n-1}$ even when its pencil changes Segre type. The
spectral divisor tower is therefore not confused with that original
common-factor locus; Remark~\ref{rem:two-divisors-v154} explains
the distinction.

The classical mechanisms used here are attributed as such. The
binary reciprocal rings have the rectangular Schur presentation
of \cite{Grinberg2019}; the block calculations use the
Stanley--Reisner and Hochster machinery of
\cite{MillerSturmfels2005}; the apolar Hilbert function and
Gorenstein theorem are those of \cite{ShafieiSymmetric}.
Classical pencil and invariant-theoretic comparisons are
\cite{FevolaMandelshtamSturmfels,MFK1994}. The discriminant line
calculation is also discussed in \cite{Krishna2026}; the higher
Fitting divisors here retain the finer block data not recorded by
that discriminant. Further comparisons with coincident-root
loci, tuple common-factor spaces and conductor squares are in
Appendix~\ref{sec:introduction-v132}.

The full theorem and proof text of \cite{Ballico93} has not been
obtained. Consequently no theorem-level assertion that its
higher-order failure-locus results do or do not anticipate the
present inverse is made. This documentary limitation concerns
historical comparison, not the formulation of a weaker inverse.

\subsection{Reading structure}
Part I proves the sharp inverse and its polarized and spectral
moduli consequences. Part II develops the common-divisor image,
its binary collision fibres and its homological invariants, together
with the full-coordinate-change common limit. The complete relative,
curve, recognition, rigidification, covering and spectral arguments
remain in the appendices. The central dependency is
\emph{first relation--coefficient inverse--intrinsic spectral
line--divisor fibres}; the full-coordinate-change boundary is a
separate application of the same incidence theory.

\part*{Part I. Reconstruction and intrinsic moduli}

% BEGIN PRESERVED PART 33-first-relations-v147.tex
\section{First relations, ungraded isomorphisms, and inverse systems}
\label{sec:first-relations-v147}

The passage from a finite local algebra to its first relation space is
independent of pencils.  We isolate it before undertaking the geometric
inverse.  In this section \(E\) denotes a space of degree-one
\emph{generators}; thus it is a cotangent space, rather than the dual
normal space used later.  Let \(e=\dim E\), \(d\ge2\), and
\(K\subset\Sym^d E\).  Set
\begin{equation}\label{eq:general-first-algebra-v147}
 A(K)=\Sym E/\bigl((K)+\mathfrak m^{d+1}\bigr),
 \qquad \mathfrak m=(E),\qquad \mathfrak n=\mathfrak m A(K).
\end{equation}
The grading in this presentation is not supplied to an isomorphism.

\begin{proposition}[The first-relation orbit principle]
\label{prop:first-relation-orbit-v147}
The abstract algebra \(A(K)\) recovers \(E\) and \(K\), up to
linear isomorphism, by
\begin{equation}\label{eq:intrinsic-first-relation-v147}
 E=\mathfrak n/\mathfrak n^2,\qquad
 K=\ker\bigl[\Sym^d(\mathfrak n/\mathfrak n^2)
                                    \longrightarrow\mathfrak n^d\bigr].
\end{equation}
Consequently \(A(K)\simeq A(K')\) as ungraded complex algebras
if and only if some \(a:E\xrightarrow{\sim}E'\) satisfies
\((\Sym^d a)K=K'\).

More precisely, after fixing the homogeneous inclusions of generators,
all isomorphisms with tangent map \(a\) are exactly the substitutions
\[
                 x\longmapsto a(x)+u(x),\qquad
                 u\in\Hom_{\C}(E,\mathfrak n'^2).
\]
These assertions about substitutions remain valid after scalar
extension to any commutative complex algebra.
\end{proposition}
\begin{proof}
The positive ideal is the unique maximal ideal, and
\(\mathfrak n^{d+1}=0\).  Multiplication of \(d\) lifts of
cotangent vectors is independent of their lifts: replacing any one
by an element of \(\mathfrak n^2\) changes the product in
\(\mathfrak n^{d+1}\).  There are no relations of smaller degree,
and the degree-\(d\) quotient is \(\Sym^d E/K\).  This proves
\eqref{eq:intrinsic-first-relation-v147}.  In particular every
isomorphism induces the stated linear equivalence.

Conversely substitute \(a(x)+u(x)\) for each generator.  A relation
of degree \(d\) changes, beyond its linear substitution, only in
degrees at least \(d+1\).  Those terms vanish.  The relations of
degree \(d+1\) also vanish.  The induced map on associated graded
algebras is the isomorphism induced by \(a\); a finite filtration
then makes the algebra map invertible.  Every map with this tangent
map has precisely such generator images.  These polynomial identities
hold over any commutative complex coefficient algebra.  The intrinsic
preservation of the augmentation under that scalar extension is
verified in the next proof, so no preservation of a chosen origin has
tacitly been added to the isomorphism problem.
\end{proof}

\begin{proposition}[The full automorphism group scheme]
\label{prop:first-relation-group-v147}
Write \(\ell=\dim_{\C}A(K)=\binom{e+d}{d}-\dim K\) and
\(H_K=\operatorname{Stab}_{\operatorname{GL}(E)}(K)\), with its
closed subgroup scheme structure.  The functor
\(D\mapsto\operatorname{Aut}_{D}(A(K)\otimes_{\C}D)\) is an
affine group scheme of finite type and fits into a split exact sequence
of affine group schemes over \(\C\):
\begin{equation}\label{eq:first-group-v147}
 1\longrightarrow U_K\longrightarrow\operatorname{Aut}(A(K))
       \longrightarrow H_K\longrightarrow1.
\end{equation}
The kernel is unipotent.  Choosing the homogeneous lifting of \(E\)
identifies its underlying scheme, but not generally its group law,
with the affine space associated to \(\Hom_{\C}(E,\mathfrak n^2)\).
Its dimension is \(e(\ell-1-e)\).  Intrinsically, \(U_K\) acts
simply transitively on the affine space of linear splittings of
\(\mathfrak n\twoheadrightarrow\mathfrak n/\mathfrak n^2\).
The splitting of \eqref{eq:first-group-v147} uses a choice of such a
lifting; it is not an additional intrinsic grading.
\end{proposition}
\begin{proof}
Choose a basis of the finite-dimensional vector space \(A=A(K)\).
Inside \(\operatorname{GL}(A)\), the equations
\(f(1)=1\) and \(f(xy)=f(x)f(y)\), imposed on basis elements,
are polynomial equations.  They represent the automorphism functor,
including its values on nonreduced complex algebras.

The augmentation can be recovered without mistaking the absolute
nilradical after scalar extension for the augmentation ideal.  If
\(M_z\) denotes multiplication by \(z\), the homogeneous filtration
gives
\begin{equation}\label{eq:trace-augmentation-v147}
                      \operatorname{tr}(M_z)=\ell\,\epsilon(z).
\end{equation}
Indeed multiplication by a positive-degree element is strictly
triangular, whereas multiplication by a scalar has that scalar on
every diagonal entry.  Identity~\eqref{eq:trace-augmentation-v147}
holds over every commutative complex algebra \(D\).  Algebra
automorphisms conjugate multiplication matrices, so they preserve
\(\epsilon=\ell^{-1}\operatorname{tr}(M_{(-)})\), its kernel
\(\mathfrak n_D\), and all its powers.  Thus the tangent map is
a morphism of group schemes into \(H_K\), not merely a map on
complex points.

Linear substitutions give a morphism \(H_K\to\operatorname{Aut}(A)\)
sectioning this tangent map.  By Proposition~\ref{prop:first-relation-orbit-v147},
its kernel has, functorially in \(D\), the unique parametrization
\(x\mapsto x+u(x)\) with \(u:E_D\to\mathfrak n_D^2\).
Reading generator images and making these substitutions are inverse
polynomial maps.  On a basis ordered by the positive filtration all
kernel matrices are unitriangular, proving unipotence.  Multiplication
is substitution, which need not be addition of corrections.  The
dimension follows from \(\dim\mathfrak n^2=\ell-1-e\).

Finally two splittings of \(\mathfrak n\twoheadrightarrow E\)
differ by exactly a map \(E\to\mathfrak n^2\).  Any splitting
generates \(A\), and its degree-\(d\) products impose exactly
\(K\); the preceding substitution argument therefore gives the
unique element of \(U_K\) relating any two splittings.  This is
the asserted coordinate-free torsor description.
\end{proof}

\subsection{The inverse-system description and its precise role}
Put \(E^{\vee}=\Hom_{\C}(E,\C)\).  Use the apolar pairing in
which dual divided monomials satisfy
\(\langle x^{\alpha},y^{\beta}/\beta!\rangle=\delta_{\alpha\beta}\).
Here the pairing means evaluation at the origin after applying the
constant-coefficient differential operator.  For
\(I=(K)+\mathfrak m^{d+1}\), its Macaulay inverse system is
\begin{equation}\label{eq:inverse-system-v147}
 I^{\perp}=\bigoplus_{j<d}\Sym^j E^{\vee}\ \oplus\ K^{\perp},
              \qquad K^{\perp}\subset\Sym^d E^{\vee}.
\end{equation}
To check this equality, the relations of degree \(d+1\) bound the
degree of every inverse-system element by \(d\).  The degree-\(d\)
relations pair only with its top component.  All lower components are
unrestricted.  Conversely that direct sum is stable under every
partial derivative: differentiating its top part lands in the full
lower-degree space.  It is therefore annihilated by \(I\), proving
the displayed equality.  Thus the orbit problem for \(K\) is
equivalently the contragredient orbit problem for \(K^{\perp}\).
This uses the classical inverse-system correspondence; it does not
assert a new canonical-grading theorem.

Elias and Rossi prove that Gorenstein local algebras with Hilbert
function \((1,e,e,1)\) are canonically graded
\cite[Theorem~3.3 and Corollaries~3.4--3.5]{EliasRossiShort}.
Their proof removes a lower-degree term of the inverse-system
generator by a nonlinear coordinate change.  In their later work,
compressed Gorenstein algebras of socle degree at most four are
canonically graded, with counterexamples beyond the stated range
\cite[Theorem~3.1 and Examples~3.3--3.5]{EliasRossiCompressed}.
In contrast, \eqref{eq:general-first-algebra-v147} is homogeneous by
construction, and its truncation makes all higher substitutions
invisible to its first relation.  We do not deduce that an arbitrary
filtered deformation with the same Hilbert function is graded.

For the pencil algebras, \(e=n^2\), \(d=n^2+2n-4\ge11\), and
\(\dim K=\binom N2\).  They are not Gorenstein: their socles contain
\(\mathfrak n^d\), whose dimension is
\(\binom{e+d-1}{d}-\binom N2>1\).  For example, the first binomial
is at least \(\binom{e+1}{2}\), while \(e=n^2>N\).  The issue
left by Proposition~\ref{prop:first-relation-orbit-v147} is accordingly
not the removal of a nonhomogeneous inverse-system term.  It is the
recovery of the matrix factors and defining coefficient data from the
\emph{unrestricted} \(\operatorname{GL}(E)\)-orbit of \(K\).
The next sections establish that recovery and its global descent.

% BEGIN PRESERVED PART 00-ruling-foundations-v145.tex
\section{Intrinsic tensor rulings}
\label{sec:ruling-foundations-v143}
The first step of the inverse problem is to recognize the two tensor
factors of a determinant cone without an ambient marking.  The following
scheme-theoretic ruling calculation and descent lemma supply precisely
this step.  They apply to the pencil argument independently of the web
and exterior-contraction constructions in the separate research archive.

\begin{lemma}[The relative scheme of maximal rank-one spaces]
\label{lem:rank-one-fano-scheme}
Let \(B\) be any complex scheme and \(E,F\) vector bundles of
rank \(n\ge2\).  In \(\Gr_B(n,E\otimes F)\), let
\(Z(E,F)\) be the zero scheme of the restriction of the quadratic
rank-one equations to the universal subbundle \(\mathcal T\).
Equivalently it is the zero scheme of
\[
 \Sym^2\mathcal T\longrightarrow
              \bigwedge^2E\otimes\bigwedge^2F.
\]
The natural morphism
\begin{equation}\label{eq:fano-rank-one-scheme}
 \Pj_B(E)\ \sqcup\ \Pj_B(F)\longrightarrow Z(E,F),
 \qquad \ell\mapsto\ell\otimes F,
 \quad m\mapsto E\otimes m,
\end{equation}
is an isomorphism.  It commutes with arbitrary change of complex
base, including change to a nonreduced base.
\end{lemma}
\begin{proof}
We first work over \(\C\).  Each displayed map to the
Grassmannian is a closed immersion.  Its composition with the
Pluecker embedding is the degree-\(n\) Veronese embedding on
\(\ell^{\otimes n}\otimes\det F\), respectively
\(\det E\otimes m^{\otimes n}\), followed by its natural linear
inclusion into \(\bigwedge^n(E\otimes F)\).  This linear inclusion
is injective: polarization of the nonzero pure-power map gives the
corresponding irreducible symmetric-power summand.  The rank-one
equations vanish identically on both universal families, so the
closed immersions factor through \(Z(E,F)\).

Every geometric point is in one of the two images.  Indeed the sum
of two rank-one tensors is rank at most one only if their first
factors or their second factors are proportional.  Starting with two
independent tensors, this observation forces a linear space consisting
of rank-one tensors to fix the same factor throughout.  A space of
dimension \(n\) is then exactly one of the displayed spaces.
The two images are disjoint: a space fixing both factors has dimension
one, not \(n\).

We verify that these pointwise statements omit no infinitesimal
structure.  At \(\ell\otimes F\), choose a generator of \(\ell\)
and a complement in \(E\).  A Grassmannian tangent vector is a map
\(\phi:F\to(E/\ell)\otimes F\).  For each coordinate of
\(E/\ell\), write its component as \(A:F\to F\).
Linearizing the restricted minors gives
\[
                       u\wedge A(u)=0\quad\text{for every }u\in F.
\]
Polarizing this diagonal identity gives
\[
               u\wedge A(v)+v\wedge A(u)=0\quad(u,v\in F).
\]
In characteristic zero the diagonal and polarized identities are
equivalent.  These are precisely all coefficients of the bilinear
linearization of the restricted quadratic map on \(\Sym^2F\),
so no infinitesimal equation is omitted by testing the diagonal.
Taking the basis vectors shows that \(A\) is diagonal; taking their
pairwise sums makes all its diagonal entries equal.  Conversely a
scalar \(A\) satisfies all the linearized equations.  Thus the
Zariski tangent space to \(Z(E,F)\) is precisely
\(\Hom(\ell,E/\ell)\), of dimension \(n-1\).
The other image has the symmetric calculation.

At every closed point the closed immersed projective space gives
local dimension at least \(n-1\), whereas the tangent calculation
gives embedding dimension exactly \(n-1\).  The local ring is
therefore regular.  A finite-type complex scheme whose closed-point
local rings are regular is reduced: a nonempty support of its
nilradical would contain a closed point.  In particular there is no
nilpotent thickening or embedded associated structure at either
image.  The disjoint union of the two closed immersions is itself a
closed immersion, surjective on geometric points.  Its ideal in the
now reduced target vanishes at every closed point, so it is zero.
This proves the scheme isomorphism over \(\C\).

For general \(B\), trivialize \(E,F\) on a common open cover.
The Grassmannian and the zero scheme of the displayed universal
quadratic map are, on such an open, the base change of the constant
construction just proved.  Thus \eqref{eq:fano-rank-one-scheme} is
an isomorphism there even if the base has nilpotents.  The maps are
natural under changes of frame and hence glue.  This argument is
an identification by equations and base change, not an inference of
relative flatness from a fibrewise tangent calculation.
\end{proof}


\begin{lemma}[Oriented relative Segre descent]
\label{lem:oriented-segre-descent}
Let \(B\) be a connected projective complex variety with
\(H^0(B,\OO_B)=\C\), and let \(V,V'\) have the same dimension
\(n\ge2\).  Let \(\mathcal U,\mathcal U'\) be vector bundles
of rank \(n\) on \(B\).  Suppose a linear bundle isomorphism
\[
 \alpha:V\otimes\mathcal U^*\xrightarrow{\sim}
                         V'\otimes\mathcal U'^*
\]
preserves the rank-one cones and the ruling whose parameter bundle
is \(\Pj(V)\times B\).  The induced isomorphism
\(\Pj(V)\times B\simeq\Pj(V')\times B\) is regular and is
induced by one constant projective isomorphism \([g]:\Pj(V)\to\Pj(V')\).
Locally \(\alpha(T)=gTh^{-1}\), with the same \([g]\) on every
chart.  The possible line-bundle twists of vector-bundle lifts do
not change the projective left coefficient lines of any homogeneous
functorial coefficient construction.
\end{lemma}
\begin{proof}
Lemma~\ref{lem:rank-one-fano-scheme}, applied to
\(E=V\otimes\OO_B\) and \(F=\mathcal U^*\), identifies the
closed relative Fano scheme with
\((\Pj(V)\times B)\sqcup\Pj_B(\mathcal U^*)\).
In particular its two components carry their displayed scheme structures,
with no additional nilpotent or embedded components.
The isomorphism \(\alpha\) induces a regular map of these relative
parameter schemes, and the orientation hypothesis preserves the
first component.

For completeness, an isomorphism of two trivial \(\Pj^{n-1}\)-bundles
over any base gives a regular projective-matrix-valued map, not
merely a fibrewise collection.  Its pullback of the relative
\(\OO(1)\) has the form \(\OO(1)\otimes p^*L\) for a line
bundle \(L\) on the base.  To see this, tensor by \(\OO(-1)\);
the resulting line bundle is trivial on every projective-space
fibre, and its pushforward and evaluation identify it with a
pullback from the base.  Pushing forward \(\OO(1)\) now gives
an isomorphism of the corresponding trivial rank-\(n\) bundles,
with that line-bundle twist.  On charts trivializing \(L\), it
is an invertible matrix \(g_i\) of regular functions.  On overlaps,
\(g_j=u_{ij}g_i\) for an invertible regular function \(u_{ij}\).
The classes \([g_i]\) thus glue to a regular morphism to the
projective-isomorphism scheme, identified after choosing fixed
bases with \(\PGL_n\).  This accounts explicitly for the possible
absence of a global linear lift.

The group \(\PGL_n\) is affine: it is the determinant-nonzero
principal open in the projective space of matrices.  A morphism
from \(B\) to an affine scheme factors through
\(\Spec H^0(B,\OO_B)=\Spec\C\), so the projective map is constant.
Choose a constant linear lift \(g\).  Each local \(g_i\) then
differs from \(g\) by a scalar unit.  After removing its two
projective-factor actions, a linear map preserving every
decomposable line is scalar: apply it to sums of two tensors with
one common factor.  Consequently the remaining action is a right
bundle map, and the scalar units may be absorbed in its local
lifts \(h\).  This proves the local tensor formula.

Finally a scalar multiplying a lift acts on a homogeneous functor
of degree \(a\) by its scalar to the power \(a\) (or \(-a\)
on its dual).  It cannot alter a projective coefficient line.
The conclusion applies to every homogeneous degree.  In particular,
no right change of frame can replace the single projective left map
by separate transformations on different coefficient spaces.
\end{proof}

% BEGIN PRESERVED PART 34-coefficient-descent-v147.tex
\section{Actual factor descent and moving coefficient spaces}
\label{sec:coefficient-descent-v147}

We now separate the linear-preserver input, the descent of vector
bundles, and the recovery of coefficients.  The resulting principle
allows several varying coefficient spaces, of arbitrary prescribed
ranks, rather than only the line associated with a pencil.

\begin{lemma}[Linear preservers of the determinant cone]
\label{lem:linear-preserver-v147}
Let \(U,V,U',V'\) be complex vector spaces of dimension \(n\ge2\).
An invertible linear map
\(\Phi:\Hom(U,V)\to\Hom(U',V')\) carrying the determinant
hypersurface onto the determinant hypersurface has one of the forms
\[
 T\longmapsto gTh^{-1},\qquad
 T\longmapsto gT^{\mathsf t}h^{-1},
\]
where the second formula is written after choosing identifications
for transposition.  The two cases respectively preserve and exchange
the two rank-one rulings.  In the first case the pair \((g,h)\)
is unique up to a common nonzero scalar.
\end{lemma}
\begin{proof}
At a matrix of rank \(r\), an invertible \(r\)-minor reduces the
local rank condition to the rank condition on a generic
\((n-r)\)-square matrix.  The singular locus of the rank-at-most
\(s\) variety is therefore the rank-at-most \(s-1\) variety for
\(1\le s<n\): on the rank-\(s\) stratum the Schur-complement
entries give independent equations, whereas on a lower stratum the
linear terms of its defining \((s+1)\)-minors vanish.  Starting
with the determinant cone and iterating reduced singular loci hence
recovers the rank-one cone.  For \(n=2\) that cone is already the
determinant cone.

If \(v\otimes u\) and \(w\otimes z\) have rank one, their sum
has rank at most one only if \(v,w\) or \(u,z\) are dependent.
It follows that each maximal linear space of rank-one tensors has
one factor fixed.  The two families are
\(\Pj(V)\) and \(\Pj(U^*)\).  The induced Segre automorphism
must preserve these families or exchange them.  In the preserving
case its restriction to a ruling gives projective linear maps on
each factor.  Remove linear lifts of both maps.  The remaining
linear transformation fixes each decomposable line.  In tensor bases
its eigenvalues at \(v_i\otimes u_j\) agree along each row and
each column by applying it to
\((v_i+v_k)\otimes u_j\) and \(v_i\otimes(u_j+u_l)\).
They are all equal, so the remaining map is scalar.  This proves
the first form and its uniqueness.  Composing with transposition
reduces the exchanging case to the first.

This is the square invertible case of the classical rank-one
preserver theorem of Marcus and Moyls
\cite[Theorem~1, pp.~1218--1219]{MarcusMoyls}; the argument above
specifies the reduction from the determinant hypersurface that is
needed here.
\end{proof}

\begin{lemma}[An actual tensor map removes the line twist]
\label{lem:actual-factor-descent-v147}
Let \(B\) be a reduced complex variety, \(V\) a complex vector
space with \(\dim V\ge2\), and \(\mathcal A_1,\mathcal A_2\)
vector bundles.  Suppose an actual bundle isomorphism
\[
 \Psi:V\otimes\mathcal A_1\xrightarrow{\sim}V\otimes\mathcal A_2
\]
preserves the rank-one cones and induces the identity on the
\(\Pj(V)\)-parameter space of rulings.  There is a unique bundle
isomorphism \(a:\mathcal A_1\to\mathcal A_2\) with
\(\Psi=1_V\otimes a\).

Equivalently, the intrinsic operation of taking the commutant of the
constant matrix algebra gives
\begin{equation}\label{eq:commutant-descent-v147}
 \mathcal Hom_{\End(V)\otimes\OO_B}
       (V\otimes\mathcal A_1,V\otimes\mathcal A_2)
             \simeq\mathcal Hom_{\OO_B}(\mathcal A_1,\mathcal A_2).
\end{equation}
Thus after a constant left projective map has been lifted linearly,
no residual line-bundle twist is left in the recovery of the right
factor from \(\Psi\).
\end{lemma}
\begin{proof}
For every constant line \(\C v\subset V\), the corresponding
ruling is fixed, so on every geometric fibre \(\Psi\) maps
\(v\otimes\mathcal A_1\) onto \(v\otimes\mathcal A_2\).
Choose a basis \(v_1,\ldots,v_n\) for the purpose of checking
this assertion as an equality of bundle maps.  The rulings for
\(v_i\) force all off-diagonal blocks of \(\Psi\) to vanish.
The rulings for \(v_i+v_j\) force its diagonal blocks to agree.
Since \(B\) is reduced, fibrewise zero morphisms of vector bundles
are zero.  Therefore \(\Psi\) commutes with every constant
matrix unit, hence with \(\End(V)\otimes\OO_B\).

For clarity, commutant recovery itself is canonical and does not
use these chosen matrix units as extra data.  Tensoring a bundle
map \(a\) with \(1_V\) gives the forward map in
\eqref{eq:commutant-descent-v147}.  Commutation with matrix units
shows locally that it is bijective, with inverse reading any diagonal
block.  This proves that it is an isomorphism of sheaves, independent
of the checking basis.  Its inverse applied to the global section
\(\Psi\) produces a global section \(a\), not projective local
sections requiring arbitrary scalar lifts.  Applying the same
operation to \(\Psi^{-1}\) proves that \(a\) is invertible.
Uniqueness follows from injectivity of \(a\mapsto1_V\otimes a\).

Projective-factor data alone would permit replacing a bundle by a
line twist.  The additional datum used here is the actual linear map
of the tensor bundles.  It is this datum, recovered from the dual of
the first nilradical quotient, that makes the commutant argument
applicable.
\end{proof}

\subsection{A general finite-order coefficient inverse}
Fix \(n\ge2\), \(q\ge1\), and integers
\(0\le r_\lambda\le m_\lambda=\dim\mathbb S_\lambda\C^n\)
for partitions \(\lambda\) of \(q\) of length at most \(n\).
Assume that for at least one partition
\begin{equation}\label{eq:strict-support-v147}
                         0<r_\lambda<m_\lambda.
\end{equation}
Let \(B\) be a smooth connected projective complex variety,
\(V\) a constant \(n\)-dimensional space, and \(\mathcal U\)
a rank-\(n\) bundle.  For each \(\lambda\), choose a subbundle
with locally free quotient
\[
 \mathcal L_\lambda\subset\mathbb S_\lambda V^*\otimes\OO_B,
                       \qquad\rank\mathcal L_\lambda=r_\lambda.
\]
Zero and full subbundles are permitted.  In the matrix bundle
\(\mathcal E=\Hom(\mathcal U,V\otimes\OO_B)\), the
multiplicity-one Cauchy decomposition supplies the subbundle
\begin{equation}\label{eq:general-moving-coeff-v147}
 \mathcal J_q=\bigoplus_\lambda
       \mathcal L_\lambda\otimes\mathbb S_\lambda\mathcal U
             \ \subset\ \Sym^q\mathcal E^*.
\end{equation}
Let \(\mathcal D_n\subset\Sym^n\mathcal E^*\) be the
determinant line, put \(D=n+q\), and define
\begin{equation}\label{eq:coefficient-thickening-v147}
 \mathcal K=\mathcal D_n\mathcal J_q,
 \quad \mathcal A=\Sym\mathcal E^*/\bigl((\mathcal K)
                                  +\mathfrak a^{D+1}\bigr),
 \quad Y=\Spec_B\mathcal A,
\end{equation}
where \(\mathfrak a\) is the positive homogeneous ideal.  The
notation \(D\) in this subsection is an integer, not a failure
locus.  The scheme \(Y\) is finite over \(B\), and need not be
zero-dimensional over \(\C\).

\begin{theorem}[Unmarked descent for moving coefficient spaces]
\label{thm:general-moving-coefficients-v147}
For fixed \(n,q,(r_\lambda)\) satisfying
\eqref{eq:strict-support-v147}, the abstract complex scheme \(Y\)
recovers its coefficient data as follows.  Two such schemes \(Y_1,Y_2\)
are isomorphic if and only if there are an isomorphism
\(\sigma:B_1\to B_2\), a constant linear isomorphism
\(g:V_1\to V_2\), and an actual bundle isomorphism
\(h:\mathcal U_1\to\sigma^*\mathcal U_2\) such that, for every
\(\lambda\),
\begin{equation}\label{eq:coefficient-equivalence-v147}
 \mathcal L_{1,\lambda}=
       (\mathbb S_\lambda g)^*\sigma^*\mathcal L_{2,\lambda}.
\end{equation}
Every unmarked isomorphism induces a normal-bundle map
\(T\mapsto gTh^{-1}\); for that map the pair \((g,h)\) is
unique up to a common nonzero constant scalar.  Conversely every
such triple lifts homogeneously to an isomorphism.  The scheme is
finite locally free over its reduced base of rank
\[
             \binom{n^2+D}{D}-\sum_\lambda r_\lambda m_\lambda.
\]
No nontriviality hypothesis on \(\Pj_B(\mathcal U^*)\) is required.
\end{theorem}
\begin{proof}
\emph{Recovering the intrinsic linear data.}
Multiplication by a determinant polynomial is an injective linear
map of polynomial spaces.  In local bundle frames its matrix is
constant over \(\C\) and its cokernel is free.  The Cauchy
summands and the \(\mathcal L_\lambda\) are subbundles, so
\(\mathcal K\) is a subbundle of rank
\(\sum r_\lambda m_\lambda\).  The graded pieces of
\(\mathcal A\) below \(D\) are the full symmetric powers;
its top piece is \(\Sym^D\mathcal E^*/\mathcal K\).
This proves local freeness and the rank formula.

The nilradical \(\mathfrak n\) of \(\OO_Y\) is its positive
ideal, since \(B\) is reduced.  Thus \(Y_{\mathrm{red}}=B\) and
\(\mathfrak n/\mathfrak n^2=\mathcal E^*\) intrinsically.
The multiplication kernel
\[
 \ker\bigl[\Sym^D(\mathfrak n/\mathfrak n^2)
                                    \longrightarrow\mathfrak n^D\bigr]
\]
is precisely \(\mathcal K\), by the proof of
Proposition~\ref{prop:first-relation-orbit-v147} in local frames.
An abstract isomorphism consequently gives the base isomorphism
\(\sigma\) and an actual linear normal-bundle isomorphism carrying
\(\mathcal K_1\) to \(\mathcal K_2\) contravariantly.  The
normal map is linear over the reduced base even when the original
scheme map does not preserve the chosen graph projection.  Indeed
the action on \(\mathfrak n/\mathfrak n^2\) is linear over
\(\OO_Y/\mathfrak n\).

\emph{Recovering and orienting the factors.}
Generate the homogeneous ideal \((\mathcal K)\) in the recovered
symmetric algebra; do not confuse its cone with the finite
truncation.  Its radical is the determinant ideal.  To verify this,
at an invertible matrix \(T\), the map
\(\mathbb S_\lambda T\) is invertible.  For a partition with
\(r_\lambda>0\), a nonzero covector in \(\mathcal L_\lambda\)
has a nonzero matrix coefficient there.  Thus the residual
coefficients cannot all vanish off the determinant.  All generators
vanish on the determinant.  The relative determinant hypersurface
is reduced over the smooth base; the equality of radicals follows
on affine charts from this equality of closed points.  Successive
reduced singular loci and Lemma~\ref{lem:rank-one-fano-scheme}
recover its two unordered tensor rulings with their relative scheme
structures.

Cancellation of the determinant line is intrinsic and recovers the
ideal generated by \(\mathcal J_q\), and hence its degree-\(q\)
piece.  In a Cauchy summand the two contraction supports of
\(\mathcal L_\lambda\otimes\mathbb S_\lambda\mathcal U\)
have ranks \((r_\lambda,m_\lambda)\).  Separate factor changes
preserve these ranks.  By Lemma~\ref{lem:linear-preserver-v147},
an exchange of the two matrix factors interchanges the ranks in
each summand.  At a partition satisfying
\eqref{eq:strict-support-v147}, the result
\((m_\lambda,r_\lambda)\) cannot equal
\((r_\lambda,m_\lambda)\) for the other datum.  Thus no fibre
permits an exchange.  Twisting local factor lifts by lines changes
no support rank.  The \(V\)-ruling is thereby oriented intrinsically,
including when both projective ruling bundles are trivial.

\emph{Descending the actual map.}
By Lemma~\ref{lem:oriented-segre-descent}, the induced morphism
\(B_1\to\operatorname{Isom}(\Pj(V_1),\Pj(V_2))\) is constant:
the target is affine and \(H^0(B_1,\OO_{B_1})=\C\).  Choose a
constant linear lift \(g\).  Remove it from the actual normal map.
Lemma~\ref{lem:actual-factor-descent-v147}, applied to
\(\mathcal A_1=\mathcal U_1^*\) and
\(\mathcal A_2=(\sigma^*\mathcal U_2)^*\), supplies the unique
global right factor.  Dualizing its inverse gives the actual
isomorphism \(h\).  This step uses more than an isomorphism of
projective bundles; it is exactly what excludes an undetected line
twist.  Replacing \(g\) by \(cg\) replaces \(h\) by \(ch\).
These are all ambiguities because the left lift is constant.

\emph{Recovering every coefficient subbundle.}
The oriented Cauchy decomposition is now functorial for the recovered
left-right map.  Project the residual degree-\(q\) module onto each
summand and contract against the full right dual factor.  The image
is precisely \(\mathcal L_\lambda\), as a subbundle of the
constant left module, not just a collection of fibrewise projective
subspaces.  The right map acts invertibly on the full Schur factor
and hence does not change this image.  Pullback by the normal map
therefore gives \eqref{eq:coefficient-equivalence-v147} for the same
constant \(g\) in all summands.  Scalar choices in the determinant
line do not change these subbundles.  This proves necessity.

Conversely \(T\mapsto gTh^{-1}\) carries the determinant line,
all coefficient subbundles, and every power of the positive ideal
to the corresponding objects.  It induces the required homogeneous
scheme isomorphism.  The stated uniqueness refers to the induced
normal map; it does not eliminate the higher-order automorphisms of
that scheme.
\end{proof}

\subsection{Identical fibres and graded bundles do not determine a family}
The global conclusion has content even when neither fibre isomorphism
classes nor the underlying graded vector bundles vary.  The following
example uses the general coefficient theorem, not a further spectral
readout of a reconstructed pencil.

\begin{example}[Two isotrivial thickenings separated by coefficient descent]
\label{ex:isotrivial-covers-v147}
Take \(B=\Pj^1\), \(V=\C^3\), \(\mathcal U=\OO_B^3\),
\(q=1\), and \(r_{(1)}=1\).  The two everywhere nonzero columns
\[
 (s^3,t^3,0)^{\mathsf t},\qquad
 (s^3+st^2,t^3,0)^{\mathsf t}
\]
define subbundles \(\mathcal L_0,\mathcal L_1\subset V^*\otimes\OO_B\),
both isomorphic to \(\OO_B(-3)\).  Let \(Y_0,Y_1\) be their
order-\(4\) thickenings \eqref{eq:coefficient-thickening-v147}.
They are finite flat of rank \(712\) over \(B\) and Zariski
locally trivial with the same Artin algebra as fibre.  Their
nilradical graded vector bundles are isomorphic in every degree,
although \(Y_0\) and \(Y_1\) are not isomorphic as abstract
complex schemes.
\end{example}
\begin{proof}
Neither column vanishes, since its second entry is nonzero when
\(t\ne0\), and its first is nonzero at \(t=0\).  The left
coefficient map to its image line \(\Pj^1\subset\Pj(V^*)\)
is, in the coordinate \(z=s/t\), respectively
\[
                        f_0(z)=z^3,\qquad f_1(z)=z^3+z.
\]
A local frame completing a generator of \(\mathcal L_i\) to a
basis of \(V^*\otimes\OO\) changes the coefficient line to a
fixed line.  This gives Zariski local product presentations of both
thickenings.  Alternatively, equality of the fibre isomorphism
classes follows from transitivity of \(\operatorname{GL}(V)\)
on nonzero coefficient lines.  The rank formula is
\(\binom{13}{4}-3=712\).

Let \(E=V^*\otimes\C^3\), of dimension nine.  For \(j<4\),
the degree-\(j\) nilradical quotient is \(\Sym^j E\otimes\OO_B\).
Multiplication by the fixed determinant embeds \(E\) as a fixed
nine-dimensional subspace of \(\Sym^4 E\).  Choosing a constant
linear complement gives, for the top quotient,
\[
 \mathfrak n_i^4\simeq
 \OO_B^{486}\oplus
           \bigl((V^*\otimes\OO_B)/\mathcal L_i\bigr)^{\oplus3}.
\]
Each column has only its first two entries nonzero, and these two
entries have no common zero.  The rank-one quotient of \(\OO_B^2\)
by \(\OO_B(-3)\) is consequently \(\OO_B(3)\), by determinants.
Thus the displayed top bundle is
\(\OO_B^{489}\oplus\OO_B(3)^{\oplus3}\) for both \(i\).
The assertions about every graded vector bundle follow.  They are
not assertions that the graded \emph{algebra} structures agree.

If an abstract isomorphism \(Y_0\simeq Y_1\) existed,
Theorem~\ref{thm:general-moving-coefficients-v147} would identify
their coefficient maps by one automorphism of \(B\) and one
constant projective transformation of \(V^*\).  Both maps have
as image the same projective line, so the latter transformation
would restrict to a projective automorphism of that line.  Hence
\(f_0\) and \(f_1\) would be equivalent by independent projective
changes in source and target.  This is impossible.  The map \(f_0\)
has two ramification points, each of index three.  The map \(f_1\)
has the two simple zeros of \(3z^2+1\), each of index two, and
infinity, of index three.  Ramification indices are invariant under
those projective changes.  Thus the unmarked algebra scheme
distinguishes these two globally different coefficient maps even
though its fibre algebras, conormal bundle, abstract first-relation
bundle, and all graded vector bundles agree.
\end{proof}

\subsection{Pencils and the category of their automorphisms}
For pencils, take \(q=2p\), \(p=N-2\), and
\(F(H)=\bigwedge^{N-2}\Sym^2 H\).  By
Lemma~\ref{lem:pencil-irreducibility} this is a single Schur module
of dimension \(M=\binom N2\), and
Lemma~\ref{lem:moving-coefficients-v146} identifies its coefficient
subbundle with
\(\mathcal L_{\mathcal R}=\bigwedge^p(\Sym^2V/\mathcal R)^*\).
Its rank is one, so the strict-support hypothesis holds for every
pencil.  Exterior duality recovers the subbundle \(\mathcal R\)
from that line.  Theorem~\ref{thm:general-moving-coefficients-v147}
therefore proves the unrestricted moving-pencil inverse, including
its actual source-bundle conclusion.  This is a simultaneous
reconstruction theorem for a map into a coefficient Grassmannian,
not just a classification of its individual fibres.

\begin{corollary}[Algebraic automorphism sequence for a pencil]
\label{cor:pencil-group-scheme-v147}
Let \(\mathfrak A_R^{[d]}\) be the unmarked local algebra of
Theorem~\ref{thm:artin-local-inverse-v146}, and let
\(G_R=\{g\in\operatorname{GL}(V):(\Sym^2g)R=R\}\).
There is a split exact sequence of complex affine algebraic groups
\begin{equation}\label{eq:pencil-algebraic-group-v147}
 1\longrightarrow U_R\longrightarrow
 \operatorname{Aut}(\mathfrak A_R^{[d]})
 \longrightarrow (G_R\times\operatorname{GL}(U))/\mathbb G_m
 \longrightarrow1,
\end{equation}
where \(\mathbb G_m\) is the diagonal scalar subgroup.  Its
intrinsic quotient is the closed stabilizer of \(K_R\) in the
general linear group of the cotangent space.  In that form the
sequence represents automorphisms after extension to every
commutative complex algebra.  Its splitting is homogeneous after
choosing generators, and \(U_R\) has the intrinsic description and
dimension in Proposition~\ref{prop:first-relation-group-v147}.
\end{corollary}
\begin{proof}
Proposition~\ref{prop:first-relation-group-v147} gives the represented
sequence with quotient \(H_{K_R}\).  The coefficient inverse and
Lemma~\ref{lem:linear-preserver-v147} identify its complex points
with the image of the left-right action; transposition is excluded
by the support ranks \((1,M)\).  Use the contragredient of the
normal action to obtain a homomorphism on the cotangent space.
Its kernel is exactly the diagonal scalar subgroup: a tensor product
acting identically has both factors scalar, with reciprocal actions
on the dual factors.  The image is a closed algebraic subgroup and
all the groups in question are smooth in characteristic zero
\cite[Lemma~39.8.2]{StacksGroups}.  Thus the stabilizer and this
image agree as closed subgroup schemes, not only on complex points.
The induced quotient identification is an isomorphism of algebraic
groups: after division by the kernel it is a bijective separable
homomorphism between smooth groups.  The splitting and kernel are
those already constructed by substitutions.
\end{proof}

For a positive-dimensional projective reduction,
Theorem~\ref{thm:automorphism-kernel-v146} is used as an exact
sequence of groups of complex scheme automorphisms, with its stated
filtered kernel.  We do not identify that sequence with a finite-type
affine group scheme or a moduli-stack equivalence.  Likewise, the
marked construction over nonreduced parameter bases later in the
paper retains its specified relative socle: it does not replace that
ideal by the absolute nilradical.  These distinctions preserve the
full unmarked geometric assertion while specifying its category.

% BEGIN PRESERVED PART 13a-criterion-v145.tex
\section{An intrinsic coefficient principle}
\label{sec:structural-coefficient-principle}

Extracting coefficient data from an abstract scheme and reconstructing
a defining relation space from those data are distinct operations.
The criterion below isolates the first.  Its hypotheses will be
verified directly for the multiplication schemes and then on curves.

\subsection{The structural criterion}

Let \(X\) be a finite-type complex scheme with a closed subvariety
\(B\) specified by a construction preserved by every abstract scheme
isomorphism.  Write \(\mathfrak b\subset\OO_X\) for its ideal and
\[
 \mathcal G=\bigoplus_{r\ge0}\mathfrak b^r/\mathfrak b^{r+1},
 \qquad \mathcal A=\Sym_{\OO_B}(\mathcal G_1),
 \qquad \mathcal I=\ker(\mathcal A\twoheadrightarrow\mathcal G).
\]
Assume the following conditions, for fixed integers \(n\ge2,m\ge1\).

\smallskip
\noindent\textup{(i)} \(B\) is connected and projective,
\(H^0(B,\OO_B)=\C\), and \(\mathcal G_1\) is locally free.
Its dual admits a presentation \(V\otimes\mathcal U^*\), where
\(\dim V=\rank\mathcal U=n\).

\smallskip
\noindent\textup{(ii)} The reduction of \(\Spec_B\mathcal G\)
in its degree-one ambient bundle is the generic determinant cone.
The bundle \(\Pj_B(\mathcal U^*)\) is not projectively trivial.
Thus precisely one of the two ruling parameter bundles of its rank-one
locus is projectively trivial.

\smallskip
\noindent\textup{(iii)} If \(\mathcal D\subset\mathcal A\)
is the ideal of that reduced cone and \(\mathcal L=\mathcal D_n\),
then \(\mathcal I=\mathcal D\mathcal J\), where
\(\mathcal J\) is generated in degree \(m\).  Under the Cauchy
decomposition in any local frame of \(\mathcal U\),
\begin{equation}\label{eq:structural-coefficient-hypothesis}
 \mathcal J_m=
 \bigoplus_{\lambda\in\Lambda}
          L_\lambda\otimes\mathbb S_\lambda\mathcal U
 \ \subset\
 \Sym^m(V^*\otimes\mathcal U),
 \qquad L_\lambda\subset\mathbb S_\lambda V^*,
\end{equation}
where \(\Lambda\) is a specified collection of partitions of \(m\)
of length at most \(n\), and the left subspaces are constant over
\(B\).  Zero subspaces are permitted.  The Cauchy inclusions, not
just their representation labels, are the coefficient realizations
\(\beta\otimes\xi\mapsto[T\mapsto\beta(\mathbb S_\lambda(T)\xi)]\).

These conditions are hypotheses on the normal cone and its residual
coefficient module.  They do not assume that an abstract isomorphism
extends to an ambient space, preserves a chosen frame, or already acts
on all the left subspaces through a common projective transformation.

\begin{theorem}[Intrinsic coefficient reconstruction]
\label{thm:structural-coefficient-reconstruction}
Suppose \(X,B\) satisfy \textup{(i)--(iii)}.  The abstract scheme
\(X\) determines the collection \((L_\lambda)_{\lambda\in\Lambda}\)
up to one common projective coordinate transformation on \(V\).
More precisely, an isomorphism between two such schemes induces a
constant projective isomorphism \([g]:\Pj(V)\to\Pj(V')\) such that
\[
                L_\lambda=\mathbb S_\lambda(g)^*L'_\lambda
                       \quad(\lambda\in\Lambda).
\]
The equality means equality of subspaces; scalar changes of a linear
lift of \([g]\) do not affect it.
\end{theorem}
\begin{proof}
An isomorphism transports \(B\), every power of \(\mathfrak b\),
\(\mathcal G_1\), and the surjection
\(\mathcal A\twoheadrightarrow\mathcal G\).  In particular it
transports \(\mathcal I\) in the intrinsic degree-one ambient
bundle.  Iterated reduced singular loci of the determinant cone give
its rank-one locus.  Lemma~\ref{lem:rank-one-fano-scheme} reconstructs
its two ruling parameter schemes, including their scheme structure.
Condition \textup{(ii)} distinguishes the trivial ruling intrinsically.
Lemma~\ref{lem:oriented-segre-descent} then gives one constant
projective map \([g]\).  After choosing a constant linear lift and
local frames on both bases, the ambient bundle map has the form
\(T\mapsto gTh^{-1}\).  These local choices are made only after
constructing the graded algebra, its degree-one part, and the ruling.

Locally \(\mathcal D=(d)\) with \(d=\det T\).  Its leading
coefficient is a unit, so it is a nonzerodivisor over every coefficient
ring.  Hence \((\mathcal I:\mathcal D)=\mathcal J\), and
\[
 \mathcal J_m\simeq\mathcal L^*\otimes\mathcal I_{n+m}.
\]
This recovers the residual module without choosing a scalar determinant
equation.  Under \(T\mapsto gTh^{-1}\), the Cauchy coefficient
map transports a tensor by
\(\mathbb S_\lambda(g)^*\otimes\mathbb S_\lambda(h^{-1})\).
Projection to the specified Cauchy summand, followed by contraction
against its whole right dual, therefore transports exactly
\(L'_\lambda\) to \(L_\lambda\).  The right map is invertible
and has no effect on that left image.  Both determinant-line transition
functions and changes of a lift of \([g]\) act by scalar units.
Since every partition has size \(m\), this ambiguity is common
in degree and disappears on subspaces.  This proves the assertion.
\end{proof}


% BEGIN PRESERVED PART 18-pencil-proof-v145.tex
\section{Quadratic relation spaces and intrinsic reconstruction of pencils}
\label{sec:natural-pencils}

The multiplication construction itself supplies a wider natural class
for the coefficient criterion.  In this section the number of quadratic
relations is allowed to vary independently of the embedding dimension.
For two relations the resulting theorem recovers every pencil of
quadrics, including singular pencils.  No extra block is added to the
multiplication map.

\subsection{Automatic coefficient extraction for quadratic multiplication}

Let \(\dim V=n\ge3\), \(W=\Sym^2V\), \(N=n(n+1)/2\), and
\(R\in\Gr(r,W)\), where \(1\le r<N-n\).  Put \(p=N-r>n\),
\(S=W/R\), and
\[
 A_R=\C\oplus V\oplus S,\qquad
 X_R=\Gr(n,V\oplus S).
\]
The only nonzero product in the maximal ideal is the quotient map
\(\gamma:\Sym^2V\to S\).  Write \(\mathcal K\) for the
tautological bundle on \(X_R\), and use the same definition
\eqref{eq:failure-intro} of \(\widehat D_R\), with \(A_R\) in
place of \(B_R\).  Thus \(\widehat D_R\) is the failure scheme
of generating \(A_R\) by a varying \(n\)-plane and the unit.
The notation in this section agrees with the web construction when
\((n,r)=(4,4)\).

Let \(\Lambda_{n,p}\) denote the partitions occurring in the
classical multiplicity-free decomposition
\begin{equation}\label{eq:general-exterior-decomposition}
 F(V):=\bigwedge^p\Sym^2V
       =\bigoplus_{\lambda\in\Lambda_{n,p}}\mathbb S_\lambda V.
\end{equation}
Every partition has size \(2p\); terms of length greater than \(n\)
are omitted.  Choose a nonzero volume form on \(S\) and let
\(\beta_R\in F(V)^*\) be its pullback by \(\bigwedge^p\gamma\).
Its component in \(\mathbb S_\lambda V^*\) is denoted by
\(\beta_{R,\lambda}\).  Only the lines of the nonzero components
will be used, so the volume choice is immaterial.

\begin{theorem}[Automatic intrinsic coefficient extraction]
\label{thm:automatic-quadratic-coefficients}
For every \(R\in\Gr(r,\Sym^2V)\) in the stated range, the
abstract scheme \(\widehat D_R\) determines
\[
       \bigl(\C\beta_{R,\lambda}\bigr)_{\lambda\in\Lambda_{n,p}}
\]
up to one common projective transformation of \(V\), with zero
components recorded as zero subspaces.  More precisely, the intrinsic
construction recovers the base \(B=\Gr(n,S)\), its degree-one
normal-cone bundle, and the residual ideal generated in degree \(2p\)
with coefficient module
\begin{equation}\label{eq:automatic-residual-module}
 \mathcal J_{2p}=
 \bigoplus_{\lambda\in\Lambda_{n,p}}
       \C\beta_{R,\lambda}\otimes\mathbb S_\lambda\mathcal U
 \ \subset\ \Sym^{2p}(V^*\otimes\mathcal U),
\end{equation}
where \(\mathcal U\) is the tautological rank-\(n\) bundle on
\(B\).  Consequently all geometric hypotheses of
Theorem~\ref{thm:structural-coefficient-reconstruction} hold for this
entire Grassmannian family, not just for smooth webs.
\end{theorem}
\begin{proof}
On a frame chart of \(\mathcal K\), write its inclusion as
\(\binom M{A_0}\), with \(M\) an \(n\times n\) matrix.  Since
\(A_R\) is cube-zero, columns of degree greater than two vanish.
After retaining the scalar, linear, and quadratic columns, multiplication
has the matrix
\begin{equation}\label{eq:general-multiplication-matrix}
 \begin{pmatrix}
  1&0&0\\
  0&M&0\\
  0&A_0&\gamma\Sym^2M
 \end{pmatrix}.
\end{equation}
A nonzero maximal minor must use the scalar column and all \(n\)
linear columns: these are the only columns meeting the first \(1+n\)
rows.  Expanding along those rows gives the identity of ideals
\begin{equation}\label{eq:general-fitting-factorization}
 I_{\widehat D_R}=(\det M)I_p(\gamma\Sym^2M).
\end{equation}
The same expansion works for every \(m\ge2\).  Where \(M\) is
invertible, the quadratic map is surjective.  Hence the reduction of
\(\widehat D_R\) is exactly the Schubert divisor where \(M\)
is singular, independent of \(R\) after identifying the vector
spaces \(S\).

Successive reduced singular loci of that divisor are the projection
rank loci.  Their last nonempty member is \(B=\Gr(n,S)\).
For clarity, the usual determinantal assertion here is local: at a
projection-rank-\(k\) point, an invertible \(k\)-minor splits
away, leaving a generic square \((n-k)\)-matrix and smooth
coordinates.  Every remaining entry can vary independently by
moving the kernel of the projection in \(V\) modulo its image.
Thus its reduced singular locus is the next rank locus.  Since
\(p>n\), the final base is nonempty, connected and projective,
with \(H^0(B,\OO_B)=\C\).

Near \(B\), a Grassmannian chart consists of a varying plane in
\(S\) and its graph \(T:\mathcal U\to V\).  Formula
\eqref{eq:general-fitting-factorization} is independent of the other
graph coordinates and homogeneous in \(T\).  It therefore gives
its actual normal-cone ideal, not merely an initial reduced support:
\begin{equation}\label{eq:general-normal-cone}
 C_B\widehat D_R\subset\Tot\Hom(\mathcal U,V),\qquad
 \mathcal I=(\det T)\mathcal J,\quad
 \mathcal J=I_p(\gamma\Sym^2T).
\end{equation}
There are no degree-one equations.  The ambient bundle in
\eqref{eq:general-normal-cone} is consequently the dual of the
intrinsically recovered degree-one conormal.  Its reduced ideal is
the determinant ideal \(\mathcal D\).  The determinant polynomial
is monic for a suitable monomial order, so cancellation gives
\((\mathcal I:\mathcal D)=\mathcal J\).  This is an intrinsic
ideal quotient; choosing a generator of the determinant line is not
part of the recovered datum.

The rank-one locus of the reduced cone determines the two ruling
parameter schemes by Lemma~\ref{lem:rank-one-fano-scheme}.  They are
\(\Pj(V)\times B\) and \(\Pj_B(\mathcal U^*)\).
Exactly the first is projectively trivial.  Indeed a Schubert line in
\(\Gr(n,S)\) restricts \(\mathcal U^*\) to
\(\OO(1)\oplus\OO^{n-1}\).  A trivial projective bundle on
that line would make this bundle a line-bundle twist of
\(\OO^n\); its determinant would then have degree divisible by
\(n\), rather than degree one.  Thus the trivial ruling is
intrinsically oriented.  Lemma~\ref{lem:oriented-segre-descent}
applies: an abstract isomorphism gives a morphism \(B\to\PGL(V)\),
constant because \(B\) is connected projective and \(\PGL(V)\)
is affine.  Local linear lifts and right changes of frame give
\(T\mapsto gTh^{-1}\) for this one fixed \([g]\).

It remains to prove \eqref{eq:automatic-residual-module}, rather
than assume it as a condition on the example.  The minors are the
image, for fixed \(\beta_R\), of the universal coefficient map
\begin{equation}\label{eq:general-coefficient-map}
 F(V)^*\otimes F(U)\longrightarrow
 \Sym^{2p}(V^*\otimes U),\qquad
 \beta\otimes\xi\longmapsto[T\mapsto\beta(F(T)\xi)].
\end{equation}
In characteristic zero, the polynomial functor \(F\) has the
multiplicity-free decomposition \eqref{eq:general-exterior-decomposition};
see \cite[Theorem~3.5]{ChengWang}.  The Cauchy decomposition of the
target contains each diagonal
\(\mathbb S_\lambda V^*\otimes\mathbb S_\lambda U\)
once, and no off-diagonal pair.  Equivariance kills every
off-diagonal pair.  On each diagonal pair the coefficient map is
nonzero: identify \(U\) with \(V\), evaluate at \(T=1\), and
choose \(\beta,\xi\) in that summand with
\(\beta(\xi)\ne0\).  Schur's lemma then identifies its image
with the specified Cauchy summand.  This is the actual matrix-coefficient
inclusion, not merely a dimension count.  Keeping \(\beta_R\)
fixed proves \eqref{eq:automatic-residual-module}.  All generators
have degree \(2p\), and at least one minor is nonzero since
\(\gamma\) is surjective.

Under \(T\mapsto gTh^{-1}\), the right map acts within the full
factor \(\mathbb S_\lambda U\).  Contracting against its dual
therefore leaves exactly the left coefficient line transformed by
\(\mathbb S_\lambda(g)^*\).  Determinant-line scalars and local
lift ambiguities do not change these lines.  This proves the theorem.
\end{proof}

The theorem does not assert that separately projectivized Schur
components always determine their relative scalars.  For webs this
is the recombination problem solved in the separate research archive.  For pencils that issue
disappears altogether.

\subsection{Every pencil, without a regularity hypothesis}

\begin{lemma}[The exterior square of the quadratic representation]
\label{lem:pencil-irreducibility}
For \(n\ge2\), \(\bigwedge^2\Sym^2V=\mathbb S_{(3,1)}V\)
is irreducible.  Consequently
\(\bigwedge^{N-2}\Sym^2V\) is irreducible as well.
\end{lemma}
\begin{proof}
The vector \(x_1^2\wedge x_1x_2\) is a highest-weight vector of
weight \((3,1)\).  By complete reducibility it generates a copy of
\(\mathbb S_{(3,1)}V\).  The hook formula gives
\[
 \dim\mathbb S_{(3,1)}V
   =\frac{n(n-1)(n+1)(n+2)}8
   =\binom{N}{2}.
\]
This exhausts the exterior square.  The last assertion follows by
exterior duality.  These are classical representation identities.
\end{proof}

\begin{theorem}[Intrinsic reconstruction of all quadratic pencils]
\label{thm:natural-pencil-reconstruction}
Let \(n\ge3\) and let \(R,R'\) be two-dimensional subspaces
of \(\Sym^2V\).  For the natural multiplication-failure schemes
of their cube-zero algebras of Hilbert function \((1,n,N-2)\),
\begin{equation}\label{eq:natural-pencil-torelli}
 \widehat D_R\simeq\widehat D_{R'}
 \quad\Longleftrightarrow\quad
 R'=gR\text{ for some }g\in\PGL(V).
\end{equation}
No pencil, base locus, or discriminant is required to be smooth or
regular.  The isomorphism on the left is abstract and carries no
supplied ambient embedding, base, or tensor-factor marking.  All the
reduced schemes are isomorphic to the same Schubert divisor.
\end{theorem}
\begin{proof}
Here \(p=N-2>n\).  By
Theorem~\ref{thm:automatic-quadratic-coefficients} and
Lemma~\ref{lem:pencil-irreducibility}, the normal cone has only one
left coefficient line to recover: the entire line \(\C\beta_R\).
The canonical duality
\begin{equation}\label{eq:pencil-exterior-duality}
 \bigwedge^{N-2}W^*\simeq
       \bigwedge^2W\otimes(\det W)^*
       =\bigwedge^2W\otimes(\det V)^{-(n+1)}
\end{equation}
sends that line to the Pluecker line of \(R\), with the same
projective transformation \(g\).  Indeed the pullback of a
quotient volume is the annihilator-volume dual of the kernel
\(R\).  The common determinant character is a scalar on
projective space.  A nonzero decomposable two-vector recovers its
plane as \(\{w\in W:w\wedge w_R=0\}\).  Thus an abstract
scheme isomorphism implies \(R'=gR\).

Conversely such a linear coordinate change induces the quotient
isomorphism \(S_R\to S_{R'}\), the algebra isomorphism
\(A_R\to A_{R'}\), and the corresponding Grassmannian isomorphism.
It intertwines multiplication and hence preserves its zeroth Fitting
ideal.  This proves the reverse implication.  The assertion about
reductions was proved in \eqref{eq:general-fitting-factorization}.
\end{proof}

\subsection{Discriminants, elementary divisors, and hyperelliptic moduli}

A regular pencil means one containing an invertible symmetric matrix;
it need not have a reduced discriminant.  The classical invariant of
such a pencil consists of its discriminant points together with their
Jordan-block partitions, usually called its Weierstrass--Segre data.
We use the root-labelled partitions, not a pairing of an unlabelled
list with an independently unlabelled divisor.  Their classification
is classical; compare \cite[Theorem~1.1, Corollary~2.1 and \S2]{FevolaMandelshtamSturmfels}.
The following consequence identifies precisely what the intrinsic
failure scheme remembers beyond the discriminant.

\begin{corollary}[Intrinsic recovery of elementary divisors]
\label{cor:intrinsic-segre-data}
For a regular pencil in any dimension \(n\ge3\), the abstract
scheme \(\widehat D_R\) determines all root-labelled elementary
divisors of a matrix pencil, up to projective reparametrization of its
parameter line.  Conversely these data determine its abstract failure
scheme.  In particular the invariant distinguishes different Jordan
partitions over a fixed multiple discriminant root.
\end{corollary}
\begin{proof}
The first implication follows by reconstructing \(R\), then
choosing an invertible member \(A\) and a complementary member
\(B\).  For completeness, the relevant classical sufficiency can
be seen directly.  Put \(C=A^{-1}B\); it is self-adjoint for
\(A\).  If \(hC=C'h\) identifies the Jordan forms for another
pair \((A',B')\), set
\[
 H=A^{-1}h^tA'h.
\]
Then \(H\) is invertible, \(A\)-self-adjoint, and commutes with
\(C\).  There is a polynomial \(k\) in \(H\) with
\(k^2=H\): on each nonzero eigenvalue choose a square root and
truncate its Taylor series to the size of the largest Jordan block,
then combine by the Chinese remainder theorem.  Thus \(k\) is
invertible, self-adjoint, and commutes with \(C\).
The map \(hk^{-1}\) is an isometry from \(A\) to \(A'\)
and intertwines \(C,C'\), so it simultaneously identifies
\(A,B\) with \(A',B'\).

Changing the basis of the pencil replaces \(C\) by a fractional
linear expression with invertible denominator.  At each eigenvalue
its derivative is nonzero; hence each Jordan-block size is preserved
while its root moves by the same projective transformation.  A
parameter change can avoid infinity for both finite spectra.
Root-labelled elementary divisors therefore suffice, and
\eqref{eq:natural-pencil-torelli} proves the converse for the
failure schemes.  This argument recalls the classical classification;
the new input is recovery from an unmarked abstract failure scheme.
\end{proof}

\begin{example}[Identical discriminant, distinct intrinsic failure schemes]
\label{ex:same-discriminant-distinct-failure}
Choose distinct nonzero \(\mu_3,\ldots,\mu_n\), and put
\[
 S_2=\begin{pmatrix}0&1\\1&0\end{pmatrix},\qquad
 J_2=\begin{pmatrix}0&1\\0&0\end{pmatrix}.
\]
Consider the symmetric pairs
\begin{align*}
 A_{\mathrm{s}}&=I_n,&
 B_{\mathrm{s}}&=\diag(0,0,\mu_3,\ldots,\mu_n),\\
 A_{\mathrm{j}}&=S_2\oplus I_{n-2},&
 B_{\mathrm{j}}&=S_2J_2\oplus\diag(\mu_3,\ldots,\mu_n).
\end{align*}
Their determinant binary forms agree up to a nonzero scalar:
\begin{equation}\label{eq:same-discriminant-pencils}
 \det(sA_{\mathrm{s}}+tB_{\mathrm{s}})
 =-\det(sA_{\mathrm{j}}+tB_{\mathrm{j}})
 =s^2\prod_{i=3}^n(s+\mu_i t).
\end{equation}
At the unique double root the ranks are respectively \(n-2\)
and \(n-1\), so no simultaneous projective equivalence exists.
The two root partitions are \((1,1)\) and \((2)\).
Their abstract failure schemes are therefore nonisomorphic although
both their reduced failure schemes and their discriminant divisors,
including multiplicities, agree.  This realizes inside the same
multiplication construction the classical loss of elementary-divisor
information under passage to the determinant.
\end{example}

\begin{corollary}[Hyperelliptic isomorphism classes]
\label{cor:hyperelliptic-failure-invariant}
Fix \(g\ge2\) and \(n=2g+2\).  On the locus of pencils with
simple discriminant, \(\widehat D_R\) gives an injective invariant
of the isomorphism classes of smooth hyperelliptic curves of genus
\(g\).  All these failure schemes have the same reduction; their
nonreduced structures distinguish a \((2g-1)\)-dimensional family
of curve classes.  The pencil invariant also distinguishes the two
regular boundary pencils in
Example~\ref{ex:same-discriminant-distinct-failure}, where their
singular discriminant double covers alone do not.
\end{corollary}
\begin{proof}
The classical curve associated with a simple-discriminant pencil is
the double cover of its parameter line branched at
\(\det(sA+tB)=0\); see \cite[\S4.2, equation (4.3)]{SmithQuadrics}.
Choose an invertible member.  Distinct eigenvalues of \(A^{-1}B\)
have mutually \(A\)-orthogonal eigenlines, each nonisotropic, so
simultaneous diagonalization identifies the pencil with
\(\langle\sum x_i^2,\sum\lambda_i x_i^2\rangle\).
The unordered branch divisor modulo \(\PGL_2\) therefore
determines the pencil, and every reduced divisor of degree \(2g+2\)
arises.  A smooth curve of genus at least two has a unique degree-two
map to a rational curve up to automorphism of that curve: two distinct
such maps would give a birational map to a curve of bidegree \((2,2)\)
in \(\Pj^1\times\Pj^1\), whose arithmetic genus is one.
Thus curve isomorphism is exactly branch-divisor equivalence.
Applying Theorem~\ref{thm:natural-pencil-reconstruction} gives the
asserted injectivity and well-definedness.

Ordered distinct branch points have dimension \(2g+2\).
Normalizing three of them removes the three-dimensional projective
action; taking the finite permutation quotient leaves dimension
\(2g-1\).  Equality of all reductions follows from
\eqref{eq:general-fitting-factorization}.  Finally the two binary
forms in \eqref{eq:same-discriminant-pencils} differ only by a
scalar, a square over \(\C\), and define isomorphic double covers,
whereas their failure schemes have already been separated.
\end{proof}

These are statements about complex isomorphism classes.  They do not
identify moduli stacks, assert a universal family over a coarse moduli
space, or eliminate the different automorphism groups of pencils and
curves.  Nor is the classical pencil--hyperelliptic correspondence a
new Torelli theorem.  The conclusion here is reconstruction of those
objects and of the elementary-divisor information on their boundary
from the abstract nonreduced multiplication-failure scheme, while all
reduced failure data are fixed.

% BEGIN PRESERVED PART 20-finite-neighbourhoods.tex
\section{The exact infinitesimal order of pencil reconstruction}
\label{sec:finite-neighbourhoods}

The all-pencil theorem admits a finite-order strengthening.  It
identifies the first infinitesimal neighbourhood that retains the
pencil, and shows that every preceding neighbourhood is independent
of it.  We use ideal-adic neighbourhoods of the deepest stratum;
these are not arc spaces or jet schemes.

Throughout this section let \(n\ge3\), \(W=\Sym^2V\), and
\begin{equation}\label{eq:pencil-detection-order}
 N=\binom{n+1}{2},\qquad p=N-2,\qquad
 d=n+2p=n^2+2n-4.
\end{equation}
For a pencil \(R\subset W\), put \(S_R=W/R\) and
\(B=\Gr(n,S_R)\subset\widehat D_R\).  This notation for the
socle Grassmannian is local to this section; the cube-zero algebra
is denoted by \(A_R\), as in Section~\ref{sec:natural-pencils}.
Write \(\mathfrak b\subset\OO_{\widehat D_R}\) for the
ideal of \(B\), and define
\begin{equation}\label{eq:finite-neighbourhood-definition}
 Z_R^{[k]}=(B,\OO_{\widehat D_R}/\mathfrak b^{k+1})
       \qquad(k\ge0).
\end{equation}
Thus \(Z_R^{[0]}=B\).  No embedding or identification of its
reduction is supplied with the abstract scheme \(Z_R^{[k]}\).

\begin{theorem}[Sharp finite-neighbourhood reconstruction]
\label{thm:sharp-finite-pencil}
For any two pencils \(R,R'\subset\Sym^2V\), the following
conditions are equivalent:
\begin{equation}\label{eq:finite-pencil-equivalences}
\begin{split}
 Z_R^{[d]}\simeq Z_{R'}^{[d]}
 &\quad\Longleftrightarrow\quad
 \widehat D_R\simeq\widehat D_{R'}\\
 &\quad\Longleftrightarrow\quad
 R'=gR\quad\text{for some }g\in\PGL(V).
\end{split}
\end{equation}
In contrast, for every \(0\le k<d\), all the schemes
\(Z_R^{[k]}\) are isomorphic.  Consequently \(d\) is the
smallest uniform order at which these neighbourhoods classify all
pencils up to projective equivalence.  Singular pencils are included.
\end{theorem}

\subsection{The first relation in the conormal algebra}

Let \(\mathcal U\) be tautological on \(B\).  The open set
of \(X_R\) where projection to \(S_R\) is injective is
canonically \(\Tot E\), where
\(E=\Hom(\mathcal U,V)\): a plane is the graph of a unique
map from its image in \(S_R\) to \(V\).  In particular this
is one open neighbourhood of all of \(B\), not merely a
collection of unrelated formal charts.  Let \(\mathfrak a\)
be its zero-section ideal.  The exact multiplication calculation
\eqref{eq:general-fitting-factorization} gives
\begin{equation}\label{eq:finite-homogeneous-ideal}
 \mathcal I_R=(\det T)I_p(\gamma_R\Sym^2T)
            \subset\Sym E^*.
\end{equation}
It is generated in degree \(d\), with a nonzero degree-\(d\)
piece \(K_R\).  Indeed each quadratic minor has degree \(2p\),
and an invertible \(T\) makes \(\gamma_R\Sym^2T\)
surjective.  There are no terms of lower degree, even on special
pencils.  Homogeneity is an equality of ideals in the actual graph
neighbourhood, not just an equality of initial supports.

\begin{lemma}[Intrinsic recovery from the first relation]
\label{lem:first-relation-cone}
The abstract scheme \(Z_R^{[d]}\) determines, as a graded
\(\OO_B\)-algebra, the complete normal-cone algebra
\(\gr_{\mathfrak b}\OO_{\widehat D_R}\).
\end{lemma}
\begin{proof}
The reduction of \(Z_R^{[d]}\) is \(B\), since its
zero-section ideal is nilpotent and \(B\) is reduced.  Hence
its nilradical \(\mathfrak n\) is intrinsic.  Put
\(L=\mathfrak n/\mathfrak n^2\).  Formula
\eqref{eq:finite-homogeneous-ideal} and \(d>1\) identify
\(L=E^*\).  Multiplication in the given finite scheme recovers
\begin{equation}\label{eq:intrinsic-first-relation}
 K_R=\ker\bigl[\Sym^d L\longrightarrow
                    \mathfrak n^d/\mathfrak n^{d+1}\bigr].
\end{equation}
The last denominator is zero, but retaining it displays the grading.
No splitting of the reduction map is needed to form this kernel.
Because \(\mathcal I_R\) is generated in this single degree,
\begin{equation}\label{eq:finite-cone-completion}
 \gr_{\mathfrak b}\OO_{\widehat D_R}
       =\Sym_{\OO_B}L/(K_R).
\end{equation}
Here \((K_R)\) denotes the homogeneous ideal generated by
\(K_R\), not its radical.  The right side is intrinsic to the
finite scheme, proving the assertion.
\end{proof}

\begin{proof}[Proof of Theorem~\ref{thm:sharp-finite-pencil}]
For \(k<d\), equations \eqref{eq:finite-homogeneous-ideal}
and \eqref{eq:finite-neighbourhood-definition} give
\[
 \OO_{Z_R^{[k]}}=\Sym E^*/\mathfrak a^{k+1}.
\]
A linear identification \(S_R\simeq S_{R'}\) identifies
\(B\), its tautological bundle, and \(E\).  It therefore
identifies every one of these lower-order neighbourhoods.

An abstract isomorphism at order \(d\) induces an isomorphism
of reductions, of conormal bundles, and, by
Lemma~\ref{lem:first-relation-cone}, of complete graded normal
cones.  We spell out why the reconstruction in
Theorem~\ref{thm:automatic-quadratic-coefficients} requires no
information outside this cone.  Its reduced support is the
square-matrix determinant cone; its iterated singular loci give
the rank-one cone.  The relative scheme of maximal linear spaces
in that cone has the two rulings
\(\Pj(V)\times B\) and \(\Pj_B(\mathcal U^*)\).
Exactly the first is projectively trivial, as is seen by restriction
to a Schubert line.  Thus it is intrinsically oriented.  The
induced projective transformation of the trivial factor is constant
on the connected projective base.  Local lifts have the form
\(T\mapsto gTh^{-1}\) for this one \([g]\).
The determinant ideal line is intrinsic, so its ideal quotient in
\eqref{eq:finite-homogeneous-ideal} recovers the residual ideal.
Finally Lemma~\ref{lem:pencil-irreducibility} identifies the entire
left coefficient line \(\C\beta_R\), and exterior duality
\eqref{eq:pencil-exterior-duality} recovers the Pluecker line of
\(R\).  These are operations on the recovered normal cone;
we do not assume that its isomorphism extends to \(X_R\) or
\(\widehat D_R\).  They imply \(R'=gR\).

Conversely such a projective transformation induces the algebra
and failure-scheme isomorphisms, carrying the socle Grassmannian
to the socle Grassmannian and preserving all its ideal powers.
This proves \eqref{eq:finite-pencil-equivalences}, also using
Theorem~\ref{thm:natural-pencil-reconstruction}.  For sharpness,
there are inequivalent pencils in every dimension under consideration:
\(\langle x_1^2,x_1x_2\rangle\) has a common linear factor,
whereas \(\langle x_1^2,x_2^2\rangle\) does not.  Their
neighbourhoods are nevertheless isomorphic at every order below
\(d\).  Thus no smaller order classifies all pencils.
\end{proof}

\subsection{A closed immersion of the full parameter scheme}

The finite-order result has a scheme-theoretic parameter version,
which records infinitesimal variation as well as closed points.
Here tensor coordinates are fixed.  This additional qualification
is important: we are not asserting an equivalence of unmarked
moduli stacks.

Fix an auxiliary \(n\)-dimensional vector space \(U\), and put
\[
 F(V)=\bigwedge^p\Sym^2V,\quad F(U)=\bigwedge^p\Sym^2U,
 \quad M=\dim F(U)=\binom N2,\quad H=\Hom(U,V).
\]
Choose a nonzero scalar representative \(\delta\) of the
determinant polynomial on \(H\); its scalar has no effect on
the subspaces below.  Let
\(K_R\subset\Sym^dH^*\) be the degree-\(d\) space of
\((\det T)I_p(\gamma_R\Sym^2T)\) in these fixed coordinates.

\begin{theorem}[Faithful finite-order parameterization]
\label{thm:finite-parameter-immersion}
The assignment
\begin{equation}\label{eq:finite-parameter-map}
 \Phi:\Gr(2,\Sym^2V)\longrightarrow\Gr(M,\Sym^dH^*),
       \qquad R\longmapsto K_R,
\end{equation}
is a closed immersion.  It is defined on the whole Grassmannian,
including singular pencils.  Its construction and its inverse on
the image commute with arbitrary base change over \(\C\).
In particular it separates families over nonreduced bases and has
injective differential at every pencil.
\end{theorem}

\begin{lemma}[Tensoring a line with a fixed space]
\label{lem:line-subspace-immersion}
For finite-dimensional vector spaces \(A,B_0\), with
\(b=\dim B_0>0\), the map
\(\Pj(A)\to\Gr(b,A\otimes B_0)\),
\(\ell\mapsto\ell\otimes B_0\), is a closed immersion.
\end{lemma}
\begin{proof}
Choose bases \(a_0,\ldots,a_t\) and a basis of \(B_0\).
Above the projective chart
\(\ell=\C(a_0+\sum x_i a_i)\), the subspace is the graph
of the block matrices \(x_i I_b\).  The corresponding
Grassmannian chart has arbitrary \(b\times b\) blocks.  The
image is cut out by their off-diagonal entries and the differences
between their diagonal entries; an inverse reads one diagonal
entry in each block.  These are scheme-theoretic equations over
any \(\C\)-algebra, not just pointwise conditions.  The
preimage of this Grassmannian chart is exactly the stated projective
chart, since the projection has determinant \(a_0^b\).
Such charts cover the image.  The map is proper because its source
is projective, so these chartwise closed immersions give a global
closed immersion.
\end{proof}

\begin{proof}[Proof of Theorem~\ref{thm:finite-parameter-immersion}]
The coefficient map \eqref{eq:general-coefficient-map} becomes
\[
 \mu:F(V)^*\otimes F(U)\longrightarrow\Sym^{2p}H^*.
\]
It is injective: by Lemma~\ref{lem:pencil-irreducibility} its
source is irreducible for \(\operatorname{GL}(V)\times
\operatorname{GL}(U)\), and evaluation at an isomorphism
\(U\simeq V\) shows that the equivariant map is nonzero.
Multiplication by the nonzero polynomial \(\delta\) is
injective.  Thus \(\nu=\delta\mu\) is a fixed linear
inclusion and
\begin{equation}\label{eq:finite-coefficient-factorization}
 K_R=\nu(\C\beta_R\otimes F(U)).
\end{equation}
In particular \(\dim K_R=M\), also on every singular pencil.
By exterior duality the quotient-volume map
\(R\mapsto\C\beta_R\) is the Pluecker closed immersion
\(\Gr(2,W)\hookrightarrow\Pj(F(V)^*)\).
Factor \(\Phi\) through this map, the closed immersion of
Lemma~\ref{lem:line-subspace-immersion}, and the Grassmannian
closed immersion induced by \(\nu\).  This proves the theorem.
All are morphisms of schemes defined by the displayed polynomial
and linear equations.  A closed immersion is a monomorphism and
remains one after arbitrary base change.  This proves the family
assertion, not merely separation of complex-valued points.
\end{proof}

\begin{corollary}[A flat family of finite normal algebras]
\label{cor:finite-flat-family}
Over \(\Gr(2,W)\), the algebras
\begin{equation}\label{eq:finite-flat-algebra}
 Q_R=\Sym H^*/\bigl((K_R)+\mathfrak m^{d+1}\bigr),
 \qquad \mathfrak m=\bigoplus_{j>0}\Sym^jH^*,
\end{equation}
form a finite locally free family of rank
\[
              \binom{n^2+d}{d}-\binom N2.
\]
In fixed coordinates this family determines every base-changed
family of pencils.  All its truncations below degree \(d\)
are constant families.
\end{corollary}
\begin{proof}
Theorem~\ref{thm:finite-parameter-immersion} makes the spaces
\(K_R\) a rank-\(M\) subbundle of the constant degree-\(d\)
bundle, with locally free quotient.  In degree \(j<d\) the
algebra is \(\Sym^jH^*\), and its degree-\(d\) piece is
\(\Sym^dH^*/K_R\); higher degrees vanish.  This proves
local freeness and base-change compatibility.  Summing the
ranks gives the displayed binomial identity.  The kernel of the
coordinate quotient in degree \(d\) recovers \(K_R\), and
\(\Phi\) recovers the pencil family.  This is a statement
about coordinate quotient algebras; it does not discard their
coordinate maps and then infer a family reconstruction theorem.
\end{proof}

\subsection{Classical moduli at a sharp order}

\begin{corollary}[Hyperelliptic classes in finite neighbourhoods]
\label{cor:finite-hyperelliptic}
Let \(g\ge2\), \(n=2g+2\), and
\(d=4g^2+12g+4\).  For pencils with simple discriminant,
\(Z_R^{[d]}\) and \(Z_{R'}^{[d]}\) are isomorphic if
and only if their associated hyperelliptic curves are isomorphic.
Every lower-order neighbourhood is independent of the curve.
Thus the same \((2g-1)\)-dimensional family of curve classes
is invisible to all these lower-order neighbourhoods and is
separated at order \(d\).
\end{corollary}
\begin{proof}
Combine Theorem~\ref{thm:sharp-finite-pencil} with
Corollary~\ref{cor:hyperelliptic-failure-invariant}, and substitute
\(n=2g+2\) in \eqref{eq:pencil-detection-order}.
\end{proof}

The order here concerns a neighbourhood of the entire projective
socle Grassmannian.  Its global ruling distinguishes the constant
\(V\)-factor from the tautological factor.  A single unmarked
Artinian normal slice does not supply this orientation; the
coordinate statement of Corollary~\ref{cor:finite-flat-family}
is deliberately different.  The result is therefore not a claim
that one unmarked local algebra replaces the projective family.
The matrix-pencil classification and the hyperelliptic
correspondence remain classical.  The assertion proved here is
that this classical information first appears at the specified
order in the intrinsic multiplication-failure neighbourhoods.

% BEGIN PRESERVED PART 32-local-inverse-v146.tex
\section{Coefficient-support orientation and a local algebraic inverse}
\label{sec:local-inverse-v146}

For pencils there is a second, entirely local way to distinguish the
two tensor rulings.  It uses the asymmetry of the residual coefficient
module rather than the topology of a projective bundle.  Consequently
the nontrivial-projectivization hypothesis in the preceding family
theorem can be removed, and a single unmarked Artin local algebra
already recovers the pencil.

\subsection{The coefficient module orients its own tensor factors}
If \(K\subset A\otimes B\) is a linear subspace, define its two
factor supports by
\[
 s_A(K)=\im(K\otimes B^*\longrightarrow A),\qquad
 s_B(K)=\im(K\otimes A^*\longrightarrow B).
\]
The maps here are contractions.  Their ranks are invariant under
separate changes of basis and are exchanged by interchanging the
factors.  In particular, if \(K=L\otimes B\), then
\(s_A(K)=L\) and \(s_B(K)=B\).  The support ranks are therefore
\((\dim L,\dim B)\), not the common dimension \(\dim K\).

\begin{lemma}[Asymmetric coefficient support]
\label{lem:coefficient-orientation-v146}
Suppose a homogeneous ideal in the matrix space \(\Hom(U,V)\), with
\(\dim U=\dim V=n\ge3\), has reduced support the determinant cone.
Suppose its intrinsic residual ideal after determinant cancellation
has, in one multiplicity-one Cauchy summand of degree \(q\), the
coefficient subspace
\[
 L\otimes\mathbb S_\lambda U
 \ \subset\ \mathbb S_\lambda V^*\otimes\mathbb S_\lambda U,
 \qquad 0<\dim L<\dim\mathbb S_\lambda U.
\]
A linear isomorphism carrying this ideal to another ideal with a
coefficient subspace of the same displayed type cannot interchange
the two matrix rulings.  The conclusion is unchanged by line twists
of the locally chosen tensor factors.  It also applies fibrewise to
subbundles of this type over a reduced base.
\end{lemma}
\begin{proof}
A linear isomorphism of determinant cones preserves their successive
rank loci and therefore the rank-one cone.  Its action on the two
ruling components either preserves them or interchanges them.  To
make the familiar two alternatives explicit, identify the rank-one
projectivization with \(\Pj(V)\times\Pj(U^*)\).  An automorphism
preserving its two families of maximal linear spaces acts separately
on the two projective factors.  Each factor map is projective linear,
by restricting the ambient linear map to a ruling.  Removing their
linear lifts leaves a map fixing every decomposable line; sums of
tensors with one factor in common show that this remaining map is
scalar.  The preserving alternative is thus a left-right matrix
map.  In the interchanging alternative, compose with transposition
to obtain the preserving alternative.

The determinant ideal is carried to the determinant ideal, so its
ideal quotient is carried to the residual ideal.  The Cauchy
decomposition is functorial for separate changes of matrix factors;
transposition exchanges the factors in its \(\lambda\)-summand.
Put \(m=\dim\mathbb S_\lambda U\), and \(r=\dim L<m\).
The displayed subspace has support ranks \((r,m)\).  Separate
factor maps preserve these ranks, while transposition changes them
to \((m,r)\).  An ideal of the prescribed type on the other side
has support ranks \((r',m)\) with \(r'<m\).  It cannot be the
transposed image.  This excludes interchanging the rulings.

Tensoring a factor by a line changes its Schur module by a line
power, not its contraction rank.  The argument therefore does not
require a specified linear lift of either projective factor.  For
bundle maps it applies to each geometric fibre; an interchange on
any fibre would contradict the same strict rank inequality there.
\end{proof}

For pencils, take
\(F(H)=\bigwedge^{N-2}\Sym^2H\).
Lemma~\ref{lem:pencil-irreducibility} makes this a single irreducible
Schur module, of dimension \(M=\binom N2>1\).  The residual module
is precisely
\(\mathcal L_{\mathcal R}\otimes F(\mathcal U)\), by
Lemma~\ref{lem:moving-coefficients-v146}.  Its support ranks are
\((1,M)\).  The relevant strict inequality holds for every pencil,
including singular pencils; it is not a generic-rank condition.

\begin{theorem}[Removal of the projective-bundle restriction]
\label{thm:unrestricted-moving-v146}
In Theorem~\ref{thm:moving-reconstruction-v146}, the hypothesis
that \(\Pj_B(\mathcal U^*)\) is nontrivial can be omitted.  The
coefficient-sub\-bundle reconstruction, the smallest uniform order,
and the automorphism exact sequence remain valid for every
rank-\(n\) source bundle on a smooth connected projective complex
base.  In particular the base is allowed to be a point.
\end{theorem}
\begin{proof}
The first-relation construction, the recovered homogeneous cone,
and cancellation of the determinant use neither projective
nontriviality nor positive dimension of the base.  They recover
both the unordered rulings and the residual coefficient module
before a ruling is chosen.  At this stage apply
Lemma~\ref{lem:coefficient-orientation-v146} with support ranks
\((1,M)\).  It excludes the transposed alternative and hence
orients the left ruling even when both ruling bundles are trivial.
This replaces exactly the nontrivial-projectivization step in the
proof of Theorem~\ref{thm:moving-reconstruction-v146}.

The left ruling is the trivial bundle \(\Pj(V)\times B\).
Projectivity and connectedness still make its projective-matrix
map constant.  Choose its constant linear lift; uniqueness of the
right bundle map on overlaps gives the same actual isomorphism
\(h\), and coefficient contraction recovers the entire pencil
subbundle.  The converse uses the same homogeneous map.  The
lower-order algebras and the two inequivalent constant pencils in
Corollary~\ref{cor:moving-sharpness-v146} are unchanged.  Finally,
the proof of Theorem~\ref{thm:automorphism-kernel-v146} uses only
the resulting classification of graded data, not the earlier way
of orienting the rulings, so it too applies without the restriction.
\end{proof}

\subsection{Every pencil from one Artin local algebra}
Choose \(n\)-dimensional vector spaces \(U,V\).  Let
\(T=(t_{ij})\) be the universal matrix, and put
\(\mathfrak m=(t_{ij})\).  For a pencil
\(R\subset\Sym^2V\), let \(\gamma_R\) be its quotient map and set
\begin{equation}\label{eq:local-algebra-v146}
 \mathfrak A_R^{[k]}=
 \C[t_{ij}]/\bigl((\det T)I_p(\gamma_R\Sym^2T)
                                      +\mathfrak m^{k+1}\bigr).
\end{equation}
An isomorphism below means an isomorphism of complex algebras,
without a grading, tangent coordinates, matrix decomposition or
chosen generators.

\begin{theorem}[Unmarked local reconstruction]
\label{thm:artin-local-inverse-v146}
For every \(n\ge3\) and every pair of quadratic pencils,
\begin{equation}\label{eq:artin-local-inverse-v146}
 \mathfrak A_R^{[d]}\simeq\mathfrak A_{R'}^{[d]}
 \quad\Longleftrightarrow\quad
 R'=(\Sym^2g)R\quad\text{for some }g\in\operatorname{GL}(V).
\end{equation}
The algebra has embedding dimension \(n^2\) and length
\(\binom{n^2+d}{d}-\binom N2\).  Every lower-order algebra
\(\mathfrak A_R^{[k]}\), \(k<d\), is independent of \(R\).
Thus \(d=n^2+2n-4\) is the smallest uniform reconstruction order
already among these single-point objects.
\end{theorem}
\begin{proof}
Theorem~\ref{thm:unrestricted-moving-v146} applies with
\(B=\Spec\C\) and the bundle \(U\).  To spell out what its
intrinsic construction uses in this case, the unique maximal ideal
\(\mathfrak n\) of the finite algebra determines its tangent
space and the multiplication kernel
\[
 K_R=\ker\left[\Sym^d(\mathfrak n/\mathfrak n^2)
                                      \longrightarrow\mathfrak n^d\right].
\]
This is precisely the degree-\(d\) ideal in the chosen presentation;
nonlinear changes of generators contribute only terms above degree
\(d\) to a degree-\(d\) relation.  Hence an ungraded algebra
isomorphism identifies the homogeneous cones generated by these
kernels.  Their determinant reductions recover the unordered
matrix factors.  The residual support ranks \((1,M)\) rule out
transposition, and the remaining left-right map carries the left
coefficient line to the left coefficient line.  Exterior duality
then gives \(R'=(\Sym^2g)R\).  Conversely such a map, together
with any chosen right identification, induces the algebra
isomorphism.  There are no relations below degree \(d\), so the
embedding dimension and the lower-order assertion are immediate.
The relation subspace at degree \(d\) has dimension \(M\), giving
the asserted length.  The inequivalent common-factor and
no-common-factor pencils prove uniform minimality as before.
\end{proof}

In particular the earlier Schubert-line construction is a sufficient
geometric realization, not a lower bound on the dimension of a
supporting reduced scheme.  The local inverse just proved retains no
curve at all.  Its additional ingredient is the coefficient-support
asymmetry, which was not used in the earlier geometric orientation
argument.  This is a stronger inverse statement, not a conclusion
from knowing the discriminant or the reduced determinant cone.

\subsection{The local automorphism kernel is explicit}
Put \(a=n^2\), \(\ell=\length\mathfrak A_R^{[d]}\), and
\(G_R=\{g\in\operatorname{GL}(V):(\Sym^2g)R=R\}\).
The symbol \(\mathfrak n^2\) in the next statement is the square
of the maximal ideal of the finite algebra.

\begin{proposition}[Tangent-identity automorphisms]
\label{prop:local-automorphisms-v146}
There is a split exact sequence of complex automorphism groups
\begin{equation}\label{eq:local-aut-sequence-v146}
 1\longrightarrow U_R\longrightarrow
 \operatorname{Aut}_{\C}(\mathfrak A_R^{[d]})
 \longrightarrow (G_R\times\operatorname{GL}(U))/\C^*
 \longrightarrow1,
\end{equation}
where the scalar subgroup is diagonal.  The kernel \(U_R\) is a
unipotent algebraic group whose underlying variety is an affine
space of dimension
\[
                            a(\ell-1-a).
\]
More explicitly, after choosing the \(a\) degree-one generators
\(x_1,\ldots,x_a\), its elements are exactly the substitutions
\(x_i\mapsto x_i+u_i\) with arbitrary \(u_i\in\mathfrak n^2\).
The affine-space parametrization is not an assertion that this
group law is additive.
\end{proposition}
\begin{proof}
The local inverse and the support-rank argument identify the linear
stabilizer of \(K_R\) with the indicated left-right quotient;
transposition is absent.  These linear maps lift homogeneously,
giving the quotient and splitting.  An automorphism in the kernel
is the identity on \(\mathfrak n/\mathfrak n^2\), and consequently
has exactly the displayed generator images.

Conversely choose arbitrary \(u_i\in\mathfrak n^2\).  Under this
substitution a homogeneous defining relation of degree \(d\)
changes only by terms in \(\mathfrak n^{d+1}=0\), so all defining
relations are respected.  The induced map is the identity on the
associated graded algebra, hence is invertible as a finite filtered
linear map and therefore as an algebra map.  There are no further
conditions on the \(u_i\).  Reading the generator images and
performing substitutions are polynomial mutually inverse maps
between the kernel and the vector space
\((\mathfrak n^2)^{\oplus a}\).  Since
\(\dim\mathfrak n^2=\ell-1-a\), this proves the variety and
dimension assertions.  On a basis ordered by degree every such
automorphism is upper unitriangular; the algebraic group is
therefore unipotent.  Its law includes higher substitution terms,
which accounts for its generally noncommutative structure.
\end{proof}

The positive-dimensional family statement remains useful after the
local inverse: it recovers an entire varying subbundle and the
actual source bundle up to one global left transformation, whereas
separate fibrewise local isomorphisms need not glue.  The geometric
automorphism statements here still do not identify unmarked moduli
stacks or all deformations over nonreduced parameter bases.

\section{The intrinsic polarization and a projective moduli completion}
\label{sec:polarization-v154}
Put $P=\Gr(2,\Sym^2V)$, $N=\binom{n+1}{2}$,
$M=\binom N2$, and $D=n^2+2n-4$. Fix temporarily an auxiliary
$n$-space $U$ and put $E=\Hom(U,V)^*$. Write $L=\OO_P(1)$ for
the Pluecker line and $\Gamma=\Phi(P)\subset\Gr(M,\Sym^D E)$ for
the closed image of Theorem~\ref{thm:finite-parameter-immersion}.
The group in this section is $\operatorname{SL}(V)$, acting by
congruence on pencils and by its induced action on $E$, with $U$
fixed. It is not the full group $\operatorname{GL}(E)$.

\begin{theorem}[Polarized completion of the effective failure moduli]
\label{thm:polarized-completion-v154}
There is an $\operatorname{SL}(V)$-equivariant isomorphism
\begin{equation}\label{eq:polarization-v154}
 \Phi^*\OO_{\Gr(M,\Sym^D E)}(1)\simeq L^{\otimes M}
                 \otimes C_U,
\end{equation}
where $C_U$ is a constant one-dimensional space on which
$\operatorname{SL}(V)$ acts trivially. Consequently the stable and
semistable loci of $\Gamma$, for the restricted Pluecker
linearization, are respectively the images of $P^s$ and $P^{ss}$.
The quotient defined by the complete restricted section ring is
canonically
\begin{equation}\label{eq:git-completion-v154}
 \overline{\mathcal M}^{\mathrm{fail}}_n
 =\Proj\!\left(\bigoplus_{q\geq0}
                  H^0(P,L^{\otimes q})^{\operatorname{SL}(V)}\right)
 \simeq\Gamma/\!/\operatorname{SL}(V).
\end{equation}
It is a normal projective variety. It contains the geometric quotient
of the simple-discriminant pencils as a dense open subvariety, and
has dimension $n-3$. The complement of that open is the image of the
semistable discriminant-zero locus.

The intrinsic inverse identifies the effective semistable failure
stack with $[P^{ss}/\PGL(V)]$ and gives a moduli morphism to
\eqref{eq:git-completion-v154}. On framed charts the failure-algebra
family extends over all $P$, remains finite locally free, and
commutes with arbitrary base change. These assertions do not require,
and do not assert, the existence of a universal ordinary algebra on
the coarse quotient or properness of the entire quotient stack.
\end{theorem}
\begin{proof}
Let $\mathcal K$ be the tautological relation bundle pulled back by
$\Phi$. The coefficient factorization
\eqref{eq:finite-coefficient-factorization} identifies it with
$\mathcal B\otimes F(U)$, where $\mathcal B$ is the quotient-volume
coefficient line. Exterior duality gives
$\mathcal B\simeq L^{-1}$ up to a fixed determinant character of
$\operatorname{GL}(V)$. That character is trivial on
$\operatorname{SL}(V)$. Taking the dual determinant gives
\eqref{eq:polarization-v154}, with the harmless constant factor
$(\det F(U))^*$ and the corresponding fixed determinant line of $V$.
A trivialization of the constant factor identifies the section rings;
a different trivialization rescales each graded piece and has no
effect on its projective spectrum. The section ring for $L^M$ is
the $M$th Veronese of the section ring for $L$. Taking invariants
commutes with this operation. The Hilbert--Mumford weights are
multiplied by $M>0$, so stability and semistability are unchanged.
This also explains why one uses the complete section ring on
$\Gamma$, rather than assuming surjectivity of unspecified ambient
invariant restrictions.

Finite generation and the projective good quotient are the reductive
GIT construction of \cite[Chapter~1]{MFK1994}. The complete section
ring of the normal projective variety $P$ is integrally closed: it
is the intersection of the valuation conditions for rational sections
of powers of $L$ at the prime divisors of $P$. Its invariant subring
is integrally closed as well. Indeed an element of its fraction
field integral over the invariants is integral over the original
ring, belongs to that ring by normality, and is invariant. This proves
normality of the projective quotient.

For completeness, the simple-discriminant open is nonempty and
stable. In a frame $(Q_0,Q_1)$ of a pencil, take the discriminant of
$\det(sQ_0+tQ_1)$. Under a frame change $h\in\operatorname{GL}_2$
it has weight $(\det h)^{n(n-1)}$, and it is invariant under
$\operatorname{SL}(V)$. Thus it is a nonzero invariant section of
$L^{n(n-1)}$. It is nonzero at a diagonal pencil with $n$ distinct
roots, so its nonvanishing locus is semistable and saturated.
A pencil there has a nonsingular member. Relative to that member,
the other member is self-adjoint with distinct eigenvalues;
orthogonal diagonalization identifies its congruence class with an
unordered configuration of $n$ distinct points of $\Pj^1$.
The latter configurations are stable for the usual binary-form
linearization: the root-multiplicity criterion is strict since
$1<n/2$ for $n\geq3$. Their quotient is geometric and separated.
It follows that the congruence orbits are closed in the saturated
nonvanishing locus. Their stabilizers are finite: the induced
projective transformation permutes a set of at least three points,
and the kernel consists of diagonal sign changes modulo scalars.
They are therefore stable. This gives a geometric dense open of
dimension $n-3$. The complement is the image of the invariant
section's zero locus by the good-quotient property. This reasoning
uses classical pencil classification and GIT, not a new classification
of binary forms; compare \cite{FevolaMandelshtamSturmfels,MFK1994}.

Finally apply Theorem~\ref{thm:pencil-stack-equivalence} and the
rigidification results of Section~\ref{sec:rigidification-v151}.
The Pluecker polarization is taken on the recovered pencil, not on
an arbitrary first-relation orbit. A sufficiently divisible power
of $L$, for example $L^n$, has trivial central character and descends
from $\operatorname{SL}(V)$ to $\PGL(V)$. A change of auxiliary frame
changes the constant factor in \eqref{eq:polarization-v154}, not the
resulting projective quotient. On the framed Grassmannian the extended
algebras are exactly those of Corollary~\ref{cor:finite-flat-family}.
Their equivariance yields the stated classifying morphisms. Removing
ineffective inertia does not turn these into a universal ordinary
algebra on the coarse space, and no such descent is used.
\end{proof}

\begin{remark}[What is compactified]\label{rem:git-scope-v154}
The completion is the classical polarized pencil quotient recovered
from the intrinsic failure invariant, with its polarization verified
by \eqref{eq:polarization-v154}. It is not a new construction of
classical GIT. Its role is to supply a proper coarse target for the
effective invariant and a specified class of semistable limiting
families. Strictly semistable points represent closed-orbit
representatives, not every orbit in their fibres. In particular Segre
data on the stack need not descend to an invariant separating all
points of a coarse-quotient fibre. For $n=3$ the quotient is a point
and its discriminant boundary can be empty. The full
$\operatorname{GL}(E)$ common limit of
Theorem~\ref{thm:universal-boundary-v152} is a different construction.
\end{remark}

\section{Spectral divisor fibres and the recovery of degeneration type}
\label{sec:spectral-divisor-v154}
The preceding completion concerns the native congruence action. To
retain degeneration type before passage to a coarse quotient, we
associate a tower of binary divisor incidences to a regular pencil.
This construction is intrinsic after the unmarked inverse and is
unchanged by a simultaneous change of the pencil and source frames.

Let $\Lambda=\Pj(R)$ be the pencil line and let
$\mathcal T_R$ be the cokernel of the universal symmetric map on
$\Lambda$. A regular pencil means that its determinant is not
identically zero. Thus $\mathcal T_R$ is a torsion sheaf. Define
$D_j(R)$ by
\begin{equation}\label{eq:spectral-Fitting-v154}
 \II_{D_j(R)}=\Fitt_j(\mathcal T_R),\qquad 0\leq j\leq n,
\end{equation}
where $D_n=\varnothing$. These are effective divisors on the smooth
curve $\Lambda$. Write $s_j=\deg D_j$. If $s_j>0$, let $G_j$ be
its binary equation and put
\[
 J_j=G_j H^0(\Lambda,\OO_\Lambda(1))
       \subset H^0(\Lambda,\OO_\Lambda(s_j+1)).
\]
Changing the scalar representative of $G_j$ does not change $J_j$.
Let $\mathcal F_j(R)$ be the whole scheme fibre at $J_j$ of the
degree-one divisor normalization $\nu_1$, with $r=2$ and $D=s_j+1$.
For $s_j=0$ put $\mathcal F_j=\varnothing$.

\begin{theorem}[A fibre-theoretic Segre invariant]\label{thm:spectral-fibres-v154}
There is a natural identification of finite schemes over $\Lambda$
\begin{equation}\label{eq:fibre-divisor-v154}
                        \mathcal F_j(R)\simeq D_j(R).
\end{equation}
At a spectral point $x$, write the positive elementary divisors of
$\mathcal T_R$ as
$\lambda_{x,1}\geq\cdots\geq\lambda_{x,r_x}>0$, and extend this
sequence by zeroes. Then
\begin{equation}\label{eq:fibre-length-v154}
 (\mathcal F_j(R))_x=\Spec\C[\epsilon]/(\epsilon^{b_{x,j}}),
 \qquad b_{x,j}=\sum_{i>j}\lambda_{x,i},
\end{equation}
with the empty-scheme convention when $b_{x,j}=0$. Consequently
\begin{equation}\label{eq:Segre-difference-v154}
                 \lambda_{x,j+1}=b_{x,j}-b_{x,j+1}.
\end{equation}
The tower, with its maps to the same intrinsic spectral line,
therefore determines the full Segre symbol. Conversely, the
elementary-divisor partitions together with their spectral support
on that line determine the tower. The tower retains the projective
configuration of the support as well as the Segre symbol.
It is an invariant of the unmarked finite failure algebra.

The construction is compatible with base change on every family
stratum where the Fitting ideals in \eqref{eq:spectral-Fitting-v154}
define relative effective Cartier divisors of fixed degrees. On such
a stratum $\mathcal F_j$ is finite locally free of degree $s_j$.
These hypotheses are part of the relative statement; no assertion of
flatness across a jump of $s_j$ is intended.
\end{theorem}
\begin{proof}
Over the discrete valuation ring $\OO_{\Lambda,x}$, Smith normal
form gives diagonal entries with exponents
$\lambda_{x,r_x},\ldots,\lambda_{x,1}$ and $n-r_x$ zero
exponents. The minors of size $n-j$ have greatest common divisor
of order $\sum_{i>j}\lambda_{x,i}$. This proves the formula for
$D_j$, including its nonreduced structure. It is the usual
Fitting-ideal description of elementary divisors, used here on the
intrinsic spectral line.

The residual two-space $H^0(\OO_\Lambda(1))$ is gcd-free. Hence
Theorem~\ref{thm:binary-fibres-v153} identifies the degree-one
normalization fibre at $J_j$ with the scheme of linear factors of
$G_j$. On an affine chart write a monic linear factor as $z-a$.
Monic division gives $G_j=(z-a)H+G_j(a)$, so the entire factor
scheme is $G_j(a)=0$, not its reduction. These descriptions agree
on overlaps of the projective line. They give
\eqref{eq:fibre-divisor-v154} and the local algebra
\eqref{eq:fibre-length-v154}. Taking successive differences proves
\eqref{eq:Segre-difference-v154}. The common map to $\Lambda$ is
essential: forgetting it separately in each member of the tower
would discard the alignment of the spectral supports.

Pencil frame changes act on $\Lambda$ by projective transformations;
source congruence changes give isomorphic cokernels and preserve all
Fitting ideals. The intrinsic inverse of
Theorem~\ref{thm:artin-local-inverse-v146} therefore defines this
invariant from the unmarked algebra. The equivalence with Segre data
uses the classical elementary-divisor convention of
\cite[\S5]{FevolaMandelshtamSturmfels}.

For a family with the stated Cartier and degree conditions, the same
monic division argument works over any base algebra, including one
with nilpotents. It represents the divisor's functor of points, so
it commutes with base change. A relative effective divisor of fixed
degree on a projective-line bundle is finite locally free of that
degree. Local choices of $G_j$ differ by units and glue. Formation
of Fitting ideals itself commutes with base change; requiring them
to remain flat Cartier divisors is precisely what permits the
additional finite-flat conclusion.
\end{proof}

\begin{example}[Equal discriminant, different fibre towers]
\label{ex:segre-fibres-v154}
At a root of multiplicity four, the partitions $(3,1)$ and $(2,2)$
have the same zeroth fibre $\C[\epsilon]/(\epsilon^4)$.
Their first fibres are respectively $\C$ and
$\C[\epsilon]/(\epsilon^2)$; all higher fibres are empty there.
Thus the discriminant alone does not detect the distinction, whereas
the divisor-fibre tower does. Both are realized by regular symmetric
pencils: for a block of size $l$, take the reversal matrix $H_l$ and
the symmetric pair $(H_l,H_lJ_l(\alpha))$, where $J_l(\alpha)$ is
a Jordan block. Direct sums realize every indicated partition.
Additional simple blocks may be appended. The example is a statement
on the regular-pencil stack, not a claim that distinct Segre strata
always remain distinct after strictly semistable GIT identification.
\end{example}

\begin{remark}[The two divisor constructions]\label{rem:two-divisors-v154}
The primitive first relation of an actual pencil has exact common
factor $(\det T)^{n-1}$, by Lemma~\ref{lem:primitive-pencil}.
Its common-factor degree does not jump when the pencil changes
Segre type. The divisors $D_j$ above instead lie on the recovered
spectral line and define auxiliary binary first-relation subspaces
$J_j$. Thus \eqref{eq:Segre-difference-v154} relates pencil data to
intrinsically induced divisor ramification, not to a nonexistent
change in the conductor of the original exact-gcd locus. The
identification supplies a moduli consequence of the all-binary-fibre
theorem while respecting this distinction.
\end{remark}

\part*{Part II. Divisor incidence and homological fibres}

% BEGIN PRESERVED PART 37-common-divisor-strata-v149.tex
\section{Common-divisor strata and intrinsic families}
\label{sec:common-divisor-strata}

The coefficient condition in Theorem~\ref{thm:recognition-dichotomy-v148}
is not imposed in this section.  We instead retain the common divisor
of a space of first relations.  This geometric condition gives a
larger smooth parameter space, includes relations without either
factor symmetry, and has a scheme-theoretic inverse.  The distinction
between a stratum and the condition on its geometric points is
essential when the parameter base is nonreduced.

Write \(S_j(E)=\Sym^j E\), where \(E\) is a complex vector space.
For \(r\ge2\), denote by
\[
 X_{h,r}(E)\subset\Gr(r,S_h(E))
\]
the locus of subspaces with greatest common divisor equal to \(1\),
up to a scalar.  Empty Grassmannians are omitted throughout.  When
\(h=0\) this notation can also be used for \(r=1\).

\begin{lemma}[Universal common divisor]\label{lem:universal-gcd}
Let \(D=g+h\).  The spaces of \(r\) degree-\(D\) forms whose exact
greatest common divisor has degree \(g\) form a locally closed
reduced subscheme \(Z_g\) of \(\Gr(r,S_D(E))\).  Multiplication
induces an isomorphism
\begin{equation}\label{eq:universal-gcd}
 \Pj(S_g(E))\times X_{h,r}(E)\ \xrightarrow{\ \sim\ }\ Z_g,
                 \qquad ([f],J)\longmapsto fJ.
\end{equation}
In particular, the common-divisor line and the residual subbundle of
a family with classifying map in \(Z_g\) are recovered uniquely and
commute with arbitrary complex base change.
\end{lemma}
\begin{proof}
For each \(a\), multiplication from
\(\Pj(S_a(E))\times\Gr(r,S_{D-a}(E))\) has closed image: its domain
is projective.  The union for \(a>g\) is closed.  Remove this union
from the image for \(a=g\), and give the resulting locally closed
set its reduced structure.  Its points are exactly those in the
statement.  The inverse image of this set under multiplication is
\(\Pj(S_g(E))\times X_{h,r}(E)\).  The same closed-image argument
with \(D=h\) shows that \(X_{h,r}(E)\) is open.

The restricted multiplication map is proper, by base change.  It
is injective on geometric points over every algebraically closed
extension of \(\C\), since a greatest common divisor is unique up
to scalar in a polynomial ring.  Its differential is injective as
well.  Indeed, a tangent vector in its kernel satisfies
\[
                   \dot f\, j+f\dot j\in fJ\quad(j\in J).
\]
Consequently \(f\mid\dot f\,j\) for every \(j\in J\).  For each
irreducible factor of \(f\), some element of \(J\) is not divisible
by that factor.  Comparison of its valuation gives
\(f\mid\dot f\).  Equality of degrees makes \(\dot f\) a scalar
multiple of \(f\); then \(\dot J=0\) in
\(\Hom(J,S_h(E)/J)\).

In characteristic zero the geometric-point injectivity and this
differential calculation imply that the map is universally
injective and unramified.  A proper, universally injective,
unramified morphism is a closed immersion
\cite[Lemma~41.7.2]{StacksClosedImmersion}.
Here its image is all of the reduced target.  The defining ideal
of a closed immersion with this support is zero in a reduced
scheme.  This proves the isomorphism, including its assertion for
nonreduced test schemes.  Equivalently, one can check properness
on the base change of the original projective multiplication map;
that base change has the same underlying open preimage.  Its source
is reduced because it is this open subscheme of the original
smooth product.
\end{proof}

\begin{remark}\label{rem:gcd-scheme-condition}
A map from a nonreduced scheme is required to factor through the
locally closed \emph{scheme} \(Z_g\), not merely to have geometric
points in its underlying set.  These are different requirements.
For example, over \(\C[\epsilon]/(\epsilon^2)\) the span of
\(x^2+\epsilon y^2,xy\) specializes to a pair with common divisor
\(x\), but has no factor \(x+\epsilon(ax+by)\).  Comparison of the
\(y^2\) coefficient in the first form proves this.  The stratum in
Lemma~\ref{lem:universal-gcd} excludes this transverse infinitesimal
family.  No use of a fibrewise radical can replace its scheme structure.
\end{remark}

Now let \(\dim V=\dim U=n\ge3\),
\(E=V^*\otimes U\), and \(\delta=\det T\).  Denote by
\(G\subset\operatorname{GL}(E)\) the stabilizer of the line
\(\C\delta\).  By Lemma~\ref{lem:linear-preserver-v147},
\begin{equation}\label{eq:determinant-group}
 G^\circ=(\operatorname{GL}(V)\times\operatorname{GL}(U))/\mathbb G_m,
                 \qquad G/G^\circ=\mathbf Z/2\mathbf Z.
\end{equation}
The second component exchanges the two tensor factors; the diagonal
scalars in the displayed product act trivially on the matrix space.
Fix \(s\ge1\) and \(D=sn+h\).  Inside \(Z_{sn}\), impose the condition
that the common-divisor line is the \(s\)-th power of a determinant
form after a linear change of the \(n^2\) variables.  Call the resulting
locally closed scheme \(Z^{\det}_{s,h,r}\).

\begin{theorem}[Intrinsic determinant-divisor stratum]
\label{thm:determinant-stratum}
There are canonical isomorphisms of schemes and algebraic stacks
\begin{equation}\label{eq:determinant-stratum}
 Z^{\det}_{s,h,r}\simeq
 \operatorname{GL}(E)\times^G X_{h,r}(E),\qquad
 [Z^{\det}_{s,h,r}/\operatorname{GL}(E)]\simeq[X_{h,r}(E)/G].
\end{equation}
They hold over arbitrary complex parameter schemes.  The stack is
smooth.  At a residual relation \(J\), its tangent complex, in degrees
\(-1,0\), is
\begin{equation}\label{eq:ambient-tangent-complex}
 [\,\operatorname{Lie}(G)\longrightarrow
                    \Hom(J,S_h(E)/J)\,],
             \qquad X\longmapsto(j\mapsto Xj\bmod J).
\end{equation}
No invariance of \(J\) under either full factor group is assumed.
The equivalence identifies stabilizers, infinitesimal families, and
specializations whose classifying maps remain in this stratum.
\end{theorem}
\begin{proof}
The power map \(\Pj(S_n(E))\to\Pj(S_{sn}(E))\),
\([f]\mapsto[f^s]\), is a closed immersion in characteristic zero.
It is proper and geometrically injective, and its differential has
kernel given by \(sf^{s-1}\dot f\in\C f^s\), hence by
\(\dot f\in\C f\).  The same closed-immersion criterion as above
applies.  The orbit of \([\delta]\) is the smooth locally closed
homogeneous space \(\operatorname{GL}(E)/G\).  The preimage of its
power orbit under the first projection in
\eqref{eq:universal-gcd} is therefore a locally closed subscheme.
It is exactly \(Z^{\det}_{s,h,r}\).

Pulling back along the \(G\)-torsor
\(\operatorname{GL}(E)\to\operatorname{GL}(E)/G\) trivializes the
common-divisor form, leaving the unrestricted residual relation
\(J\in X_{h,r}(E)\).  Descent gives the first isomorphism in
\eqref{eq:determinant-stratum}; descent of equivariant maps and
morphisms gives the second.  This proof uses morphisms of schemes,
not a comparison only on complex points.  It also supplies local
lifts of all isomorphisms over an fppf cover.

The space \(X_{h,r}(E)\) is open in a Grassmannian and \(G\) is
smooth.  The quotient is algebraic and smooth
\cite[Theorem~97.17.2]{StacksQuotient}; see also
\cite[Lemma~101.33.7]{StacksSmoothQuotient}.
For completeness, write a first-order residual subspace as the
graph of \(\epsilon\phi:J\to S_h(E)/J\).  An infinitesimal change
\(1+\epsilon X\) replaces \(\phi\) by \(\phi+X|_J\bmod J\).
This proves the tangent complex.  Its kernel records infinitesimal
automorphisms and its cokernel records first-order isomorphism
classes.  The same torsor description proves compatibility with
specialization.  No assertion is made about limits outside the
specified common-divisor stratum.
\end{proof}

The next step connects these parameter spaces with ungraded finite
algebras.  It is useful to state the family condition intrinsically,
since an absolute nilradical is unsuitable over a nonreduced base.
Put \(e=\dim E\) and
\(\ell=\binom{e+D}{D}-r\).  Consider finite locally free commutative
unital algebras \(\mathcal A\) of rank \(\ell\) on a complex scheme
\(T_0\), subject to the following conditions.  The map
\(\epsilon(a)=\ell^{-1}\operatorname{tr}(M_a)\) is multiplicative.
For \(\mathcal N=\ker\epsilon\), one has
\(\mathcal N^{D+1}=0\); all quotients of its power filtration are
locally free; \(\mathcal E=\mathcal N/\mathcal N^2\) has rank \(e\);
and multiplication gives isomorphisms
\(\Sym^j\mathcal E\simeq\mathcal N^j/\mathcal N^{j+1}\) for
\(j<D\).  In degree \(D\) its kernel has rank \(r\).  These
conditions specify the maximal-lower-Hilbert-function stratum; they
do not assert that arbitrary smoothings of a local algebra remain
in it.

\begin{proposition}[Relative first-relation algebra]
\label{prop:relative-first-relation}
Every such algebra is locally isomorphic, over its base, to
\[
 \mathcal A(\mathcal K)=
 \Sym\mathcal E/((\mathcal K)+\mathfrak m^{D+1}),\qquad
 \mathcal K=\ker(\Sym^D\mathcal E\longrightarrow\mathcal N^D).
\]
Its trace augmentation, \(\mathcal E\), and \(\mathcal K\) are
intrinsic and compatible with arbitrary base change in this
stratum.  For two such algebras, taking the linear part gives a
locally surjective map of isomorphism sheaves to the sheaf of
isomorphisms of \((\mathcal E,\mathcal K)\).  Its kernel is the
smooth substitution group whose underlying scheme in a local
presentation is
\(\Hom(\mathcal E,\mathcal N^2)\).

After quotienting morphisms by this tangent-identity kernel and
fppf stackifying, the algebras with first-relation map in
\(Z^{\det}_{s,h,r}\) have stack \([X_{h,r}(E)/G]\).
\end{proposition}
\begin{proof}
Locally lift a basis of \(\mathcal E\) to \(\mathcal N\).
Nilpotence and multiplication on the successive quotients make the
induced map from the truncated symmetric algebra surjective.
A relation with a nonzero term of least degree \(j<D\) would give
a nonzero kernel in degree \(j\) of the associated graded algebra,
which is impossible.  Its entire kernel therefore lies in degree
\(D\) and is precisely \(\mathcal K\).  Changing a lift by an
element of \(\mathcal N^2\) does not alter a product of \(D\)
generators, since the difference has degree at least \(D+1\).
This also proves the asserted uniqueness of the first relation.

In this presentation multiplication by a positive-degree element
is strictly filtration-increasing.  Its trace is zero over any
base ring.  Multiplication by a scalar has trace \(\ell\) times
that scalar, proving that the presentation's augmentation is the
specified intrinsic trace augmentation.  All filtration sequences
are locally split sequences of vector bundles, so their kernels
and quotients commute with base change.

An isomorphism must preserve trace and induces a linear map carrying
\(\mathcal K\) to \(\mathcal K'\).  Conversely such a map lifts
locally to an isomorphism by the homogeneous presentations.  All
other lifts are obtained by substituting generators plus arbitrary
elements of \(\mathcal N'^2\); the inverse is constructed degree by
degree in a finite filtration.  This gives the stated kernel and
its smoothness, without replacing its substitution law by addition.
After identifying morphisms with the same linear part, frames of
\(\mathcal E\) give the quotient
\([Z^{\det}_{s,h,r}/\operatorname{GL}(E)]\).
Theorem~\ref{thm:determinant-stratum} completes the proof.
\end{proof}

\begin{remark}\label{rem:algebra-inertia}
The original algebra stack is not this quotient stack: it has the
additional nonlinear inertia just computed.  The phrase
\emph{effective first-relation stack} will refer to the specified
quotient on morphisms followed by stackification.  The construction
retains the full automorphism group before taking that quotient.
It does not identify algebras with the same geometric fibres, nor
does it discard nilpotent parameters.  In particular, its relation
to a stratum of first relations is stronger than a bijection of
isomorphism classes over \(\C\).
\end{remark}

% BEGIN PRESERVED PART 40-canonical-divisor-boundary-v151.tex
\section{The common-divisor boundary and its normalization}
\label{sec:canonical-boundary-v151}

The exact-common-divisor stratum has a canonical completion which need
not be normal.  Its normalization records a divisor, and its failure to
be an isomorphism is controlled by two different phenomena: choices of
that divisor and overlap with the residual common factor.  We establish
these assertions for arbitrary spaces of forms.  No determinant, Cauchy,
or Pluecker equation is imposed in this section.

Let \(E\) have dimension \(e\geq2\), put \(S_j=\Sym^j E\) and
\(s_j=\dim S_j=\binom{e+j-1}{j}\), and fix
\[
             g,h\geq1,\qquad D=g+h,\qquad 2\leq r\leq s_h.
\]
Write \(\mathscr G=\Gr(r,S_D)\).  A \emph{degree-\(g\) divisor
incidence} over a complex scheme \(T_0\) consists of a line subbundle
\(\mathcal L\subset S_g\otimes\OO_{T_0}\) and a rank-\(r\)
subbundle \(\mathcal K\subset S_D\otimes\OO_{T_0}\) such that
\begin{equation}\label{eq:divisor-incidence-v151}
 \mathcal K\longrightarrow
 (S_D\otimes\OO_{T_0})/(\mathcal L\cdot S_h)=0.
\end{equation}
Subbundle always means that the quotient is locally free.  In
particular, the form generating \(\mathcal L\) is nonzero on every
geometric fibre.  Condition~\eqref{eq:divisor-incidence-v151} is the
vanishing of a bundle map, not a condition imposed only on closed points.

\begin{proposition}[Canonical divisor incidence]
\label{prop:incidence-v151}
The incidence functor is represented, with its full scheme structure, by
\[
 \mathscr X_g=\Pj(S_g)\times\Gr(r,S_h).
\]
Its forgetful morphism is multiplication
\begin{equation}\label{eq:normalization-map-v151}
 \nu_g:\mathscr X_g\longrightarrow\mathscr G,
                   \qquad ([f],J)\longmapsto fJ.
\end{equation}
This morphism is finite.  Its scheme-theoretic image \(Y_g\) is an
integral projective scheme.  If \(K\) is a geometric point of \(Y_g\)
and \(F_K=\gcd(K)\), the underlying set of \(\nu_g^{-1}(K)\) is
\begin{equation}\label{eq:divisor-fibre-set-v151}
       \{[f]: f\mid F_K,\ \deg f=g\},
\end{equation}
with residual space \(J=K/f\).  The incidence construction and its
scheme-theoretic fibres commute with arbitrary complex base change.
\end{proposition}
\begin{proof}
On \(\Pj(S_g)\), multiplication from the tautological line tensored
with \(S_h\) to \(S_D\) has constant rank \(s_h\).  Indeed,
multiplication by a nonzero polynomial over a field is injective.
A nonzero maximal minor locally splits this map, so its image and
cokernel are vector bundles, also after a nonreduced base change.
The closed incidence~\eqref{eq:divisor-incidence-v151} is therefore
the relative Grassmannian of rank-\(r\) subbundles of
\(\mathcal L\otimes S_h\).  Tensoring by \(\mathcal L^{-1}\)
identifies it canonically with the stated product.  In more elementary
terms, \(\mathcal K\) has a unique inverse image in
\(\mathcal L\otimes S_h\); its quotient there is locally free
because the quotient of \(S_D/\mathcal K\) onto
\(S_D/(\mathcal L\cdot S_h)\) is a surjection of vector bundles.
This proves the functorial assertion, including its scheme structure.

The source of \(\nu_g\) is projective.  Unique factorization in the
polynomial ring over an algebraically closed field gives
\eqref{eq:divisor-fibre-set-v151}.  A fixed polynomial has only finitely
many divisors of a fixed degree, up to scalar.  The morphism is thus
proper with finite fibres and is finite
\cite[Tag~02LS]{StacksFiniteV151}.
Define \(Y_g\) by the kernel of
\(\OO_{\mathscr G}\to(\nu_g)_*\OO_{\mathscr X_g}\), as in
\cite[Tag~01R5]{StacksImageV151}.  Its integrality follows from the
integrality of \(\mathscr X_g\): on an affine open meeting the image,
the coordinate ring of the image injects into the domain which is the
coordinate ring of its finite inverse image.  Reducedness is therefore
a consequence of this incidence construction, not an extra choice.
Representability of incidence gives the base-change assertion.
\end{proof}

\begin{remark}\label{rem:degree-divisor-v151}
The condition defining \(Y_g\) is the existence of a divisor of
\emph{degree \(g\)}.  In several variables it is not the condition
\(\deg\gcd(K)\geq g\): an irreducible polynomial of larger degree
need not have a degree-\(g\) factor.  Nor do we assert that formation
of the scheme-theoretic image, or normalization itself, commutes with
arbitrary nonflat base change.  The assertion for arbitrary bases in
Proposition~\ref{prop:incidence-v151} concerns the incidence functor
and the base change of its representing morphism.
\end{remark}

We next retain the nilpotents of a fibre.  This will distinguish a
ramified divisor choice from several distinct divisor choices.

\begin{lemma}[Fibre equations and the overlap space]
\label{lem:fibre-equations-v151}
Let \(K=fJ\), set \(H=\gcd(J)\), and put \(c=\gcd(f,H)\) and
\(k=\deg c\).  Then the tangent space of the scheme-theoretic fibre
at \(([f],J)\) is naturally, after the displayed representatives have
been chosen,
\begin{equation}\label{eq:overlap-kernel-v151}
 \ker(d\nu_g)_{([f],J)}\simeq S_k/\C c,
                     \qquad \dim\ker(d\nu_g)=s_k-1.
\end{equation}
More explicitly, choose a projective coefficient chart
\(f(u)=f+\sum u_\alpha m_\alpha\), with one coefficient normalized
to its value in \(f\).  Let \(M(f(u))\) be the matrix of
multiplication \(S_h\to S_D\).  Choose a set \(A_0\) of \(s_h\)
rows for which \(M_{A_0}(f)\) is invertible, and let \(B_0\) denote
the remaining rows.  For a matrix \(C\) whose columns form a basis
of \(K\), the full fibre on this chart is defined by
\begin{equation}\label{eq:fibre-schur-v151}
 C_{B_0}-M_{B_0}(f(u))M_{A_0}(f(u))^{-1}C_{A_0}=0,
             \qquad \det M_{A_0}(f(u))\ne0.
\end{equation}
In particular its completed local ring is the power-series ring in
the \(s_g-1\) variables \(u_\alpha\), modulo the entries in
\eqref{eq:fibre-schur-v151}; no radical is taken.
\end{lemma}
\begin{proof}
For each column of \(C\), its possible quotient by \(f(u)\) is forced
by the \(A_0\) rows to be \(M_{A_0}^{-1}C_{A_0}\).  The remaining
rows are exactly the condition that this quotient works.  These
quotients span a subbundle wherever the equations hold, by the
incidence argument above.  Thus elimination has introduced neither
additional solutions nor discarded nilpotents, proving the fibre
presentation over arbitrary coefficient rings in the chart.

A tangent vector is a pair \((\dot f\bmod\C f,\dot J)\).  It lies
in the fibre precisely when
\[
              \dot f\,j+f\dot j\in fJ\quad\hbox{for all }j\in J.
\]
It is necessary and sufficient that \(f\mid\dot f\,j\) for every
\(j\in J\); the class of \(\dot j\) modulo \(J\) is then forced.
Taking the minimum valuation over a basis of \(J\), for each
irreducible polynomial, makes this equivalent to \(f\mid\dot f H\).
Writing \(f=cf_1\), \(H=cH_1\), with \(\gcd(f_1,H_1)=1\), gives
\(\dot f=f_1w\) with \(w\in S_k\).  The scalar tangent direction
\(\dot f\in\C f\) is exactly \(w\in\C c\).  This proves
\eqref{eq:overlap-kernel-v151} and its dimension assertion.
\end{proof}

\begin{theorem}[Normalization and the exact normal locus]
\label{thm:boundary-normalization-v151}
The finite map
\[
 \nu_g:\Pj(S_g)\times\Gr(r,S_h)\longrightarrow Y_g
\]
is the normalization of the natural degree-\(g\) divisor locus.
Its dimension is
\begin{equation}\label{eq:boundary-dimension-v151}
                    \dim Y_g=s_g-1+r(s_h-r).
\end{equation}
For a geometric point \(K\), the scheme \(Y_g\) is normal at \(K\)
if and only if \(F_K\) has exactly one degree-\(g\) divisor \([f]\)
and
\begin{equation}\label{eq:normal-criterion-v151}
                         \gcd(f,F_K/f)=1.
\end{equation}
Normality at a point is also equivalent to smoothness there.
Thus the two precise obstructions are more than one divisor choice,
or overlap between the selected divisor and the residual common
factor.  The latter gives the fibre tangent dimension in
\eqref{eq:overlap-kernel-v151}, and
\eqref{eq:fibre-schur-v151} determines its higher nilpotents.
\end{theorem}
\begin{proof}
The locus of gcd-free \(r\)-planes in \(S_h\) is nonempty: a space
containing \(x_1^h,x_2^h\) is one example, and it can be extended
to every allowed rank \(r\).  It is open, since the loci possessing
a positive-degree divisor are the finitely many closed images of
projective multiplication maps.  It is therefore dense in the
irreducible Grassmannian.  Its product with \(\Pj(S_g)\) maps to
the locus where the exact gcd has degree \(g\).

This last locus is open in \(Y_g\): remove the finitely many closed
images for divisor degrees \(a>g\).  Call it \(U_g\).  Every fibre
over a geometric point of \(U_g\) has one point and, by
Lemma~\ref{lem:fibre-equations-v151} with \(H=1\), zero tangent
space.  A finite local algebra over an algebraically closed field
with zero cotangent space is the field itself, by Nakayama's lemma.
The fibre consequently has length one, not merely one point.

Here is the local argument which also settles the scheme structure.
At a closed point of \(U_g\), write \(A\) for the local ring of
\(Y_g\), and \(B\) for its finite inverse-image algebra.  The map
\(A\to B\) is injective, and \(B/\mathfrak m_A B=\C\), generated
by \(1\).  Nakayama's lemma, applied to the finite module \(B/A\),
gives \(B=A\).  The non-isomorphism locus of a finite morphism is
closed, so the same conclusion holds throughout \(U_g\).  Hence
\(\nu_g\) is birational.  The source is smooth and integral, thus
normal.  On affine opens, a finite birational extension \(A\subset B\)
with \(B\) integrally closed is the integral closure of \(A\) in
the common fraction field: every element integral over \(A\) is
integral over \(B\) and hence belongs to \(B\).  This proves the
normalization assertion; compare \cite[Tag~035E]{StacksNormalV151}.
The dimension follows from the product description.

If \(Y_g\) is normal at \(K\), its normalization is an isomorphism
there, so its fibre is one reduced point.  Formula
\eqref{eq:divisor-fibre-set-v151} gives uniqueness of \([f]\), and
\eqref{eq:overlap-kernel-v151} gives \(s_k-1=0\), equivalent to
\(k=0\) because \(e\geq2\).  Conversely uniqueness and \(k=0\)
again make the entire fibre a single reduced point.  The same finite
local Nakayama argument shows that the normalization is an isomorphism
near \(K\).  Its source is normal, so \(Y_g\) is normal at \(K\).
The source is smooth, so every normal point of the target is smooth;
the converse holds for every complex variety.  This proves both
directions without replacing a nonreduced fibre by its set of points.
\end{proof}

\begin{corollary}[Universality of the exact stratum]
\label{cor:canonical-exact-v151}
The open subscheme \(U_g\subset Y_g\) represents the unmarked
relation families in the divisor incidence whose residual space is
gcd-free on every geometric fibre.  Over it the divisor line and
residual subbundle are unique and commute with arbitrary base change.
Moreover \(U_g\), as a locally closed subscheme of \(\mathscr G\),
is the reduced exact-gcd stratum of Lemma~\ref{lem:universal-gcd}.
\end{corollary}
\begin{proof}
The inverse image of \(U_g\) is
\(\Pj(S_g)\times X_{h,r}(E)\), and the proof above shows that its
finite map to \(U_g\) is an isomorphism.  Its source is reduced,
so \(U_g\) has the reduced induced structure on its locally closed
support.  The incidence functor supplies the asserted universal
property.  A family in the ambient Grassmannian which only has the
right geometric gcd need not factor through \(Y_g\).  In particular,
the transverse dual-number family of
Remark~\ref{rem:gcd-scheme-condition} is not admitted by this
corollary.  The corollary is a scheme-theoretic identification, not a
claim that fibrewise gcd tests detect nilpotent equations.
\end{proof}

\begin{example}[A nonreduced fibre at a divisor collision]
\label{ex:fat-divisor-fibre-v151}
Take \(E=\langle x,y,z\rangle\), \(g=1,h=2,r=2\), and
\(K=\langle x^2y,x^2z\rangle\).  Its gcd is \(x^2\), so the only
point of the normalization fibre has divisor \(x\).  In the chart
\(f=x+ay+bz\), division by this monic linear form amounts to
substitution \(x=-ay-bz\).  The two remainders are
\[
             (ay+bz)^2y,\qquad (ay+bz)^2z.
\]
Their coefficient ideal is \((a^2,ab,b^2)\).  Thus the full fibre is
\begin{equation}\label{eq:fat-fibre-v151}
                     \Spec\C[a,b]/(a,b)^2,
\end{equation}
of length three and tangent dimension two.  It is not a reduced
point, despite uniqueness of the geometric divisor.  The family
\[
 K_t=x\langle xy+t y^2,\ xz+t z^2\rangle
\]
has independent generators for every \(t\), exact gcd \(x\) when
\(t\ne0\), and the displayed collision at \(t=0\).  Consequently
this nonreduced fibre occurs on the boundary of the exact stratum,
not in an unrelated example.
\end{example}

\begin{example}[Several branches and a normal gcd jump]
\label{ex:branches-normal-jump-v151}
For the same \(g,h,r,E\), let
\(K=xy\langle x,z\rangle\).  The divisors \(x\) and \(y\)
give two points of the normalization fibre.  At both points the
overlap is one, so the fibre consists of two reduced points.
The point of \(Y_1\) is nonnormal.  By contrast, put
\(q=x^2+yz\) and \(K=xq\langle y,z\rangle\), now with \(h=3\).
The polynomial \(q\) is irreducible and is not divisible by \(x\).
There is exactly one linear divisor of \(\gcd K=xq\), and its
overlap with the residual gcd \(q\) is one.  This point of \(Y_1\)
is normal although its gcd degree is three.  These two examples
show why neither gcd degree nor the number of geometric lifts alone
is a sufficient normality criterion.
\end{example}

\begin{corollary}[An intrinsic boundary for first-relation algebras]
\label{cor:algebra-boundary-v151}
On the maximal-lower-Hilbert-function stratum of
Proposition~\ref{prop:relative-first-relation}, the algebras whose
intrinsic degree-\(D\) kernel has classifying map in \(Y_g\)
have, after removing only tangent-identity substitutions, stack
\([Y_g/\operatorname{GL}(E)]\).  Its normalization is the finite
representable morphism
\begin{equation}\label{eq:quotient-normalization-v151}
 [\mathscr X_g/\operatorname{GL}(E)]
                    \longrightarrow[Y_g/\operatorname{GL}(E)].
\end{equation}
All members of the universal relation family give finite locally
free algebras of rank \(\binom{e+D}{D}-r\).  This construction imposes
no pencil image equations and permits the divisor and residual
common factor to collide.
\end{corollary}
\begin{proof}
For every subbundle \(\mathcal K\) classified by the Grassmannian,
\(\Sym E/((\mathcal K)+\mathfrak m^{D+1})\) is, as a module,
the sum of \(S_j\) for \(j<D\) and \(S_D/\mathcal K\).
This proves finite local freeness and the rank.  The first-relation
proposition identifies the quotient after the specified nonlinear
inertia is removed.  Incidence and multiplication are
\(\operatorname{GL}(E)\)-equivariant.  Pulling
\eqref{eq:quotient-normalization-v151} back along the smooth frame
atlas \(Y_g\to[Y_g/\operatorname{GL}(E)]\) gives precisely
\(\nu_g\).  Finiteness, representability, and the normalization
identification follow by descent; equivalently normality and integral
closure may be checked on these smooth atlases.  This quotient is
not claimed to classify smoothings outside the specified Hilbert
stratum.
\end{proof}

For the pencil first relations one can take
\(e=n^2\), \(g=n(n-1)\), \(h=3n-4\), and
\(r=\binom{\binom{n+1}{2}}2\).  Every pencil then lies in the
exact-gcd open of this canonical boundary.  The boundary theorem
concerns the much larger space of all degree-\(g\) divisorial first
relations, rather than a closure defined by recovering a pencil and
reimposing its Pluecker equations.  It determines its normalization
and normal locus, but does not assert a classification of all its
linear orbits or of all finite-flat algebra deformations.

% BEGIN NEW MATHEMATICS v153
\section{Formal images and marked divisor germs}
\label{sec:formal-images-v153}
All rings and schemes in the following sections are over $\C$. Write
$S_j=\Sym^j E$, and let $Y_g\subset\Gr(r,S_D)$ be the
scheme-theoretic image of
\[
 \nu_g:\Pj(S_g)\times\Gr(r,S_{D-g})\longrightarrow\Gr(r,S_D),
 \qquad (f,J)\longmapsto fJ,
\]
where $\dim E\geq2$, $g\geq1$, and $2\leq r\leq\dim S_{D-g}$.
The source is smooth, and $\nu_g$ is its finite normalization: a fibre
records the degree-$g$ divisors of the common divisor, and on the
exact-degree-$g$ locus there is one reduced lift. We shall retain the
scheme structures of these fibres and of the image.

\begin{lemma}[Completed finite images]\label{lem:formal-image-v153}
Let $X\to M$ be a finite morphism of complex finite-type schemes, with
$X$ reduced and $M$ smooth, and let $Y$ be its scheme-theoretic image.
Fix a closed point $y\in Y$, put $R=\widehat{\OO}_{M,y}$, and write
$B_i=\widehat{\OO}_{X,x_i}$ for the points over $y$. Then
\[
 \widehat{\OO}_{Y,y}=R/\bigcap_i\ker(R\to B_i).
\]
This completed image is reduced. If the cotangent map from $R$ to
$B_i$ is surjective, then $R\to B_i$ is surjective. If there are two
such points, writing $I_i=\ker(R\to B_i)$ gives
\[
 \widehat{\OO}_{Y,y}=(R/I_1)\mathbin{\times}_{R/(I_1+I_2)}(R/I_2).
\]
The formal schemes at the different $x_i$ are disjoint. A local Artin
point of the finite inverse image belongs to exactly one of them.
\end{lemma}
\begin{proof}
On an affine neighbourhood the image ring is the image of the
ambient ring in a finite algebra. Exactness of completion on finite
modules preserves this kernel. The completed finite algebra is the
product of the $B_i$: the orthogonal idempotents of its Artin closed
fibre lift uniquely through the powers of the maximal ideal. Reduced
finite-type complex rings are excellent and analytically unramified,
so their completions are reduced; alternatively this applies directly
to the reduced image ring. For cotangent surjectivity, lift a basis
of the maximal ideal of $B_i$ modulo its square to $R$. The resulting
lifts generate every successive associated-graded piece. Subtracting
successive homogeneous lifts and taking their convergent sum expresses
any element of $B_i$ as the image of an element of $R$. The finite map
factors through the completed image, so this is also a closed immersion
into that image, not only into $M$. Finally the elementary exact sequence
\[
0\longrightarrow R/(I_1\cap I_2)\longrightarrow R/I_1\oplus R/I_2
 \longrightarrow R/(I_1+I_2)\longrightarrow0
\]
proves the fibre-product assertion. The idempotent decomposition gives
the last assertion, including over nonreduced Artin bases.
\end{proof}

\begin{lemma}[Coprime lifts and the split codimension]
\label{lem:coprime-lifts-v153}
If homogeneous forms $a,b$ have coprime reductions over a local Artin
base $A$, then $(a)\cap(b)=(ab)$ in $A[E]$. The relevant graded
quotients and Koszul kernels are $A$-flat. In the split-divisor model,
with $s_j=\binom{e+j-1}{j}$, $e\geq2$, $h>g\geq1$, and $r\geq2$,
\[
 c=r(s_h-s_{h-g})-(s_g-1)\geq(r-1)(s_g-1)\geq1.
\]
\end{lemma}
\begin{proof}
Over the residue field the two forms form a regular sequence. In
each degree their Koszul complex consists of finite free modules and
has its prescribed ranks. The acyclic-complex criterion over an Artin
local ring, proved by lifting splittings successively along its maximal
ideal, gives exactness and flatness of the lifted cokernels. Thus $b$
is a nonzerodivisor modulo $a$, proving the intersection identity.
For the inequality, express $s_h-s_{h-g}$ as the sum of $g$ consecutive
first differences of the sequence $s_j$. These differences are
nondecreasing for $e\geq2$. Their sum is at least $s_g-s_0=s_g-1$.
\end{proof}

In particular the two formal marked lifts in a coprime split model
are separated before their images are compared. Their intersection
functor is precisely $K_A=a_A b_A L_A$. Division gives a subbundle:
multiplication by a fibrewise nonzero form is locally split, and the
quotient of the ambient bundle by $K_A$ is locally free. Combining
this Artin-functor identification with Lemma~\ref{lem:formal-image-v153}
justifies the completed-image fibre product without any assumption
that a tangent-space comparison determines an intersection scheme.

\section{A global factorization theorem for binary normalization fibres}
\label{sec:binary-fibres-v153}
For this section $E$ has dimension two. Set $N_0=r(D+1-r)$ and
$m=r-1$. Let
\[
 \mathcal T_s=Y_s\setminus Y_{s+1},\qquad 1\leq s\leq D+1-r,
\]
with the reduced locally closed structure; an inadmissible $Y_j$ is
empty. Binary forms split into linear factors, so the $Y_j$ are
nested. The multiplication map identifies
\[
 \mathcal T_s\simeq\Pj(S_s)\times X_{D-s,r},
\]
where $X_{D-s,r}$ is the open Grassmannian of gcd-free residual spaces.
The first projection is the algebraic common-divisor morphism, not a
choice made independently on individual fibres.

For nonnegative $a,b$ define $H_{a,b}=\C$ when $ab=0$. Otherwise,
put $k=\min(a,b)$ and $l=\max(a,b)$, and set
\begin{equation}\label{eq:H-ab-v153}
 H_{a,b}=\C[q_1,\ldots,q_k]/(c_{l+1},\ldots,c_{l+k}),\quad
 \sum_{j\geq0}c_jz^j=(1+q_1z+\cdots+q_kz^k)^{-1}.
\end{equation}
The weight of $q_i$ is $i$. This grading is not the maximal-ideal
adic grading when $k>1$.

\begin{theorem}[Stratified finite flatness and all binary fibres]
\label{thm:binary-fibres-v153}
For $g\leq s$, the restriction of $\nu_g$ over $\mathcal T_s$ is the
base change, along the common-divisor morphism, of
\[
 \Pj(S_g)\times\Pj(S_{s-g})\longrightarrow\Pj(S_s),\qquad (f,H)\mapsto fH.
\]
It is finite locally free of degree $\binom{s}{g}$. If
$\gcd(K)=\prod_{i=1}^{b}\ell_i^{s_i}$, with distinct $\ell_i$ and
$\sum_i s_i=s$, its entire normalization fibre is
\begin{equation}\label{eq:binary-fibre-product-v153}
 \prod_{\substack{0\leq a_i\leq s_i\\\sum_i a_i=g}}
                   \bigotimes_{i=1}^{b} H_{a_i,s_i-a_i}.
\end{equation}
Every local factor is an Artin complete intersection and is Gorenstein.
Its length, weighted Hilbert series, and socle weight are, respectively,
\begin{gather*}
 \prod_i\binom{s_i}{a_i},\qquad
 \prod_i\prod_{j=1}^{\min(a_i,s_i-a_i)}
       \frac{1-z^{\max(a_i,s_i-a_i)+j}}{1-z^j},\qquad
 \sum_i a_i(s_i-a_i).
\end{gather*}
The total length is $\binom{s}{g}$, including at arbitrary simultaneous
root collisions. Over fixed $s$, the closure order of common-root
partition strata is the order generated by merging parts.
\end{theorem}
\begin{proof}
Over $\mathcal T_s$ the common divisor $G$ and residual subbundle $L$
are defined scheme-theoretically by the inverse to the exact-divisor
incidence. Locally on any base change, choose a member of $L$ coprime
to $G$ on each relevant geometric fibre. There is such a member over
$\C$, since only finitely many roots must be avoided. Regular-sequence
cancellation, or monic division and its flat quotient, gives
$f\mid GL$ if and only if $f\mid G$. This equivalence is an equality
of incidence functors over nonreduced bases. The residual subspace is
then forced. This proves the base-change description.

The displayed multiplication map is finite by properness and the
finiteness of the divisors of a fixed binary form. Its source and
target are smooth of dimension $s$. A finite map between these smooth
spaces is flat: a regular system of parameters in the target is a
system of parameters in each source local ring, hence a regular
sequence; the local flatness criterion applies. Its degree can be
computed at a squarefree form, where the $\binom{s}{g}$ divisors are
reduced and distinct.

Coprime root clusters split uniquely over Artin bases by Hensel
factorization. On the cluster with polynomial $x^{s_i}$, choosing a
degree-$a_i$ factor is the equation $f_i H_i=x^{s_i}$ with both factors
monic. Eliminating one factor by reciprocal expansion gives
\eqref{eq:H-ab-v153}. This proves the full product decomposition,
not only its closed points. Its equations have their sole common zero
at the origin: a monic factor of $x^{a+b}$ over $\C$ is a power of
$x$. Thus its $k$ equations in $k$ variables are a homogeneous system
of parameters and a regular sequence. The Hilbert series is the
quotient of the corresponding weighted polynomial Hilbert series.
Taking its value at one gives $\binom{a+b}{a}$; its top weight is
$ab$. A zero-dimensional complete intersection has a one-dimensional
socle. Equivalently $q_k^l$, up to a nonzero scalar, represents the
rectangular Schur class. The identification with the usual
Grassmannian cohomology presentation is classical; a purely algebraic
Schur-basis statement is \cite[Theorem~2.7]{Grinberg2019}.
Tensoring the local factors proves all three formulas. The total
length identity is the coefficient of $z^g$ in
$\prod_i(1+z)^{s_i}$. Finally the exact-gcd stratum is the product
with $X_{D-s,r}$ above. In the symmetric power of the projective line,
root clusters specialize by coalescence, and every merging is realized
by moving the corresponding roots together. This gives precisely the
stated closure order within fixed gcd degree.
\end{proof}

\begin{remark}
The finite flat statement is on the exact-gcd stratum with its indicated
scheme structure. It does not assert that normalization or formation of
scheme-theoretic images commutes with an arbitrary nonflat base change
of the ambient Grassmannian. It identifies a specific base change of
the already constructed incidence morphism.
\end{remark}

\section{Conductors and all branch intersections on the squarefree locus}
\label{sec:squarefree-v153}
Let $\mathcal U\subset\Gr(r,S_D)$ be the open set on which the common
divisor is squarefree. Its complement is the proper image of the
incidence $K=\ell^2L$. This condition concerns the common divisor,
not every member of the subspace.

\begin{theorem}[Squarefree conductor stratification]
\label{thm:squarefree-v153}
Let $K\in Y_g\cap\mathcal U$, let $s=\deg\gcd(K)$, and set
\[
 t=s-g,\qquad m=r-1,\qquad q=N_0-ms.
\]
Introduce $s$ disjoint blocks $u_i=(u_{i1},\ldots,u_{im})$, and let
$J_j$ be the ideal generated by products of one variable from each of
$j$ distinct blocks. Put $J_0=(1)$ and $J_{s+1}=0$. Then
\begin{equation}\label{eq:squarefree-ring-v153}
 \widehat{\OO}_{Y_g,K}\simeq
 \C[[z_1,\ldots,z_q,u_{11},\ldots,u_{sm}]]/J_{t+1}.
\end{equation}
The branches are indexed by the subsets $A\subset\{1,\ldots,s\}$
of cardinality $g$. The branch $A$ has ideal
$I_A=(u_{ij}:i\in A)$ in the ambient power-series ring. For any
collection of branches, their full scheme-theoretic intersection has
ideal $I_{\cup A}$; it is smooth of dimension
$q+m(s-|\cup A|)$.

The conductor in \eqref{eq:squarefree-ring-v153} is $J_t/J_{t+1}$.
Its restriction to branch $A$ is
$\prod_{i\notin A}(u_{i1},\ldots,u_{im})$. Consequently the global
conductor sheaf on $Y_g\cap\mathcal U$ is exactly
\begin{equation}\label{eq:global-conductor-v153}
 \mathfrak c_{Y_g}|_{\mathcal U}
        =\II_{Y_{g+1}/Y_g}|_{\mathcal U}.
\end{equation}
In particular it is radical there. These local rings are seminormal.
Their dimension, depth and multiplicity are
\[
 q+mt,\qquad q+t,\qquad \binom{s}{g}.
\]
They are Cohen--Macaulay exactly when $s=g$ or $r=2$.
The Hilbert series of the nonsmooth factor in its standard grading is
\begin{equation}\label{eq:block-hilbert-v153}
 \sum_{j=0}^{t}\binom{s}{j}\bigl((1-z)^{-m}-1\bigr)^j.
\end{equation}
The whole collection of branch intersections and conductor strata
is invariant under permutation of the roots and descends from any
ordered-root etale chart.
\end{theorem}
\begin{proof}
Choose $F_0\in K$ with a simple zero at each of the $s$ common roots,
and complete it to a frame $F_0,\ldots,F_m$. Hensel factorization gives
moving roots $x_1,\ldots,x_s$ in the completion. Use as coordinates
the $s$ root parameters and the $ms$ evaluations
$u_{ij}=F_j(x_i)$, with harmless invertible normalizing factors.
These are part of a regular parameter system: the evaluation map
$S_D/K\to\C^s$ is surjective because $K$ vanishes on the common-root
scheme, and $D\geq s-1$. Thus
$\Hom(K,S_D/K)$ allows independent variation of all $rs$ jets.
The passage to the displayed coordinates is triangular with invertible
diagonal, since $F_0$ has simple zeros. The other $N_0-rs$ parameters
and the $s$ root parameters give exactly $q=N_0-ms$ smooth parameters.

A marked degree-$g$ divisor near this fibre chooses a subset $A$ of
these roots. Its incidence equations are exactly $u_{ij}=0$ for
$i\in A$, with no nilpotents suppressed. There are no other divisor
lifts. Lemma~\ref{lem:formal-image-v153} therefore identifies the
completed image ideal with $\bigcap_{|A|=g}I_A$. A monomial belongs
to this intersection exactly when its block support meets every
$g$-element subset, or equivalently has at least $s-g+1=t+1$ blocks.
This proves \eqref{eq:squarefree-ring-v153}. Sums of the coordinate
ideals give all intersections, including intersections of three or
more branches.

For completeness compute the conductor inside the product of the
branch rings. Its component on branch $A$ consists of functions whose
extension by zero on every other branch belongs to the image ring.
The image of $\bigcap_{B\ne A}I_B$ modulo $I_A$ is precisely
$\prod_{i\notin A}(u_i)$. Indeed a surviving monomial has support
outside $A$, and must meet every $g$-subset other than $A$. Replacing
one element of $A$ by each $i\notin A$ forces all those blocks to
occur. Conversely their product meets every such subset. Hence a
conductor monomial has exactly $t$ surviving blocks; monomials on more
blocks are zero. This is $J_t/J_{t+1}$. The same coordinates identify
$Y_{g+1}$ with $J_t$, proving the sheaf equality by faithful flatness
of completion. The case $t=0$ means that the conductor is the unit
ideal and the next stratum is absent locally.

The ring is the ring of tuples on its coordinate branches that agree
on every pairwise intersection. To verify this directly, compare the
coefficient of each monomial on all branches where it survives; the
pairwise agreements give a unique common coefficient. If an element
$b$ of the total fraction ring has $b^2,b^3$ in the image, it first
belongs to the normalization, by integrality. Its two restrictions
on any intersection have equal squares and cubes, hence are equal
in that reduced ring. Thus $b$ lies in the image. This proves
seminormality by the square-and-cube criterion.

Formula~\eqref{eq:block-hilbert-v153} follows by sorting nonzero
monomials by their block support. The top-dimensional contributions
come from the $\binom{s}{t}$ supports of size $t$, giving dimension
$mt$ and multiplicity $\binom{s}{t}$. We supply the homological step
for the depth, since branch counting alone does not give it. The
simplicial complex of the squarefree monomial ideal $J_{t+1}$ consists
of sets using at most $t$ blocks. Any induced subcomplex meeting $b$
blocks is homotopy equivalent to the $(t-1)$-skeleton of the simplex
on those $b$ blocks: collapse each block to a representative vertex.
The two composites are contiguous to the identity, because adjoining
the representatives never increases the number of blocks. If $b\leq t$
the complex is a simplex; otherwise its only reduced homology is in
degree $t-1$, with rank $\binom{b-1}{t}$. This last assertion follows
also by the elementary induction that adds the last vertex to a
simplex skeleton.

The squarefree multidegree-$W$ part of the Koszul complex on all $sm$
variables is the augmented simplicial cochain complex of this induced
subcomplex, with degree shifted by $|W|-1$. Non-squarefree multidegrees
are contractible in the relevant summands, by pairing faces using a
variable whose exponent exceeds one. Consequently a nonzero positive
Betti number has homological degree $|W|-t$. The full vertex set gives
the nonzero degree $sm-t$, since $s>t$ when $g>0$. The projective
dimension is therefore $sm-t$, and the Auslander--Buchsbaum formula
gives depth $t$. When $t=0$ the quotient is simply $\C$. Adjoining
$q$ power-series parameters gives depth $q+t$ in every case, and
comparison with $q+mt$ proves the Cohen--Macaulay criterion.
All the constructions commute with relabelling roots, proving descent.
\end{proof}

\begin{corollary}[Binary divisor-image intersections]
\label{cor:binary-intersections-v153}
For admissible binary divisor degrees $a_1,\ldots,a_v$, their
scheme-theoretic intersection is $Y_{\max a_i}$. On the squarefree
locus all multiple intersections of normalization branches, rather
than only the degree images, are given by the union-of-subsets rule
in Theorem~\ref{thm:squarefree-v153}.
\end{corollary}
\begin{proof}
The reduced binary images are nested as closed subschemes: a binary
form of degree $b$ has a divisor of each degree $a\leq b$, and
reduced closed subschemes are determined by their closed points.
Their ideals are thus nested, so their sum is the largest ideal.
The assertion about branches has already been proved as an identity
of coordinate ideals, without taking radicals.
\end{proof}

\section{An incidence atlas for all binary root collisions}
\label{sec:collision-atlas-v153}
The squarefree theorem has a uniform replacement when any number of
root clusters collide. This replacement specifies the full completed
image, not merely a transverse family in it.

\begin{theorem}[Simultaneous collision atlas]\label{thm:atlas-v153}
Let $K\in Y_g$ be binary with exact common divisor
$\prod_{i=1}^b\ell_i^{s_i}$ and $s=\sum_i s_i$. There are formal
ambient coordinates consisting of the coefficients of monic
polynomials $P_i$ of degree $s_i$, polynomials $R_{ij}$ of degree
less than $s_i$ ($1\leq j\leq m$), and $N_0-rs$ smooth parameters
$z$. They are centred at $P_i=x_i^{s_i}$ and $R_{ij}=0$.
For every $\alpha=(a_i)$ with $0\leq a_i\leq s_i$ and $\sum a_i=g$,
let $B_\alpha$ be the regular power-series ring in the coefficients
of monic $A_i,B_i'$ of degrees $a_i,s_i-a_i$, polynomials $U_{ij}$
of degree less than $s_i-a_i$, and the same $z$. The completed
normalization and image are exactly
\begin{equation}\label{eq:collision-image-v153}
 T\longrightarrow\prod_\alpha B_\alpha,\qquad
 P_i\mapsto A_iB_i',\quad R_{ij}\mapsto A_iU_{ij},\qquad
 \widehat{\OO}_{Y_g,K}=T/\bigcap_\alpha\ker(T\to B_\alpha).
\end{equation}
Here $T=\C[[\operatorname{coeff}(P_i),\operatorname{coeff}(R_{ij}),z]]$.
Degree-zero monic factors are one, and negative-degree remainder
spaces are zero. The formula is independent, up to isomorphism of
completed incidence diagrams, of the frame, the selected $F_0$, and
the local root coordinates. It covers every higher and simultaneous
binary collision.
\end{theorem}
\begin{proof}
Choose $F_0\in K$ having exact order $s_i$ at every common root.
The required nonvanishing of its residual value excludes finitely
many hyperplanes in $K$, so a choice exists. Complete it to a frame.
Weierstrass preparation at the distinct roots produces $P_i$; division
of the remaining frame members produces $R_{ij}$. The map
\[
 S_D/K\longrightarrow\bigoplus_i\C[x_i]/(x_i^{s_i})
\]
is surjective: its kernel contains $K$, and the usual interpolation
map from $S_D$ is surjective. Thus the Grassmannian tangent space
permits independent variation of all $rs$ jets. Division and preparation
change these jets by an invertible triangular matrix, whose diagonal
contains the nonzero residual values of $F_0$. Their differentials
are independent. Formal inverse coordinates supply the remaining
$N_0-rs$ parameters; this integer is nonnegative since
$r\leq D-s+1$.

A marked divisor near a chosen normalization point splits uniquely
among the coprime root clusters. In cluster $i$ its monic representative
has degree $a_i$ and divides both $P_i$ and all $R_{ij}$. Monic division
is equivalent, over every local Artin base, to the unique equations
$P_i=A_iB_i'$ and $R_{ij}=A_iU_{ij}$. Conversely these equations give
the marked divisor and the residual subbundle. Thus the normalized
incidence functor is represented by the displayed regular ring.
It is finite over $T$: monic factor coefficients are integral over the
coefficients of $P_i$, while the remaining factors and quotients are
uniquely determined by division. Lemma~\ref{lem:formal-image-v153}
now gives the exact completed image and the intersection of kernels.

Changing a frame or a coordinate changes only a presentation of the
same functor of a subbundle with a marked divisor. The mutually inverse
changes of incidence data therefore induce isomorphisms of its complete
representing rings and of the resulting finite-image diagrams. This
also proves that the calculation describes the entire completed image.
A point of that image over an Artin ring need not lift to its
normalization; no such lifting assertion is used or implied.
\end{proof}

\section{The singular Fitting ideal of every double-root model}
\label{sec:fitting-v153}
The atlas specializes for $s=2$, $g=1$ to the following ring, after
completion of the square and suppression of smooth parameters:
\[
 A_m=R\oplus It\subset B=\C[[u_1,\ldots,u_m,t]],\qquad
 R=\C[[u_1,\ldots,u_m,\Delta]],\ I=(u_1,\ldots,u_m),\ \Delta=t^2.
\]
Put $w_i=u_it$. The defining equations are
$u_iw_j-u_jw_i$ and $w_iw_j-\Delta u_iu_j$.

\begin{theorem}[Closed singular-ideal formula]\label{thm:fitting-v153}
For every $m\geq1$, the intrinsic singular Fitting ideal is
\begin{equation}\label{eq:fitting-v153}
 \Fitt_{m+1}\Omega^{\mathrm{cont}}_{A_m/\C}
  =(I^{m+1}+\Delta I^m)\oplus I^m t
  = I^{m-1}(I^2,(w_1,\ldots,w_m),\Delta I).
\end{equation}
The first equality is a decomposition as an $R$-submodule of $B$;
the right side is an ideal of $A_m$. It is not merely an ideal of
Jacobian minors left unevaluated. For $m=1$ it is
$(u^2,w,\Delta u)$, and its radical is the conductor $(u,w)$ for
all $m$.
\end{theorem}
\begin{proof}
There are $2m+1$ ambient variables and dimension $m+1$, so the Fitting
ideal consists of the $m$-minors of the Jacobian of the displayed
relations. Evaluate this matrix in $B$ using $w_i=u_it$ and
$\Delta=t^2$. A minor not using the $\Delta$ column has $u$-degree
$m$. It vanishes at $t=0$: the only remaining rows are the Koszul
rows in the $w$ columns, of rank at most $m-1$. Hence its terms have
a positive power of $t$, and belong to $\Delta I^m\oplus I^m t$.
A minor using the $\Delta$ column has $u$-degree $m+1$, and hence
lies in $I^{m+1}\oplus I^{m+1}t$. This proves the first inclusion.

Use the $m$ rows
\[
 u_1w_j-u_jw_1\ (2\leq j\leq m),\qquad w_1^2-\Delta u_1^2.
\]
The columns $(w_1,w_2,\ldots,w_m)$, $(u_1,w_2,\ldots,w_m)$,
and $(\Delta,w_2,\ldots,w_m)$ give, up to sign and nonzero scalar,
$u_1^m t$, $\Delta u_1^m$, and $u_1^{m+1}$. The ideal is invariant
under simultaneous linear change of $u$ and $w$. In characteristic
zero the powers of linear forms span each symmetric power, so these
three minors generate the three indicated spaces of all $u$-monomials.
They generate the claimed $R$-module; it is an $A_m$-ideal since
multiplication by $w_i=u_it$ preserves it. This proves equality.
The formula also immediately gives its radical, without discarding
the nonreduced Fitting structure in the assertion itself.
\end{proof}

\section{A structural dichotomy for pure-power divisor fibres}
\label{sec:reciprocal-v153}
Let $E=\C a\oplus W$, $d=\dim W\geq1$, and $g,k\geq1$.
The factorization fibre of $a^{g+k}$ records monic factors of degrees
$g$ and $k$. Interchanging these factors identifies its two descriptions;
we may and shall take $g\geq k$. Write
\[
 Q_i\in\Sym^iW\quad(1\leq i\leq k),\qquad
 C_0=1,\quad C_j=-\sum_{i=1}^kQ_i C_{j-i},\quad C_j=0\ (j<0).
\]
Let $P=\Sym(\bigoplus_{i=1}^k(\Sym^iW)^*)$, with its generators
of weight $i$, and let $\mathfrak n$ be its variable ideal. Put
\begin{equation}\label{eq:reciprocal-algebra-v153}
 F_{g,k}(W)=P/\mathcal I,\qquad
 \mathcal I=(\operatorname{coeff}C_{g+1},\ldots,
                                    \operatorname{coeff}C_{g+k}).
\end{equation}
It is an Artin local algebra. Indeed its complex points are exactly
the factorizations of $a^{g+k}$, hence the single origin by unique
factorization. The polynomial and completed presentations coincide.
If $K=a^{g+k}L$ with $\gcd(L)=1$, cancellation by a member of $L$
not divisible by $a$, as in Lemma~\ref{lem:coprime-lifts-v153},
identifies this with the full degree-$g$ normalization fibre at $K$.

\begin{theorem}[Minimal equations and complete-intersection classification]
\label{thm:reciprocal-structure-v153}
The embedding dimension and minimal number of equations of
$F_{g,k}(W)$ in a minimal regular local presentation are
\begin{align}
 \operatorname{edim}F_{g,k}(W)&=\binom{d+k}{k}-1,\label{eq:edim-v153}\\
 \mu(\mathcal I)&=\sum_{j=g+1}^{g+k}\binom{d+j-1}{j}
                  =\binom{d+g+k}{d}-\binom{d+g}{d}.
                  \label{eq:minrels-v153}
\end{align}
More precisely the minimal relation module has the equivariant
weighted decomposition
\begin{equation}\label{eq:relation-module-v153}
 \mathcal I/\mathfrak n\mathcal I
                 \simeq\bigoplus_{j=g+1}^{g+k}(\Sym^jW)^*.
\end{equation}
Consequently this fibre is a complete intersection if and only if
$d=1$. For $d=1$, it has length $\binom{g+k}{k}$, weighted Hilbert
series $\prod_{i=1}^k(1-z^{g+i})/(1-z^i)$, and one-dimensional
socle generated by $Q_k^g$, of weight $gk$.
For $k=1$ and arbitrary $d$,
\begin{equation}\label{eq:k-one-v153}
 F_{g,1}(W)=\Sym(W^*)/\mathfrak n^{g+1},\quad
 \length F_{g,1}(W)=\binom{d+g}{g},\quad
 \operatorname{Soc}F_{g,1}(W)=\Sym^g(W^*).
\end{equation}
In particular the last algebra is Gorenstein precisely when $d=1$.
No Gorenstein classification for $d>1,k>1$ is needed for the theorem.
\end{theorem}
\begin{proof}
Since $g\geq k$, each defining form has no linear term in the
coefficient variables. Thus the presentation is minimal and
\eqref{eq:edim-v153} follows by summing $\dim\Sym^iW$.
For $d=1$, the regular-sequence argument in the proof of
Theorem~\ref{thm:binary-fibres-v153} applies to the $k$ reciprocal
relations. In particular none of them belongs to $\mathfrak n\mathcal I$.
It gives the stated Hilbert series, length, and the classical rectangular
socle class.

For general $W$, taking coefficients of $C_j$ gives an equivariant
map $(\Sym^jW)^*\to\mathcal I/\mathfrak n\mathcal I$.
Its source is irreducible. This map is nonzero: choose one coordinate
$x$ of $W$, keep only the coefficients of $x^i$ in every $Q_i$, and
set every other coefficient to zero. The coefficient of $x^j$ maps
to the $j$th binary reciprocal relation. Were it in
$\mathfrak n\mathcal I$ before specialization, its specialization
would be in the binary $\mathfrak n\mathcal I$, contradicting the
preceding minimality. Irreducibility now makes the entire coefficient
map injective. The different $j$ have different weights, so their
images are independent. They generate the relation module by the
chosen presentation, proving \eqref{eq:relation-module-v153} and
\eqref{eq:minrels-v153}.

If $d\geq2$, the sequence $\binom{d+j-1}{j}$ strictly increases in
$j$. Comparing its terms at $j=g+i$ and $j=i$ shows that the minimal
number of relations is strictly greater than the embedding dimension.
A zero-dimensional complete intersection in a minimal regular local
presentation has exactly that many minimal relations as variables;
hence these fibres are not complete intersections. This proves both
directions of the classification. Finally when $k=1$ the sole form
is $C_{g+1}=(-Q_1)^{g+1}$. Its coefficients are, with nonzero multinomial
factors, every monomial of degree $g+1$ in the $d$ variables. This gives
\eqref{eq:k-one-v153}, including the socle and length.
\end{proof}

\begin{corollary}[The ternary-pencil ramification fibre]
\label{cor:ternary-fibre-v153}
At the common limiting relation algebra with $e=9$, $g=6$ and $k=2$,
the normalization fibre has embedding dimension $44$ and exactly
\[
 \binom{14}{7}+\binom{15}{8}=9867
\]
minimal scalar equations. Their two weighted pieces have dimensions
$3432$ and $6435$, in weights seven and eight. In particular this
actual pencil-orbit boundary fibre is not a complete intersection.
\end{corollary}
\begin{proof}
Here $d=e-1=8$. Apply \eqref{eq:edim-v153} and
\eqref{eq:relation-module-v153}. These are minimal equation counts,
not the number of coefficient variables or the length of the fibre.
\end{proof}

The binary factorization rings in this section are classical. The
additional point of the multivariate argument is that polarization
preserves every full irreducible coefficient space in the minimal
relation module. Neither the Grassmannian identification nor a finite
calculation in a few variables is being used as a claim of priority
or as a substitute for that argument.

\section{Homological invariants of the divisor fibres}
\label{sec:homological-v154}
We first extract all graded Betti numbers from the squarefree local
model and then identify a determinant-apolar quotient of the
multivariate fibres. A separate exact calculation treats an entire
multivariate fibre, rather than only a quotient or its minimal
relations. The grading conventions will be specified in each case.

\begin{theorem}[Betti numbers of the squarefree block model]
\label{thm:block-betti-v154}
Let $S=\C[u_{ij}:1\leq i\leq s,1\leq j\leq m]$ with its standard
multigrading, let $0\leq t<s$, and put $A=S/J_{t+1}$, with the
block-support ideal of Theorem~\ref{thm:squarefree-v153}.
For a nonempty squarefree multidegree $F$ let $b(F)$ be the number
of blocks it meets. Besides $\beta_{0,0}=1$, the only nonzero
multigraded Betti numbers are
\begin{equation}\label{eq:block-multibetti-v154}
 \beta_{i,F}^S(A)=\binom{b(F)-1}{t}
       \quad\text{if }b(F)>t\text{ and }i=|F|-t.
\end{equation}
For $t=0$ this is the Koszul resolution of $\C$.
For all $t$, the singly graded numbers are
\begin{equation}\label{eq:block-betti-v154}
 \beta_{i,i+t}^S(A)=
 \sum_{b=t+1}^{\min(s,i+t)}\binom{s}{b}\binom{b-1}{t}
       [z^{i+t}]\big((1+z)^m-1\big)^b\quad(i\geq1).
\end{equation}
The ideal $J_{t+1}$ has a $(t+1)$-linear resolution. The projective
dimension of $A$ is $sm-t$. In the Cohen--Macaulay case $m=1$,
its type is $\binom{s-1}{t}$; it is Gorenstein exactly when
$t=0$ or $t=s-1$. The same Betti numbers over the completed ambient
regular ring apply to the local model after adjoining its smooth
parameters.
\end{theorem}
\begin{proof}
The Stanley--Reisner complex consists of sets meeting at most $t$
blocks. If $t\geq1$, its induced subcomplex on $F$ has, in each
nonempty block, a simplex all of whose vertices have identical
external links. Contract that simplex to one of its vertices: the
map and the inclusion are contiguous, because adjoining the chosen
vertex to any face does not change its block support. Applying this
one block at a time gives a homotopy equivalence with the
$(t-1)$-skeleton of the simplex on $b(F)$ vertices. If $b(F)\leq t$
this simplex is contractible. Otherwise its reduced homology is
concentrated in degree $t-1$, with dimension $\binom{b(F)-1}{t}$.
The latter assertion follows by attaching its top faces successively,
or from the augmented simplicial chain complex and the binomial
identity for its Euler characteristic. Hochster's multigraded
formula now gives \eqref{eq:block-multibetti-v154}; see
\cite[Chapter~5]{MillerSturmfels2005}. If $t=0$, every variable is
in the ideal and the Koszul formula gives the same result directly.

For a support of size $b$, the coefficient in
\eqref{eq:block-betti-v154} counts choices of $i+t$ vertices meeting
each selected block. Summing proves the formula. All ideal syzygies
therefore have the indicated linear degrees. Taking $F$ to be all
$sm$ vertices gives the nonzero last Betti number
$\binom{s-1}{t}$ in homological degree $sm-t$; no larger degree
occurs. When $m=1$, Auslander--Buchsbaum gives depth $t$, equal to
dimension, and the last Betti number is the Cohen--Macaulay type.
Completion and adjoining independent regular parameters are flat
operations preserving the minimal free resolution. These are
classical squarefree-monomial methods. The assertion here is their
full specialization to the normalization image model, not a claim
of priority for Hochster's formula or skeleton resolutions.
\end{proof}

For the remaining results use the weighted reciprocal algebra
$F_{g,k}(W)$ of \eqref{eq:reciprocal-algebra-v153}, where
$d=\dim W$ and $g\geq k\geq2$. Set all $Q_i$ except $Q_2$ equal
to zero and, in addition, impose every coefficient of $Q_2^2$.
Denote the resulting algebra by
\begin{equation}\label{eq:apolar-quotient-v154}
 B_d=\Sym((\Sym^2W)^*)/(\operatorname{coeff}Q_2^2).
\end{equation}
Its coefficient variables have reciprocal weight two. When referring
to its ordinary polynomial grading, we instead give them degree one.

\begin{theorem}[A determinant-apolar quotient of every higher fibre]
\label{thm:apolar-quotient-v154}
There is a $\operatorname{GL}(W)$-equivariant graded surjection
$F_{g,k}(W)\twoheadrightarrow B_d$ for every $g\geq k\geq2$.
After fixing the standard coefficient/differential pairing, $B_d$
is the apolar algebra of the determinant of a generic symmetric
$d\times d$ matrix. Its polynomial-degree Hilbert function is
\begin{equation}\label{eq:Narayana-v154}
 \dim(B_d)_p=\frac{1}{d+1}\binom{d+1}{p}\binom{d+1}{p+1}
                 \quad(0\leq p\leq d).
\end{equation}
It is Gorenstein, its reciprocal-weight socle is $2d$, and
\begin{equation}\label{eq:Catalan-v154}
 \length F_{g,k}(W)\ \geq\ \length B_d
       =\frac{1}{d+2}\binom{2d+2}{d+1}.
\end{equation}
For $k=2$ and $g=2$ or $3$, the quotient by the coefficients of
$Q_1$ alone is already $B_d$.
In particular the actual ternary-pencil fibre $F_{6,2}(\C^8)$ has
length at least $4862$, independently of its embedding dimension
$44$ and its $9867$ minimal equations.
\end{theorem}
\begin{proof}
Under the stated substitution, the reciprocal series becomes
$(1+Q_2z^2)^{-1}$. Its odd coefficients vanish and its even
coefficients are powers of $Q_2$. Since $g+1\geq3$, imposing
$Q_2^2=0$ annihilates all defining coefficients. This proves the
surjection, functorially for linear changes of $W$. When $k=2$ and
$g=2$ or $3$, the first even defining coefficient is $Q_2^2$ and
there are no further independent equations after $Q_1=0$.

Choose coordinates $x_1,\ldots,x_d$ on $W$ and write
$Q_2=\sum_i y_{ii}x_i^2+\sum_{i<j}y_{ij}x_ix_j$.
Let $X=(X_{ij})$ be symmetric, with independent entries for $i\leq j$,
and interpret $y_{ij}$ as $\partial/\partial X_{ij}$. Then
$Q_2(x)$ acts as directional differentiation in the rank-one
symmetric direction $xx^t$. The determinant of $X+zxx^t$ is linear
in $z$, so every coefficient of $Q_2(x)^2$ annihilates $\det X$.
These coefficients are, up to nonzero scalars, $y_{ii}^2$,
$y_{ii}y_{ij}$, $y_{ij}^2+2y_{ii}y_{jj}$,
$y_{ii}y_{jl}+y_{ij}y_{il}$ for distinct $i,j,l$, and
$y_{ij}y_{kl}+y_{ik}y_{jl}+y_{il}y_{jk}$ for four distinct indices.
They are exactly the quadratic generators in
\cite[Proposition~3.1 and Theorem~3.10]{ShafieiSymmetric}.
That theorem proves equality with the full apolar ideal, not merely
containment. Its Corollary~2.7 gives \eqref{eq:Narayana-v154} and
the Catalan length. Alternatively, that Hilbert function counts the
independent minors obtained by differentiating $\det X$.
Apolarity to one homogeneous form gives a perfect pairing between
complementary degrees and a one-dimensional socle in degree $d$.
Replacing polynomial degree by twice that degree proves the weight
statement. The length inequality follows from the surjection.
For $d=8$ the Catalan number is $4862$.

The apolar Hilbert function and its Gorenstein property are classical
results of the cited determinant theory. The new identification in
this setting is the functorial quotient of the entire family of
normalization fibres. A quotient's Gorenstein property is not being
attributed to the original fibre.
\end{proof}

\begin{proposition}[The complete first multivariate two-factor fibre]
\label{prop:F22-v154}
Let $A=F_{2,2}(\C^2)$ and write
$Q_1=ax+by$, $Q_2=cx^2+dxy+ey^2$.
Then
\begin{gather}\label{eq:F22-hilbert-v154}
 H_A^{\mathrm{wt}}(z)=1+2z+6z^2+6z^3+7z^4,\qquad
 \length A=22,\qquad\dim\operatorname{Soc}A=7,\\
 \label{eq:F22-local-hilbert-v154}
 H_{\gr_{\mathfrak m}A}(z)=1+5z+6z^2+5z^3+5z^4.
\end{gather}
The weight grading gives $a,b$ weight one and $c,d,e$ weight two.
In the minimal weighted free resolution over
$S=\C[a,b,c,d,e]$, the shift polynomials
$B_i(z)=\sum_j\beta_{i,j}^S(A)z^j$ are
\begin{align*}
 B_0&=1,& B_1&=4z^3+5z^4,\\
 B_2&=4z^5+12z^6+12z^7,&
 B_3&=21z^8+20z^9+z^{10},\\
 B_4&=21z^{10}+8z^{11},&B_5&=7z^{12}.
\end{align*}
In particular the total Betti numbers are $(1,9,28,42,29,7)$.
This fibre is not Gorenstein. The assertion concerns this whole
five-variable fibre, not the determinant-apolar quotient.
\end{proposition}
\begin{proof}
The nine initial equations are the coefficients of
$-Q_1^3+2Q_1Q_2$ and $Q_1^4-3Q_1^2Q_2+Q_2^2$.
For graded reverse lexicographic order $e>d>c>b>a$, a Groebner basis
is the following list (a scalar multiple does not change a member):
\begin{gather*}
 (c,d,e)^3,\quad b(c,d,e)^2,\quad a(c,d,e)^2,\\
 b^2e-e^2,\quad b^2d-de,\quad b^2c-a^2e,\quad b^3-2be,\\
 2abe-de,\quad 2a^2e+2abd-2ce-d^2,\quad 2abc-cd,\\
 3ab^2-2ae-2bd,\quad a^2d-cd,\quad a^2c-c^2,\\
 3a^2b-2ad-2bc,\quad a^3-2ac.
\end{gather*}
Here a monomial ideal in the list means all its minimal monomial
generators. There are $10+6+6+12=34$ polynomials. Polynomial
division gives both equality with the ideal of the nine equations
and zero remainders for all their basis S-pairs. This is a finite
exact arithmetic verification; the explicit input and all reduction
and rank tests are supplied in the accompanying calculation.
The standard monomials, grouped by weight, are
\begin{equation}\label{eq:F22-basis-v154}
\begin{array}{c|l}
0&1\\
1&a,b\\
2&a^2,ab,b^2,c,d,e\\
3&ac,bc,ad,bd,ae,be\\
4&c^2,cd,d^2,a^2e,ce,de,e^2.
\end{array}
\end{equation}
Reduction of products of these monomials gives
$\dim A/\mathfrak m=1$ and
$\dim(\mathfrak m^i)=(21,16,10,5,0)$ for $1\leq i\leq5$.
The common kernel of multiplication by $a,b,c,d,e$ is exactly the
last row of \eqref{eq:F22-basis-v154}. This proves both distinct
Hilbert series and the type.

For precision, the resolution calculation uses no numerical rank
thresholds. In weight $w$ and homological degree $i$, take basis
pairs $(J,m)$ with $J\subset\{a,b,c,d,e\}$ of size $i$, $m$ a
standard monomial, and $\operatorname{wt}(J)+\operatorname{wt}(m)=w$.
The differential is
\begin{equation}\label{eq:F22-koszul-v154}
 d(J,m)=\sum_{v\in J}(-1)^{\operatorname{pos}(v,J)-1}
                  (J\setminus\{v\},\operatorname{NF}(vm)).
\end{equation}
All entries are rational, explicitly determined by the displayed
34 polynomials. Exact row reduction and
$\beta_{i,w}=\dim C_{i,w}-\rank d_{i,w}-\rank d_{i+1,w}$ give the
six shift polynomials above. The calculation is fully specified by
\eqref{eq:F22-basis-v154}--\eqref{eq:F22-koszul-v154}; the archived
rank table makes each finite elimination independently repeatable.
As additional checks, their alternating sum is
$H_A^{\mathrm{wt}}(z)(1-z)^2(1-z^2)^3$ and the last Betti space is
the socle shifted by the total variable weight eight. These checks
do not replace the exact rank computations.
\end{proof}

\section{Descent and complete-local details}
\label{sec:formal-details-v154}
The following statements supply the local-to-global steps used in
Sections~\ref{sec:formal-images-v153}--\ref{sec:reciprocal-v153}.
They also specify which assertions concern a completed image, a
normalization fibre, or a family on an exact-divisor stratum.

\begin{lemma}[Complete-local surjectivity and reduced completion]
\label{lem:complete-local-v154}
In the cotangent-surjectivity assertion of
Lemma~\ref{lem:formal-image-v153}, the rings are Noetherian complete
local $\C$-algebras with residue field $\C$, and the homomorphism
is local, continuous, and the identity on that residue field. Under
these hypotheses cotangent surjectivity implies ring surjectivity.
A reduced local ring essentially of finite type over $\C$ has
reduced completion.
\end{lemma}
\begin{proof}
Choose preimages in the source maximal ideal of generators of the
target cotangent space. Their degree-$q$ products span every quotient
$\mathfrak m_B^q/\mathfrak m_B^{q+1}$. Given an element of $B$, first
lift its residue, then its error successively in degrees $q=1,2,\ldots$.
The correction in degree $q$ lies in $\mathfrak m_R^q$. Their sum
converges in $R$, and continuity and separatedness imply that its
image is the desired element of $B$. This is not a statement about
arbitrary noncomplete local homomorphisms.

Finite-type complex rings are excellent. Quasi-excellent local
rings are Nagata and have geometrically reduced formal fibres;
see \cite[Tags 07QV and 0BJ0]{StacksLocalV154}. For a reduced local
ring $A$, the injection into the product of its finitely many
minimal-prime quotients remains injective after flat completion.
Each quotient is a Nagata domain and its completion is reduced.
Thus $\widehat A$ is a subring of a product of reduced rings.
This justifies the reduced completed-image step of the earlier lemma.
\end{proof}

\begin{lemma}[Coprime cluster products over Artin bases]
\label{lem:hensel-products-v154}
Let $A$ be a local Artin $\C$-algebra. If the reduction of a monic
polynomial has a factorization into pairwise coprime monic factors,
there is a unique factorization into monic lifts of the same degrees.
In a factorization fibre, the functor of divisors with prescribed
degrees in the distinct residual clusters is the product of the
cluster divisor functors. Its coordinate algebra is therefore the
tensor product of their Artin coordinate algebras.
\end{lemma}
\begin{proof}
For two coprime factors $p,q$, the linear correction map
\[
 (\dot p,\dot q)\longmapsto q\dot p+p\dot q,
 \quad\deg\dot p<\deg p,\quad\deg\dot q<\deg q,
\]
is an isomorphism on coefficient spaces over the residue field.
Its kernel is zero by coprimeness, and source and target have equal
dimension. It is therefore an isomorphism over $A$ as well. Across
a square-zero extension of Artin rings it uniquely corrects a
lifted product to the given polynomial. Induction on the nilpotence
order and then on the number of clusters proves existence and
uniqueness. Apply the same argument to both the marked divisor and
its complement. The cluster product and decomposition constructions
are mutually inverse on every Artin base; Yoneda gives the asserted
product of finite schemes. Passing to inverse limits gives the
corresponding assertion on complete local bases.
\end{proof}

\begin{lemma}[The local choice and descent in the binary stratum]
\label{lem:residual-selection-v154}
The base-change identification of
Theorem~\ref{thm:binary-fibres-v153} holds as an identification of
incidence functors, not only on closed fibres. The required choice
of a member of the residual space coprime to the common divisor can
be made on a Zariski open cover of the exact-gcd stratum; the choices
are auxiliary and the resulting identifications glue uniquely.
\end{lemma}
\begin{proof}
Write a point of the stratum as $(G,L)$ in
$\Pj(S_s)\times X_{D-s,r}$. The vanishing of a member of $L$ at
one of the finitely many roots of $G$ is a proper linear condition
on that member. Over $\C$ a finite union of these proper hyperplanes
does not cover $L$. Choose a linear combination avoiding them, lift
its coefficients on a bundle-trivializing neighbourhood, and invert
its resultant with $G$. This gives the required Zariski neighbourhood.
These neighbourhoods cover the stratum and may be pulled back by
any base change.

On the marked-divisor incidence, pass locally to a monic coordinate
for the proposed divisor $f$. The quotient by $f$ is finite free.
On every geometric fibre $f$ divides $G$, so the chosen residual
member is invertible in that quotient, by coprimeness. Its
multiplication matrix has invertible determinant locally, hence it
is invertible also over a nonreduced base. Thus $f\mid G L$ implies
$f\mid G$ as an equation on that base, and the converse is immediate.
Monic division forces the remaining subspace. The identification is
independent of which coprime residual member established cancellation,
so it glues on overlaps. No lifting of arbitrary image points to the
normalization is asserted.

The finite factorization map used there is flat by miracle flatness:
its target local rings are regular, its source local rings are
Cohen--Macaulay, it is finite, and dimensions agree. These are exactly
the hypotheses of \cite[Tag 00R4]{StacksLocalV154}. This names the
flatness input separately from the incidence-functor argument.
\end{proof}

\begin{lemma}[The coordinate equalizer and coherent conductor]
\label{lem:equalizer-v154}
For the block ideals $I_A$ of Theorem~\ref{thm:squarefree-v153},
the completed image ring is the equalizer of the branch rings with
their restrictions to all pairwise intersections. This assertion
uses the coordinate monomial structure. The completed conductor
calculation there implies equality of the corresponding coherent
ideal sheaves on the squarefree open.
\end{lemma}
\begin{proof}
A monomial has a coefficient on each branch where its support
survives. If it survives on two branches, it also survives on their
coordinate intersection. Pairwise agreement therefore assigns it
one well-defined coefficient. Collecting these coefficients gives
one formal series modulo the intersection of the $I_A$; conversely
such a series has compatible restrictions. Formal convergence causes
no issue, since there are finitely many monomials in each degree.
The claim would be false for general configurations of ideals without
this coordinate argument.

Write $A\subset B$ for the local finite normalization algebra.
The conductor is $\Ann_A(B/A)$, a coherent ideal. Flat completion
preserves the exact finite-module sequence and its annihilator:
compute the latter as the intersection of the kernels of multiplication
on a finite generating set. The completed finite algebra is the
product of the normal branch rings. Thus the conductor and the ideal
of $Y_{g+1}$ have equal completed stalks at every squarefree closed
point. Faithful flatness gives equal stalks, and coherence on a
finite-type complex scheme gives equality of the ideal sheaves.
\end{proof}

\begin{proposition}[Finiteness and uniqueness of the collision diagram]
\label{prop:atlas-functor-v154}
Every branch map in Theorem~\ref{thm:atlas-v153} is finite, and the
completed incidence diagram is determined, up to the unique
isomorphism respecting its incidence data, by its marked-divisor
functor. It is independent of the chosen frame and root coordinates.
\end{proposition}
\begin{proof}
For each cluster $P_i=A_iB_i'$, the coefficients of the monic $A_i$
are integral over the coefficients of $P_i$. This follows from the
finite factorization morphism, or by expressing them as elementary
symmetric functions of a selected subset of the roots in a finite
splitting algebra. Division by a monic polynomial expresses every
coefficient of $B_i'$ polynomially in those of $A_i$ and $P_i$.
Likewise $R_{ij}=A_iU_{ij}$ expresses all quotient coefficients in
terms of $A_i$ and $R_{ij}$. Consequently the branch algebra is
generated over the ambient coefficient algebra by finitely many
integral coefficients of the $A_i$, so it is finite. Completing
preserves this finite-module assertion.

The equations represent the functor of a subbundle with a marked
divisor and its uniquely determined quotient. A change of frame or
local coordinate gives inverse natural transformations of this same
functor and of its forgetful map to the ambient Grassmannian.
The complete representing rings and their maps are uniquely
isomorphic by Yoneda. Their scheme-theoretic images consequently
agree as completed diagrams. This proves coordinate independence
without identifying arbitrary Artin points of an image with lifted
Artin points of its normalization.
\end{proof}

\paragraph{Partition closures.}
For a fixed partition $(s_1,\ldots,s_b)$, the proper map
$(\Pj^1)^b\to\Pj(S_s)$ sending $(\ell_i)$ to
$\prod\ell_i^{s_i}$ has closed image. Its open set of distinct
$\ell_i$ is dense. At the boundary precisely some of these points
coincide, and the associated parts add. Every such coincidence can
be approached from distinct points. This proves exactly the merging
closure order in the fixed-degree symmetric product; compare the
coincident-root constructions of \cite{Chipalkatti2003,Kurmann2012}.
It is not a closure classification for all analytic singularities
of the completed collision atlas.

\paragraph{Equivariance of the singular Fitting calculation.}
The group used in Theorem~\ref{thm:fitting-v153} is
$\operatorname{GL}_m(\C)$, acting simultaneously on the columns
$u$ and $w$ and fixing $\Delta$. It preserves the presentation and
therefore its intrinsic Fitting ideal. In the normalization the
three displayed minors span, under this group, respectively
$\Sym^m(\C^m)t$, $\Delta\Sym^m(\C^m)$, and
$\Sym^{m+1}(\C^m)$, with the coordinate/dual convention determined
by the action on $u$. To see spanning directly, their translates
are nonzero scalar multiples of the corresponding powers of all
linear forms; polarization in characteristic zero spans each
symmetric power. Multiplication by $R$ gives exactly the modules
asserted in \eqref{eq:fitting-v153}. This is stronger than invoking
invariance without identifying the representations.

\paragraph{Specialization of minimal relations.}
For the surjective coefficient specialization $\pi:P\to P'$ to
one coordinate in Theorem~\ref{thm:reciprocal-structure-v153}, one
has $\pi(\mathfrak n)=\mathfrak n'$ and
$\pi(\mathcal I)=\mathcal I'$. Consequently
$\pi(\mathfrak n\mathcal I)=\mathfrak n'\mathcal I'$.
A coefficient whose image is a minimal binary relation cannot lie
in $\mathfrak n\mathcal I$. This explicitly justifies the nonzero
irreducible coefficient map used in that proof. All reciprocal
series and socle weights in these arguments use the weight of a
$Q_i$ coefficient equal to $i$. They do not claim ordinary
maximal-ideal Hilbert functions; \eqref{eq:F22-local-hilbert-v154}
gives one such function separately.

% BEGIN PRESERVED PART 44-conductor-and-collisions-v152.tex
\section{Conductors and local models of divisor collisions}
\label{sec:conductor-v152}

We retain the notation of Section~\ref{sec:canonical-boundary-v151}.
All completed local rings below are completed at closed complex points.
For a reduced local ring $A$ with finite normalization $B$, its
\emph{conductor} is $\mathfrak c=\Ann_A(B/A)$, regarded also as an
ideal of $B$.  We use the intrinsic singular subscheme defined by
$\Fitt_{\dim A}\Omega_{A/\C}$; on completions, differentials mean
continuous differentials.  The conductor and this Fitting ideal need
not agree.  Our first result determines both on a multivariate
stratum.  The second identifies the way two branches merge when a
common binary root becomes double.

The formal-image and separated-germ steps below are isolated in
Lemma~\ref{lem:formal-image-v153}. Lemma~\ref{lem:coprime-lifts-v153}
gives the Artin cancellation and a direct positive lower bound for the
split codimension. Theorem~\ref{thm:atlas-v153} supplies the full
Weierstrass incidence functor in every collision order;
Theorem~\ref{thm:fitting-v153} evaluates the singular ideal for all ranks.

\subsection{Two coprime divisor choices}
\begin{theorem}[The split-divisor stratum]\label{thm:split-conductor-v152}
Suppose $h>g$, $2\leq r\leq s_{h-g}$, and
\[
 K=abL,\qquad [a],[b]\in\Pj(S_g),\qquad
 L\in\Gr(r,S_{h-g}),\qquad \gcd(L)=1.
\]
Assume that $a,b$ are coprime and that the only degree-$g$ divisors of
$ab$, up to scalar, are $a$ and $b$.  Put
\begin{align*}
 t&=2(s_g-1)+r(s_{h-g}-r),\\
 c&=r(s_h-s_{h-g})-(s_g-1),\qquad d_g=t+c.
\end{align*}
Then $c\geq1$, and there are formal coordinates in which
\begin{equation}\label{eq:split-ring-v152}
 \widehat{\OO}_{Y_g,K}=
 \C[[z_1,\ldots,z_t,x_1,\ldots,x_c,y_1,\ldots,y_c]]/
                       (x_i y_j:1\leq i,j\leq c).
\end{equation}
Its normalization is
$\C[[z,x]]\times\C[[z,y]]$.  The conductor is
\begin{equation}\label{eq:split-conductor-v152}
 \mathfrak c=(x_1,\ldots,x_c,y_1,\ldots,y_c),\qquad
 \widehat{\OO}_{Y_g,K}/\mathfrak c=\C[[z]],
\end{equation}
and the intrinsic singular subscheme has ideal $\mathfrak c^{\,c}$.
The two smooth branches meet scheme-theoretically in a smooth
$t$-dimensional subscheme.  The local ring has multiplicity two and
depth $t+1$; it is Cohen--Macaulay exactly when $c=1$.
\end{theorem}

\begin{proof}
We first determine the intersection as a scheme, not by comparing
closed points or tangent dimensions.  The normalization has precisely
two points above $K$, namely $(a,bL)$ and $(b,aL)$.  At either point,
the overlap in Lemma~\ref{lem:fibre-equations-v151} is zero.  The
map from the completed local ring of the ambient Grassmannian to
each completed source ring is consequently surjective: its cotangent
map is surjective, and completeness and Nakayama's lemma lift
successive homogeneous terms.  Thus each source completion is a
closed smooth formal branch of the image.

Here is the required assertion over infinitesimal bases.  Let $A_0$
be a local Artin $\C$-algebra.  A point on both formal branches carries
unique marked lifts $a_{A_0},b_{A_0}$ near $a,b$.  If homogeneous
polynomials have coprime reductions, those reductions form a regular
sequence in the polynomial ring over $\C$.  Its graded Koszul
complex is exact.  In each degree the polynomial modules are finite
free; lifting a splitting of the complex, or applying the local
criterion for flatness successively, shows that its lift over $A_0$
is exact and that the quotient is $A_0$-flat.  In particular
\begin{equation}\label{eq:coprime-artin-v152}
 (a_{A_0})\cap(b_{A_0})=(a_{A_0}b_{A_0}).
\end{equation}
This also follows by noting that $b_{A_0}$ is a nonzerodivisor
modulo $a_{A_0}$ and solving a relation between them.  Applying
\eqref{eq:coprime-artin-v152} to the columns of the relation
subbundle gives, uniquely,
\[
                  K_{A_0}=a_{A_0}b_{A_0}L_{A_0}.
\]
The residual module is a subbundle: multiplication by
$a_{A_0}b_{A_0}$ is a split injection in degree $D$, and the
quotient of the ambient free module by $K_{A_0}$ is free.  The
kernel of the induced surjection onto the complementary quotient
is therefore free.  This proves that the formal intersection is
represented by the product of the two divisor charts and the
Grassmannian chart for $L$, with no nilpotents omitted.

The differential of this product parametrization is injective.
Indeed, an infinitesimally constant product gives
\[
       (\dot a b+a\dot b)l+ab\dot l\in abL\quad(l\in L).
\]
Taking valuations at the factors of $a$, and using
$\gcd(a,b)=\gcd(L)=1$, forces $a\mid\dot a$; similarly
$b\mid\dot b$.  These are the two projective scalar directions,
and then $\dot L=0$.  The intersection is thus a smooth formal
closed subscheme of dimension $t$.

Let $R_0$ be the ambient completed local ring and $I_a,I_b$ the
ideals of these branches.  Finiteness of the normalization, exactness
of completion on finite modules, and the definition of the
scheme-theoretic image give
\[
 \widehat{\OO}_{Y_g,K}=R_0/(I_a\cap I_b)
   = (R_0/I_a)\mathbin{\times}_{R_0/(I_a+I_b)}(R_0/I_b).
\]
This is the elementary conductor-square construction; see also
\cite[Lemma~1.3 and \S1.3.2]{Ferrand2003}.  All three quotient rings
are regular.  Choosing compatible regular parameters for the two
surjections gives \eqref{eq:split-ring-v152}.  The branch dimension
is $d_g=s_g-1+r(s_h-r)$, so its codimension along the intersection
is the asserted $c$.  It is positive: if $c=0$, both branch ideals
would coincide locally with the intersection ideal, contradicting
the two distinct normalization branches of a reduced local ring.

In the fibre-product description, an element belongs to the
conductor exactly when both its restrictions to the intersection
vanish.  This gives \eqref{eq:split-conductor-v152}.  The Jacobian
of the relations $x_i y_j$ has nonzero entries only among the $x_i$
and $y_j$.  The Fitting ideal in question is its ideal of $c$-minors.
Every such minor belongs to $\mathfrak c^{\,c}$.  Conversely,
choosing the $c$ columns indexed by the $y_j$ and one row $(i_j,j)$
for each $j$ gives every monomial $x_{i_1}\cdots x_{i_c}$.
The symmetric choice gives every degree-$c$ monomial in the $y$'s.
Mixed monomials vanish in the ring.  Thus the ideal is exactly
$\mathfrak c^{\,c}$, including its nonreduced structure.

Finally use
\[
 0\longrightarrow \widehat{\OO}_{Y_g,K}
 \longrightarrow\C[[z,x]]\oplus\C[[z,y]]
 \longrightarrow\C[[z]]\longrightarrow0,
\]
where the last arrow is the difference of restrictions.  The depth
lemma gives depth $t+1$; for $c=1$ this equals the dimension.
For $c>1$ it is strictly smaller.  Additivity of leading Hilbert
coefficients gives multiplicity two, since the last module has
smaller dimension than the two regular branches.
\end{proof}

\begin{remark}\label{rem:split-scope-v152}
The hypotheses are nonempty in every number of variables: for $e=2$
one may take $g=1$, and for $e\geq3$ one may take two distinct
irreducible forms of degree $g$.  The theorem treats an explicitly
specified stratum, not all pairs of overlapping degree-$g$ divisors.
In particular, it does not substitute a binary factorization rule
for multivariate divisibility.  Its high-codimension branch
intersections need not be normal-crossing divisors.
\end{remark}

\subsection{A double common root}
For $m\geq1$, put $R_m=\C[[u_1,\ldots,u_m,\Delta]]$ and
$I=(u_1,\ldots,u_m)\subset R_m$.  Define
\begin{equation}\label{eq:pinch-algebra-v152}
 A_m=R_m\oplus I t\ \subset\ B_m=\C[[u_1,\ldots,u_m,t]],
                    \qquad \Delta=t^2.
\end{equation}
The right side of \eqref{eq:pinch-algebra-v152} is a subring,
not a square-zero extension: $(It)^2=\Delta I^2$.  For $m=1$
it is the ordinary pinch-point ring.

\begin{theorem}[Collision of two binary divisor choices]
\label{thm:binary-pinch-v152}
Let $e=2$, $g=1$, $D=h+1$, and $2\leq r\leq D-1$.
At a point $K_0=\ell^2 L$ with
$L\in\Gr(r,S_{D-2})$ and $\gcd(L)=1$, set
\[
             m=r-1,\qquad q=r(D-1-r)+1.
\]
Then
\begin{equation}\label{eq:binary-completion-v152}
              \widehat{\OO}_{Y_1,K_0}\simeq A_m[[z_1,\ldots,z_q]].
\end{equation}
Writing $w_i=u_i t$, a complete presentation of $A_m$ is
\begin{equation}\label{eq:pinch-equations-v152}
 \C[[u_1,\ldots,u_m,\Delta,w_1,\ldots,w_m]]/
 \bigl(u_iw_j-u_jw_i,\ w_iw_j-\Delta u_i u_j\bigr),
\end{equation}
where the first relations have $i<j$ and the second have $i\leq j$.
Its normalization is $B_m$, and its conductor is
\[
 \mathfrak c=(u_1,\ldots,u_m,w_1,\ldots,w_m),\qquad
 \mathfrak c B_m=I B_m.
\]
The induced map on conductor quotients is
\begin{equation}\label{eq:conductor-cover-v152}
              \C[[\Delta]]\longrightarrow\C[[t]],\qquad\Delta\longmapsto t^2.
\end{equation}
Thus the two reduced divisor lifts merge through a ramified double
cover of the conductor.  The completed local ring has dimension
$q+m+1$ and depth $q+2$.  It is Cohen--Macaulay if and only if $r=2$.
For $r=2$, its intrinsic singular ideal, with the smooth variables
suppressed, is $(w,\Delta u,u^2)$, whereas its conductor is $(u,w)$.
\end{theorem}

\begin{proof}
Choose an affine coordinate $x$ with the root of $\ell$ at zero.
Some member $F_0\in K_0$ has multiplicity exactly two there, since
$\gcd(L)=1$.  Extend it to a frame $F_0,\ldots,F_{r-1}$.
In the completed local ring of the ambient Grassmannian the formal
Weierstrass theorem writes the universal $F_0$ as a unit times
\[
                         P(x)=x^2+a x+b.
\]
Divide each $F_i$, $i>0$, by this monic polynomial.  Its remainder
is $u_i x+v_i$.  The $2r$ functions
$a,b,u_1,v_1,\ldots,u_m,v_m$ are part of a regular parameter
system.  To see this rather than assume versality, note that
$K_0$ has zero two-jet at the root, while
$S_D/K_0\to\C[x]/(x^2)$ is surjective.  The Grassmannian tangent
space $\Hom(K_0,S_D/K_0)$ therefore allows independent variations
of both jets in every frame column.  Passing from these jets to
the displayed remainders is an invertible triangular operation,
with leading coefficient $F_0/x^2(0)\ne0$.  Their differentials
are consequently independent.

Set
\[
       t=x+a/2,\quad \Delta=a^2/4-b,\quad
       w_i=a u_i/2-v_i.
\]
The incidence of a common root near zero is now exactly
\[
                       \Delta=t^2,\qquad w_i=u_i t.
\]
There are no other normalization points above $K_0$, because its
exact common divisor is $\ell^2$.  Consequently the image of this
completed incidence is the completed image defining $Y_1$.
The ambient dimension is $N_0=r(D+1-r)$.  Of the $2r$ parameters
above, $a$ is free, as are the other $N_0-2r$ parameters.  This
accounts for $q=N_0-2r+1$ smooth variables in
\eqref{eq:binary-completion-v152}.

It remains to justify elimination scheme-theoretically.  The
relations in \eqref{eq:pinch-equations-v152} reduce any series to
$p+\sum_i w_i p_i$ with $p,p_i\in R_m$.  If its image in $B_m$
is zero, the decomposition $B_m=R_m\oplus R_m t$ gives
$p=0$ and $\sum_i u_i p_i=0$.  The first syzygies of the regular
sequence $u_1,\ldots,u_m$ are generated by the Koszul syzygies.
They are precisely the relations $u_iw_j-u_jw_i$.  This proves
the presentation without taking a radical.  The same argument
works first over polynomial rings and then by exact completion.

The extension $A_m\subset B_m$ is finite, since $t^2=\Delta$,
and birational, since $t=w_1/u_1$ in the fraction field.  The
regular ring $B_m$ is thus the normalization.  For
$p+v t\in R_m\oplus I t$, multiplication by $t$ stays in
$A_m$ if and only if $p\in I$.  As $1,t$ generate $B_m$,
the conductor is $I\oplus I t$, proving the quotient map
\eqref{eq:conductor-cover-v152}.  Equivalently,
\begin{equation}\label{eq:pinch-square-v152}
 A_m=B_m\mathbin{\times}_{\C[[t]]}\C[[\Delta]],
\end{equation}
where the first map sets $u=0$ and the second sends $\Delta$ to
$t^2$.  This is a specific conductor square, not an assumption
that the normalization descends through its reduced fibres.

As an $R_m$-module, $A_m=R_m\oplus I$.  For $m=1$, the ideal
$I$ is free.  For $m>1$, the exact sequence
$0\to I\to R_m\to\C[[\Delta]]\to0$ gives
$\operatorname{depth}_{R_m}I=2$.  Hence $A_m$ has depth two
for every $m\geq1$.  The maximal ideal of $R_m$ generates an
ideal with maximal radical in $A_m$, so this is also its local
ring depth.  Adding the $q$ smooth parameters gives the assertion.
When $m=1$, the equation is $w^2-\Delta u^2=0$; its partial
derivatives generate $(w,\Delta u,u^2)$ in characteristic zero.
\end{proof}

\begin{corollary}[The split and ramified parts of the same model]
\label{cor:collision-strata-v152}
On the binary locus where the common divisor has degree exactly two,
there are two strata.  Over two distinct roots the formal local ring
has the split form \eqref{eq:split-ring-v152}, with
$t=q+1$ and $c=r-1$.  Over a double root it has
\eqref{eq:binary-completion-v152}.  The support of the singular
locus in the latter model is the conductor, and the normalization
fibre at a collision is $\Spec\C[t]/(t^2)$.
For arbitrary $m$, its singular subscheme is given without
radicalization by the $m$-minors of the Jacobian of the displayed
$m^2$ relations in \eqref{eq:pinch-equations-v152}.
\end{corollary}
\begin{proof}
The assertion at distinct roots is Theorem~\ref{thm:split-conductor-v152}.
In the collision model the normalization is an isomorphism wherever
some $u_i$ is invertible.  On $u=w=0$, a nonzero value of $\Delta$
gives two distinct points after an \'etale square-root extension,
while $\Delta=0$ gives the stated length-two fibre.  The exact
normality criterion proves the assertion about singular support.
There are $2m+1$ nonsmooth coordinates and dimension $m+1$;
the cotangent presentation therefore gives the $m$-minors as
claimed.  This formula specifies the whole singular scheme, even
where it is thicker than its conductor support.
\end{proof}

\begin{remark}[Descent to intrinsic first relations]
\label{rem:conductor-stack-v152}
On the maximal-lower-Hilbert-function stratum of
Corollary~\ref{cor:algebra-boundary-v151}, these statements compute
the completed smooth-atlas singularities of the intrinsic
first-relation algebra stack after tangent-identity rigidification.
The conductor is equivariant, and formation of its finite-module
annihilator and of the Fitting ideals commutes with flat base change.
Thus the calculations descend through the frame atlas.  This does
not identify these rings with completed local rings of a coarse
orbit space, where stabilizers can introduce further singularities.
\end{remark}

% BEGIN PRESERVED PART 46-divisor-intersections-v152.tex
\section{Intersections of divisor-degree images}
\label{sec:degree-intersections-v152}

Here the total relation degree $D$ is fixed.  Write $Y_a\subset
\Gr(r,S_D)$ for the image admitting a divisor of degree $a$;
empty Grassmannians are omitted.  In several variables these images
are not generally nested.  Their reduced intersections nevertheless
have a finite description by common parts of two marked divisors.

\begin{proposition}[The common-part formula]
\label{prop:degree-intersections-v152}
For $1\leq a,b<D$ and $0\leq j\leq\min(a,b)$, let $M_j$ be the
scheme-theoretic image of
\begin{equation}\label{eq:intersection-map-v152}
\begin{split}
 \Pj(S_j)\times\Pj(S_{a-j})\times\Pj(S_{b-j})
       \times\Gr(r,S_{D-a-b+j})&\longrightarrow\Gr(r,S_D),\\
                   (c,u,v,L)&\longmapsto cuvL.
\end{split}
\end{equation}
A negative last degree or insufficient dimension gives the empty
image, and $\Pj(S_0)$ is a point.  Then
\begin{equation}\label{eq:intersection-reduced-v152}
             (Y_a\cap Y_b)_{\mathrm{red}}
                         =\left(\bigcup_j M_j\right)_{\mathrm{red}}.
\end{equation}
On the locus of two marked divisors $f_a,f_b$ with exact common
part of degree $j$, the corresponding factorization is
$c=\gcd(f_a,f_b)$, $f_a=cu$, $f_b=cv$, with $\gcd(u,v)=1$.
In particular, at a relation $K=a_0^mL$ with $a_0$ a nonzero
linear form and $\gcd(L)=1$, the
set of positive divisor degrees is exactly $\{1,\ldots,m\}$.
For the boundary point in Theorem~\ref{thm:universal-boundary-v152},
$m=D-q_n$.
\end{proposition}
\begin{proof}
Each map in \eqref{eq:intersection-map-v152} is defined on a
projective scheme: multiplication by a nonzero form has constant
rank, so the subspace construction is a morphism.  Its image is
closed and admits the two divisors $cu$ and $cv$, proving one
inclusion.  Conversely, if $K$ has divisors $f_a,f_b$, unique
factorization gives their common part $c$ of some degree $j$.
Their least common multiple is $cuv$ and divides every member of
$K$.  Division yields $L$ of degree $D-a-b+j$, of the same rank.
This proves equality on geometric points.  Over $\C$, a reduced
closed subscheme of a finite-type scheme is determined by its closed
points, giving \eqref{eq:intersection-reduced-v152}.  The last
assertion follows because every homogeneous divisor of a power
of one linear form is a power of that form.
\end{proof}

The reduction in \eqref{eq:intersection-reduced-v152} is essential
to the statement of this global formula.  It is not a prescription
to reduce the individual $Y_a$, their normalization fibres, or their
singular subschemes.  The full intersection of the two formal
branches on the coprime stratum was computed over Artin bases in
Theorem~\ref{thm:split-conductor-v152};
Theorem~\ref{thm:binary-pinch-v152} computes the thickening when
those branches collide.  Thus the degree-poset description and
the local scheme calculations have distinct, explicit roles.

% BEGIN PRESERVED PART 45-pencil-orbit-boundary-v152.tex
\section{A common ramified boundary of pencil failure orbits}
\label{sec:orbit-boundary-v152}

The fixed tensor presentation of a pencil has exact common divisor
$\delta^{n-1}$.  Varying the pencil alone therefore stays in the
exact-divisor locus.  The intrinsic finite algebra permits a larger,
linear coordinate-change action: changing the basis of its entire cotangent
space.  We now describe a boundary point produced by such changes.
The generic members are genuine pencil failure algebras.  The central
member has a larger common divisor, and its normalization fibre is
nonreduced.  Thus the boundary of Section~\ref{sec:canonical-boundary-v151}
occurs inside the orbit closures of the motivating algebras themselves.

Let $n=\dim V=\dim U\geq3$, $E=V^*\otimes U$, and put
\[
 e=n^2,\quad N=\binom{n+1}{2},\quad r=\binom N2,\quad
 g=n(n-1),\quad h=3n-4,\quad D=g+h.
\]
For a pencil $R$, write $K_R=\delta^{n-1}C_R\subset\Sym^D E$
as in Lemma~\ref{lem:primitive-pencil}.  Define the projective orbit
closure
\[
 Z_R=\overline{\operatorname{GL}(E)\cdot K_R}
                  \ \subset\ \Gr(r,\Sym^D E).
\]
It is independent of a chosen basis of the cotangent space.  Since
$Y_g$ is closed and invariant, $Z_R\subset Y_g$.

The group in $Z_R$ is the full $\operatorname{GL}(E)$, not the
$\PGL(V)$ congruence group of the pencil. Fixing the standard flag,
scalar direction and complement specifies the displayed common point;
no choice-independent distinguished moduli point is asserted.

\begin{theorem}[A common point in full first-relation orbit closures]
\label{thm:universal-boundary-v152}
Set
\[
 q_n=\begin{cases}3,&n=3,\\4,&n\geq4,\end{cases}
 \qquad k_n=h-q_n.
\]
There is an explicitly constructed point $K_\infty$ belonging to
$Z_R$ for every quadratic pencil $R$, with
\begin{equation}\label{eq:boundary-gcd-v152}
 K_\infty=a^{g+k_n}L_n,\qquad
 L_n\subset\Sym^{q_n}E,\quad \rank L_n=r,\quad\gcd(L_n)=1,
\end{equation}
where $a\in E$ is a nonzero linear form.  The flag pencil
$R_*=\langle x_1^2,x_1x_2\rangle$ admits a finite-flat family of
first-relation algebras over $\A^1$, with special relation space
$K_\infty$ and every nonzero fibre isomorphic to its failure algebra.
Together with Theorem~\ref{thm:closed-pencil-v151}, this gives a
chain of two explicit finite-flat specializations from the failure
algebra of every $R$ to the same boundary algebra.

The fibre of $\nu_g$ over $K_\infty$ has a single geometric point,
marked by $a^g$, and is nonreduced.  Its tangent dimension is
\begin{equation}\label{eq:boundary-tangent-v152}
                 \binom{e+k_n-1}{k_n}-1.
\end{equation}
In particular $K_\infty$ lies in the singular locus of $Y_g$ and
outside the exact-degree-$g$ locus.  For $n=3$ this tangent dimension
is $44$.
\end{theorem}

The construction requires the exact osculating order of one elementary
flag embedding.  We include the calculation to distinguish the theorem
from a computationally suggested degeneration.

\begin{lemma}[The flag coefficient filtration]
\label{lem:flag-jets-v152}
Write a matrix of independent linear coordinates as
\[
                    T=aI_n+B,\qquad B_{11}=0.
\]
For $C_*=C_{R_*}$, the filtration by vanishing order at $T=I_n$
on the chart $T_{11}=1$ has last nonzero graded piece in degree $q_n$.
It has a nonzero degree-zero piece.  Hence an adapted basis
$c_1,\ldots,c_r$ has orders $0\leq d_j\leq q_n$, attaining both
endpoints, and
\begin{equation}\label{eq:flag-initial-v152}
 \lim_{\tau\to0}\tau^{-d_j}c_j(aI_n+\tau B)
                           =a^{h-d_j}p_j(B),
\end{equation}
where the displayed initial forms are linearly independent.
\end{lemma}
\begin{proof}
The complementary-minor identity for $\Sym^2T$ identifies the vector
of primitive coefficients, up to a fixed invertible coordinate
change, with
\begin{equation}\label{eq:flag-coefficients-v152}
 \delta(T)^3\bigl(v^2\wedge vw\bigr),\qquad
                 v=T^{-1}e_1,\quad w=T^{-1}e_2.
\end{equation}
Indeed $\det(\Sym^2T)=\delta^{n+1}$, the complementary minors have
order two, and division of $\delta I_{N-2}$ by $\delta^{n-1}$ leaves
$\delta^3$.  This identity is initially on $\delta\ne0$; both sides
are the same polynomial vector of degree $h$.

A standard affine flag chart has representatives
\[
       v=e_1+\sum_{i=2}^n z_i e_i,
       \qquad w=e_2+\sum_{i=3}^n y_i e_i.
\]
Every coordinate of $v^2\wedge vw$ is a polynomial of degree at most
four.  The map from $\operatorname{GL}_n$ to this flag chart is
smooth, and restricting to $T_{11}=1$ remains smooth: the scalar
subgroup is contained in each flag stabilizer and is transverse to
this chart.  Normalizing $v,w$ only multiplies the whole wedge by
a common unit.  Smooth pullback preserves the order of a nonzero
function, as does multiplication by a unit.  The coordinates in
\eqref{eq:flag-coefficients-v152} are linearly independent by Lemma~\ref{lem:pencil-irreducibility} and the nonzero
matrix-coefficient inclusion in \eqref{eq:general-coefficient-map}.  Thus no nonzero linear
combination can have order exceeding four.  For $n\geq4$, the
coordinate indexed by $e_3^2\wedge e_3e_4$ is
\[
                       z_3^2(z_3y_4-z_4y_3),
\]
which is nonzero and has order exactly four.

For $n=3$, put $z=z_2$, $w=z_3$, $t=y_3$.  In the ordered basis
$e_1^2,e_1e_2,e_1e_3,e_2^2,e_2e_3,e_3^2$, the two factors are
\[
 (1,2z,2w,z^2,2zw,w^2),\qquad
 (0,1,t,z,zt+w,wt).
\]
Their fifteen wedge coordinates have the same linear span as
\begin{equation}\label{eq:ternary-flag-basis-v152}
\begin{gathered}
 1,\ z,\ w,\ t,\ z^2,\ zw,\ w^2,\ zt,\ wt,\ z^2t,\ zwt,\ w^2t,\\
 z^2w-z^3t,\quad zw^2-z^2wt,\quad w^3-zw^2t.
\end{gathered}
\end{equation}
This is obtained by subtracting the two coordinates with a common
quadratic term and rescaling by $2$.  Their lowest homogeneous terms
are distinct monomials, all of degrees at most three, and some have
degree three.  Thus the third jet is injective on this coefficient
space and the second is not.  This proves $q_3=3$.  The coordinate
$e_1^2\wedge e_1e_2$ is $1$ in every dimension, proving the
order-zero assertion.  Finally, homogeneity of degree $h$ gives
\eqref{eq:flag-initial-v152}; an adapted basis is by definition a
basis on each associated-graded piece, so its initial forms are
independent.
\end{proof}

\begin{proof}[Proof of Theorem~\ref{thm:universal-boundary-v152}]
The linear substitution $T\mapsto aI_n+\tau B$ is invertible for
$\tau\ne0$.  Divide the adapted coefficient $c_j$ by its exact
power $\tau^{d_j}$ after substitution.  There are no negative powers:
all its terms have $B$-degree at least $d_j$.  The resulting columns
span a rank-$r$ subbundle $J_\tau\subset\Sym^h E\otimes\C[\tau]$.
Independence at zero follows from the lemma, and away from zero from
the invertible substitution.  The ideal of its maximal minors in $\C[\tau]$ has no common
closed zero, and is therefore the unit ideal. On each maximal-minor
open a corresponding square submatrix is invertible and splits the
image. These opens cover $\operatorname{Spec}\C[\tau]$, proving the
subbundle assertion over the whole line, including all scheme points.  Set
\[
                 K_\tau=\det(aI_n+\tau B)^{n-1}J_\tau.
\]
Multiplication by the nonzero homogeneous polynomial in every fibre
again has constant rank.  Consequently $K_\tau$ is a subbundle, and
\[
 \mathcal A_\tau=\Sym E\otimes\C[\tau]\big/
                           (K_\tau,\mathfrak m^{D+1})
\]
is finite locally free of rank $\binom{e+D}{D}-r$.
For $\tau\ne0$, its relation space is a linear transform of $K_{R_*}$.

At zero the marked factor becomes $a^g$, while
$J_0=\langle a^{h-d_j}p_j(B)\rangle$.  It contains a nonzero multiple
of $a^h$, so its gcd is a power of $a$.  It also contains a polynomial
with exact $a$-adic order $h-q_n$, and every column is divisible by
that power.  Therefore $\gcd(J_0)=a^{k_n}$, proving
\eqref{eq:boundary-gcd-v152}.  Notice that $0<k_n<g$ for every
$n\geq3$.

The closed-flag specialization of Theorem~\ref{thm:closed-pencil-v151}
and its failure-algebra lift imply
$K_{R_*}\in Z_R$.  A closed invariant subset containing a point
contains that point's orbit closure.  Hence the just-constructed
$K_\infty\in Z_{R_*}$ belongs to every $Z_R$.  This proves the
assertion about the two-stage specialization without replacing it
by an unproved assertion about one fixed one-parameter subgroup
acting on every initial pencil.

The sole degree-$g$ divisor of $a^{g+k_n}$ is $a^g$.  Its overlap
with the residual gcd is $a^{k_n}$.  Lemma~\ref{lem:fibre-equations-v151}
gives \eqref{eq:boundary-tangent-v152}; the fibre is supported at one
point and has positive tangent space.  The normalization criterion
then proves singularity.  With $n=3$, $e=9$ and $k_n=2$, the dimension
is $\binom{10}{2}-1=44$.
\end{proof}

\subsection{The complete normalization fibre at the limiting algebra}
The last theorem is not restricted to the tangent space of its fibre.
The full local algebra can be given uniformly.

\begin{proposition}[Pure-power divisor fibre]
\label{prop:pure-power-fibre-v152}
Let $E=\C a\oplus W$, $g\geq k\geq1$, and
$K=a^{g+k}L$ with $\gcd(L)=1$.  Introduce coefficient variables for
$q_i\in\Sym^i W$, $1\leq i\leq k$, and define forms $c_j$ recursively by
\begin{equation}\label{eq:reciprocal-v152}
 c_0=1,\qquad c_j=-\sum_{i=1}^k q_i c_{j-i},\qquad c_j=0\ (j<0).
\end{equation}
The entire fibre of $\nu_g$ over $K$, supported at its unique point,
is the Artin algebra
\begin{equation}\label{eq:pure-power-fibre-v152}
 \C[[\text{coefficients of }q_1,\ldots,q_k]]\big/
     (\text{coefficients of }c_{g+1},\ldots,c_{g+k}).
\end{equation}
Its tangent dimension is $s_k(E)-1$.
\end{proposition}
\begin{proof}
Choose $p\in L$ not divisible by $a$.  Over any local Artin base,
a divisor $f$ near $a^g$ is monic in $a$, and the quotient by $f$
is flat over the base.  The reductions of $f,p$ are coprime.
The regular-sequence argument used in
\eqref{eq:coprime-artin-v152} makes $p$ a nonzerodivisor modulo $f$.
Thus $f\mid a^{g+k}p$ is equivalent to $f\mid a^{g+k}$, functorially
on these bases.  Once this condition holds, it holds for every
member of $K$.

Polynomial division gives a unique monic residual factor
$H=a^k+q_1a^{k-1}+\cdots+q_k$.  The equation $fH=a^{g+k}$ determines
the coefficients of $f$ successively as $c_1,\ldots,c_g$ in
\eqref{eq:reciprocal-v152}.  The remaining $k$ product coefficients
vanish if and only if $c_{g+1},\ldots,c_{g+k}$ vanish.  This follows
successively by comparing the product with the reciprocal series
$(1+\sum q_i a^{-i})^{-1}$.  Hence
\eqref{eq:pure-power-fibre-v152} represents the full infinitesimal
fibre, without reduction.  Finiteness of $\nu_g$ makes it Artin;
there is no other geometric divisor, so completion loses no part
of the fibre.  Since $g\geq k$, these equations have no linear
terms in the coefficient variables.  The number of coefficient variables is
$\sum_{i=1}^k\dim\Sym^iW=\dim\Sym^kE-1$, as asserted.
\end{proof}

\begin{example}[The nine-generator limit]\label{ex:n3-boundary-v152}
For $n=3$, the construction can be written with four of the nine
linear coordinates, denoted $a,u,v,w$.  The other five coordinates
remain generators of the algebra.  Up to nonzero scalar changes
of the relation basis, $K_\infty=a^8L_3$, where
\begin{equation}\label{eq:n3-limit-v152}
\begin{split}
 L_3=\langle& a^3,\ a^2u,\ a^2v,\ a^2w,\ au^2,\ auv,\ av^2,
             \ auw,\ avw,\\
            &u^2v,\ u^2w,\ uv^2,\ uvw,\ v^3,\ v^2w\rangle.
\end{split}
\end{equation}
Here one may take $u=B_{21}$, $v=B_{31}$, $w=B_{32}$; the monomials
are the lowest terms in \eqref{eq:ternary-flag-basis-v152}.
The relation space has dimension $15$ and exact gcd $a^8$.
The selected divisor has degree $g=6$, and its fibre has $44$
coefficient variables: eight coefficients of a linear form $Q_1$
and thirty-six of a quadratic form $Q_2$ in $W$.
Its defining ideal consists of all coefficients of the two forms
\begin{align*}
 c_7&=-Q_1^7+6Q_1^5Q_2-10Q_1^3Q_2^2+4Q_1Q_2^3,\\
 c_8&=Q_1^8-7Q_1^6Q_2+15Q_1^4Q_2^2-10Q_1^2Q_2^3+Q_2^4.
\end{align*}
This is a full presentation of the normalization fibre.  Its embedding
dimension $44$ is not being identified with its length, nor is this
fibre algebra being identified with the completed local ring of $Y_g$.
\end{example}

\begin{remark}[Relation to pencil invariants]\label{rem:segre-boundary-v152}
The action defining $Z_R$ is the intrinsic $\operatorname{GL}(E)$
action on first relations, not merely the congruence action on $V$.
The flag pencil is singular: every member has rank at most two.
It is therefore outside the regular-pencil open set treated by
\cite[Theorem~1.1 and \S5]{FevolaMandelshtamSturmfels}, where Segre
symbols organize classification and degenerations.  The theorem
above does not purport to classify that poset or all orbits in $Z_R$.
Its conclusion is instead a uniform common boundary and an exact
ramification algebra for genuine finite failure algebras.  The
increase from $g$ to $g+k_n$ in the intrinsic gcd is precisely what
is invisible in a fixed-determinant pencil specialization.
\end{remark}

\appendix

% BEGIN PRESERVED PART 00-principal-introduction-v152.tex
\section{Detailed reconstruction statements and comparisons}
\label{sec:introduction-v132}

A degeneracy locus usually forgets the map that defines it.  We study
an inverse problem in which neither the ambient embedding nor the
normal directions are retained: can the abstract nonreduced failure
scheme recover a quadratic pencil?  The answer is affirmative for
every complex pencil, including singular pencils.  The first relation, rather than its reduced determinant support,
carries this information.  The inverse is also global: a thickening over a projective reduction
recovers an entire varying pencil subbundle and its actual source
bundle, up to one constant left transformation.  No fibrewise choice
of coordinates or projective line-bundle ambiguity is left.  A
single unmarked local algebra suffices for the constant-pencil
problem; a Schubert line gives another geometric realization.

Let \(V\) be a complex vector space of dimension \(n\ge3\), and put
\[
 W=\Sym^2V,\qquad N=\binom{n+1}{2},\qquad
 R\in\Gr(2,W),\qquad S_R=W/R,\qquad p=N-2.
\]
The cube-zero algebra \(A_R=\C\oplus V\oplus S_R\) has multiplication
\(\Sym^2V\twoheadrightarrow S_R\), with
\((V\oplus S_R)S_R=0\).  On \(X_R=\Gr(n,V\oplus S_R)\), let
\(\mathcal K\) denote the tautological bundle.  For any \(m\ge2\), set
\begin{equation}\label{eq:failure-intro}
 \widehat D_R=V\!\left(\Fitt_0\coker[
 \Sym^m(\OO\oplus\mathcal K)\longrightarrow A_R\otimes\OO]\right).
\end{equation}
This is the maximal-minor scheme of multiplication, not merely its
support.  It is independent of \(m\ge2\).  Its reduction is the
Schubert divisor where \(\mathcal K\to V\) drops rank.  Write
\(B_R=\Gr(n,S_R)\), let \(\mathfrak b_R\) be its ideal in
\(\widehat D_R\), and define
\[
 Z_R^{[k]}=V(\mathfrak b_R^{k+1})\subset\widehat D_R,
       \qquad d=n+2p=n^2+2n-4.
\]
An isomorphism of these neighbourhoods means an abstract complex
scheme isomorphism; it is supplied with no ambient or tensor marking.

\begin{theorem}[Sharp intrinsic reconstruction]
\label{thm:principal-finite-v143}
For all quadratic pencils \(R,R'\subset\Sym^2V\),
\[
 Z_R^{[d]}\simeq Z_{R'}^{[d]}
       \quad\Longleftrightarrow\quad
 R'=gR\quad\text{for some }g\in\PGL(V).
\]
For every \(0\le k<d\), the isomorphism class of \(Z_R^{[k]}\)
is independent of \(R\).  Thus \(d\) is the smallest uniform
reconstruction order.  The same statements hold after restriction,
using the graph projection, to any Schubert line in \(B_R\).
The resulting abstract curve-supported neighbourhood is independent
of the choice of that line.
\end{theorem}

The Grassmannian statement is Theorem~\ref{thm:sharp-finite-pencil}.
The curve-supported statement is
Corollary~\ref{cor:schubert-line-reconstruction-v145}.  Its more general
form, Theorem~\ref{thm:curve-reconstruction-v145}, applies over any
smooth connected projective curve with a rank-\(n\) bundle having
nontrivial projectivization.  For example,
\(\OO_C(-P)\oplus\OO_C^{n-1}\) works on every such pointed curve.
The curve is intrinsic as the reduction; a chosen embedding of the
curve is not part of the isomorphism data.

\subsection{A moving family, not only its individual fibres}
The same finite object also reconstructs a pencil that varies over its
projective reduction.  Let \(B\) be a smooth connected projective
complex variety, let \(\mathcal U\) have rank \(n\), and let \(\mathcal R\) be a rank-two subbundle of
\(\Sym^2V\otimes\OO_B\) with locally free quotient.  Use the local
ideal \((\det T)I_p(\gamma\Sym^2T)\) in
\(\Hom(\mathcal U,V)\) and truncate at order \(d\).
Theorem~\ref{thm:unrestricted-moving-v146} states that two resulting
abstract finite schemes are isomorphic exactly when their data are
related by an isomorphism of reduced bases, one constant linear left
map, and an actual isomorphism of the source bundles, carrying one
pencil subbundle to the other.  The projection and both tensor factors
are forgotten in the input.  The varying coefficient line is recovered
as a subbundle, not by pointwise choices of projective coordinates.

This strengthens the constant-pencil curve theorem in two directions:
the reduced base and source bundle may vary in the comparison, and the
pencil need not be constant on that base.  Example~\ref{ex:nonconstant-pencil-v146}
has a nonconstant unordered cross-ratio and six spectral collision
points, all retained by the same finite scheme.  Theorem~\ref{thm:automorphism-kernel-v146}
then identifies the remaining freedom in an isomorphism.  Its quotient
is the group of the recovered linear data; its kernel acts trivially
on the associated graded scheme and has a terminating commutator
filtration.  Explicit nonlinear shears show why that kernel cannot be
omitted.  These statements are consequences of the unmarked inverse,
not additional claims of novelty for the classical Pluecker map or
for spectral cross-ratios.

\subsection{A single unmarked local algebra suffices}
The strongest form of the inverse needs no positive-dimensional support.
Choose an \(n\)-dimensional source space \(U\), write \(T=(t_{ij})\)
for its universal map to \(V\), and let \(\mathfrak m=(t_{ij})\).
Define
\[
 \mathfrak A_R^{[d]}=
 \C[t_{ij}]/\bigl((\det T)I_p(\gamma_R\Sym^2T)
                                      +\mathfrak m^{d+1}\bigr).
\]
\begin{theorem}[Local algebraic inverse]\label{thm:principal-local-v146}
For every pair of complex quadratic pencils, an ungraded algebra
isomorphism \(\mathfrak A_R^{[d]}\simeq\mathfrak A_{R'}^{[d]}\)
exists if and only if \(R'=gR\) for some \(g\in\PGL(V)\).
The order \(d=n^2+2n-4\) is the smallest uniform order already for
these single-point objects.
\end{theorem}
The full statement and proof are
Theorem~\ref{thm:artin-local-inverse-v146}.  Proposition~\ref{prop:first-relation-orbit-v147}
separates the general truncation argument from the pencil-specific
orbit rigidity: nonlinear coordinate changes form an explicitly described
unipotent kernel, whereas the first relation space carries the
reconstruction data.  The extra local ingredient
is that the residual coefficient module has factor-support ranks
\((1,\binom N2)\).  Transposition reverses these unequal ranks and
therefore cannot carry one pencil ideal to another.  This orients the
unordered matrix factors without a nontrivial projective bundle.
Theorem~\ref{thm:unrestricted-moving-v146} consequently removes that
restriction from the moving-family theorem as well.  This is not a
claim that the reduced determinant cone alone remembers the pencil:
the nonreduced first relation kernel and its residual coefficients
are indispensable.  Proposition~\ref{prop:local-automorphisms-v146}
also identifies the entire tangent-identity kernel, including its
affine-space parametrization and nonadditive group law.

\subsection{Conductors, collision singularities, and actual failure orbits}
Normalization alone does not describe the gluing of its branches.
Theorems~\ref{thm:split-conductor-v152} and
\ref{thm:binary-pinch-v152} determine two successive local models.
On the coprime two-divisor stratum the completed image is
\[
       \C[[z,x_1,\ldots,x_c,y_1,\ldots,y_c]]/(x_i y_j).
\]
Its conductor is $(x,y)$, its intrinsic singular ideal is $(x,y)^c$,
and its depth is $\dim(z)+1$.  When two common binary roots collide,
the model is instead
\[
 \C[[u,\Delta]]\oplus(u_1,\ldots,u_m)t
          \ \subset\ \C[[u,t]],\qquad \Delta=t^2,
\]
up to a specified smooth factor.  Its conductor has a ramified
double cover in the normalization.  These formulas determine the
Cohen--Macaulay cases and distinguish the singular scheme from its
reduced support.  Proposition~\ref{prop:degree-intersections-v152}
organizes the reduced intersections of different divisor-degree
images by the common part of two chosen divisors.

The boundary also has a direct role in the pencil inverse.
Theorem~\ref{thm:universal-boundary-v152} constructs a common point
in every intrinsic first-relation orbit closure of a pencil failure
algebra.  A scalar-matrix contraction of the flag-pencil relation
has gcd $a^{D-q_n}$, where $q_3=3$ and $q_n=4$ for $n\geq4$.
The selected degree-$n(n-1)$ divisor is unique but ramified; its
normalization fibre has tangent dimension
$\binom{n^2+3n-5-q_n}{3n-4-q_n}-1$.
Proposition~\ref{prop:pure-power-fibre-v152} gives its entire Artin
algebra by reciprocal-polynomial equations.  This specializes
\emph{actual} failure algebras through flat families; it does not
merely place an unrelated common-factor example beside the inverse
theorem.  The exact order of the flag coefficient filtration, proved
in Lemma~\ref{lem:flag-jets-v152}, is the link between the pencil
representation and the divisor collision.

\subsection{Factorization loci and the geometric comparison}
The normalization mechanism has classical precedents.
Chipalkatti's coincident-root locus parametrizes one binary form
with a prescribed root partition.  His Proposition~5.1, Theorem~5.4,
and Corollary~5.8 distinguish multiple normalization points from
failure of tangent injectivity and describe the singular support
\cite{Chipalkatti2003}.  Kurmann gives local incidence equations,
proves their scheme-theoretic equality in Theorem~1.5, and expresses
smoothness through splittings of partitions in Proposition~2.1
\cite{Kurmann2012}.  In particular, uniqueness plus tangent
injectivity is not a new general principle of the present paper.

Here the datum is an $r$-plane of forms in $e$ variables, not one
binary form; the selected divisor has prescribed degree, not a
root partition.  The residual datum is a Grassmannian subspace,
and for $e>2$ divisor degrees need not fill an interval.  These
changes produce high-codimension conductor strata and non-Cohen--Macaulay
branch gluings.  Theorems~\ref{thm:split-conductor-v152} and
\ref{thm:binary-pinch-v152} compute the image rings, not merely
the incidence graph or singular support.  Ferrand's conductor-square
formalism \cite{Ferrand2003} is the standard algebraic framework;
the particular conductor ideals, singular Fitting ideals, and depths
are calculated here for these Grassmannian strata.

There is a second, closer binary comparison.  Hu--Lin--Shao
\cite[\S2.2, Corollary~3.6, and Theorem~1.1]{HuLinShao2007}
study tuples of binary forms, define common-factor strata using
resultant homomorphisms, and resolve the map-space boundary by
successive blow-ups.  On the open set of linearly independent
$r$-tuples, passing to their span is the frame quotient onto
$\Gr(r,S_D)$; the common-divisor condition is invariant under this
quotient.  Thus our binary loci genuinely meet that theory rather
than forming an unrelated example.  We do not identify its
resultant scheme structure with our image structure without a
scheme comparison, nor claim a new general resolution.  Our
multivariate local rings and the pencil-orbit boundary are the
specific additional conclusions.

Rigidification is used as infrastructure, not as an independent
headline innovation.  The appropriate primary general result is
Abramovich--Olsson--Vistoli, Appendix~A, Theorem~A.1
\cite{AOV2008}, for compatible flat finitely presented normal
inertia subgroups, without a centrality assumption.  Our argument
identifies the particular substitution subgroup and twisted
right-factor subgroup and retains their actual group laws.  The
boundary statement is on the maximal-lower-Hilbert-function
first-relation stratum throughout.

The comparison with regular-pencil Segre strata is made in
Remark~\ref{rem:segre-boundary-v152}.  The common closed flag orbit
is not presented as a new classification of congruence orbits.
The new pencil-side conclusion is its explicit ramified boundary
under intrinsic cotangent changes, together with the complete
normalization-fibre algebra.  The separate documentary status of
Ballico's 1993 paper, discussed below, remains unchanged: absent
its full theorem text, no nonanticipation conclusion is drawn.

\subsection{The intrinsic mechanism}
In the graph neighbourhood of \(B_R\), with tautological bundle
\(\mathcal U\) and graph map \(T:\mathcal U\to V\), the ideal is
\[
                 (\det T)I_p(\gamma_R\Sym^2T).
\]
Its generators have degree \(d\).  This degree calculation alone
only explains why lower orders are constant.  The inverse requires
three further intrinsic operations.  Multiplication in the finite
scheme recovers the first relation kernel and hence the whole
homogeneous cone.  The rank-one Fano scheme of its reduced determinant
cone then recovers the two tensor rulings.  On the socle Grassmannian
one may distinguish them by projective triviality, already on a
Schubert line.  In general the unequal coefficient-support ranks
orient them without that restriction.  Properness forces the induced
map to \(\PGL(V)\) to be constant.  After choosing its linear lift,
the actual normal-bundle map lies in the commutant of \(\End(V)\),
which recovers the right bundle without a line twist
(Lemma~\ref{lem:actual-factor-descent-v147}).  Finally, cancellation of the intrinsic
determinant ideal recovers the coefficient module.  For pencils its
exterior representation is irreducible, and exterior duality identifies
its single coefficient line with the Pluecker line of \(R\).

Theorem~\ref{thm:structural-coefficient-reconstruction} separates this
geometric mechanism from the particular multiplication algebra.  The
curve theorem shows why the mechanism is not tied to the dimension of
the original parameter Grassmannian.  The geometric orientation used here has a local replacement: coefficient-support asymmetry.  Theorem~\ref{thm:artin-local-inverse-v146} uses it to reconstruct the pencil from one unmarked Artin local algebra.

Theorem~\ref{thm:general-moving-coefficients-v147} gives the broader
principle.  It reconstructs simultaneously varying subspaces in
several Schur modules from the unmarked first-relation thickening,
provided one coefficient support is nonzero and proper.  The
coefficient spaces may have arbitrary prescribed ranks; they need
not be lines coming from quadratic pencils.  Its proof isolates the
actual factor-descent step from classical linear-preserver theory.
Example~\ref{ex:isotrivial-covers-v147} shows why the global descent is
not equivalent to fibrewise reconstruction: two Zariski locally
trivial families have the same Artin fibre and isomorphic nilradical
graded vector bundles in every degree, but their abstract total
schemes differ.  The inverse detects the inequivalent degree-three
coefficient maps.  This is the mechanism behind the moving-pencil
theorem, rather than
an assertion of novelty for canonical grading or for the numerical
value of the first relation degree.


\subsection{An exact image criterion and a pencil-native global consequence}
The coefficient inverse admits an intrinsic converse.
Theorem~\ref{thm:recognition-dichotomy-v148} starts with an
unmarked first relation whose homogeneous radical is a determinant
cone.  After recovering its unordered rulings, invariance under one
full factor group is necessary and sufficient for the residual
relation to come from coefficient subspaces in the Cauchy modules.
Both orderings satisfy this test precisely when every coefficient
support is zero or full.  Thus the strict-support condition is
exactly the criterion for a unique admissible orientation in the
recognized class.  Theorem~\ref{thm:exact-coefficient-stabilizer-v148}
computes its entire affine stabilizer, including the transpose
component and the universal nonlinear automorphism kernel.
Corollary~\ref{cor:intrinsic-orbits-v148} expresses the inverse as
an equivalence of geometric orbit sets, with those ambiguities
specified rather than suppressed.

A global consequence lies inside the quadratic-pencil problem
itself.  Fix a plane \(H=\langle x,y\rangle\subset V\).
For a degree-\(m\) map \(f=[a:b]:\Pj^1\to\Pj(H)\), take
the pencil subbundle
\[
                  \mathcal R_f=(ax+by)H.
\]
Every fibre is equivalent to \(\langle x^2,xy\rangle\).
Nevertheless, Theorem~\ref{thm:covering-pencil-moduli-v148} proves
\[
 Y_f\simeq Y_{f'}\quad\Longleftrightarrow\quad
 f'=\alpha\circ f\circ\sigma,\qquad
                         \alpha,\sigma\in\PGL_2(\C),
\]
for their unmarked order-\(d\) failure thickenings with trivial
source bundle.  Their fibre algebras and every nilradical graded
vector bundle depend only on \(n,m\), not on \(f\).
Their abstract pencil, quotient and first-relation bundles are
fixed as well.  Thus the entire covering equivalence problem, not
only one chosen pair, is retained in the global multiplication and
is invisible to those additive data.  The maps \(z^3\) and
\(z^3+z\) give an explicit separating pair of quadratic pencils
(Corollary~\ref{cor:pencil-native-pair-v148}).

Neither complete reducibility of Cauchy modules, linear preservers
of a Segre cone, nor division by the common factor of a pencil is
new in isolation.  Their synthesis here has a precise output:
recognition of the image and its exact orientation ambiguity after
all matrix markings are forgotten, followed by recovery of a
whole covering map from a locally isotrivial pencil thickening.
A projective factor map alone would not recover the actual source
bundle; the intrinsic normal map and its commutant do so.  A
pointwise pencil inverse would not recover a covering map; one
constant left transformation over the entire proper base does so.
These distinctions specify the geometric content beyond the
classical ingredients.  The recognized coefficient class is not the
class of all determinant-radical relations.  Its geometric orbit
statement alone does not describe nonreduced deformation groupoids;
the relative comparison below uses a further scheme-theoretic inverse.


\subsection{Intrinsic strata, infinitesimal parameters, and pencil moduli}
The relative inverse begins with a geometric condition on the first
relation, not invariance under a factor group.  For a space \(K\) of
homogeneous forms, the divisor-incidence functor defines a canonical
scheme-theoretic image in the relation Grassmannian.  Its exact-gcd open
is the locally closed stratum of Lemma~\ref{lem:universal-gcd}, now
identified by its incidence universal property in
Corollary~\ref{cor:canonical-exact-v151}.  Multiplication by the
universal common-divisor line is an isomorphism onto this open.
Consequently the common factor is recovered even over a nonreduced
parameter base.  The condition is scheme-theoretic: merely asking all
geometric fibres to have the specified common factor is insufficient.
When that divisor is a power of a determinant, the remaining
parameter is an arbitrary gcd-free residual subspace, modulo the
full determinant-preserver group
(Theorem~\ref{thm:determinant-stratum}).  This is a smooth quotient
stack, and neither of the two factor symmetries is required.

For pencils the exact common divisor is \(\delta^{n-1}\), and the
primitive residual relation has degree \(3n-4\)
(Lemma~\ref{lem:primitive-pencil}).  The Cauchy projection and the
Pluecker equations recognize its image as a locally closed scheme,
not just a set of complex points.  The resulting relative theorem is
\begin{equation}\label{eq:intro-effective-moduli}
 \mathscr F_{\rm pen}^{\rm eff}\simeq
                [\Gr(2,\Sym^2V)/\PGL(V)].
\end{equation}
Here the superscript records two explicit operations on morphisms:
first remove tangent-identity nonlinear substitutions, then remove
the ineffective right matrix factor, and fppf stackify.
Theorem~\ref{thm:pencil-stack-equivalence} proves this equivalence
over arbitrary complex parameter schemes, retaining all singular
pencils.  It does not identify the unrigidified algebra stack with
the pencil stack.  Its tangent complex is
\([\mathfrak{pgl}(V)\to\Hom(R,\Sym^2V/R)]\), so the effective
stack is smooth of dimension \(n-3\), with precisely the pencil
stabilizer jumps.  Rank-incidence and spectral Fitting strata retain
their scheme structures and deformation equations by
Corollary~\ref{cor:pencil-specialization-stack}.  This is an exact-image and rigidification comparison; the displayed
rank and spectral consequences are its functorial restrictions, not an
independent classification of those loci.  The two quotient operations
have the universal properties of
Theorem~\ref{thm:rigidification-v151} and
Corollary~\ref{cor:two-rigidifications-v151}.

The ambient determinant-divisor stratum is substantially larger.
Theorem~\ref{thm:transverse-pencil-directions} computes its normal
space to the pencil stratum in three parts: non-Pluecker lines,
non-scalar right-factor directions, and other Cauchy summands.
Theorem~\ref{thm:fixed-support-grassmannians} constructs through every
pencil a projective Grassmannian of relations with the same determinant
radical, containing a projective line whose other points admit neither
factor symmetry.  Thus the family inverse describes how the
coefficient and pencil loci sit in, and can be left within, a larger
intrinsic stratum.  It does not falsely infer deformation stability
from fibrewise recognition.

\subsection{The normalization of a natural common-divisor boundary}
There is a different geometric question beyond the pencil image:
what happens when a divisor of a first relation is no longer uniquely
recoverable in a family?  Let \(E\) have dimension \(e\geq2\), set
\(S_j=\Sym^jE\) and \(s_j=\dim S_j\), and fix
\(g,h\geq1\), \(D=g+h\), and \(2\leq r\leq s_h\).
Let \(Y_g\subset\Gr(r,S_D)\) be the scheme-theoretic image of the
incidence of a line of degree-\(g\) divisors of a rank-\(r\) relation
bundle.  This definition uses the vanishing of the universal quotient
map; it does not first recover a pencil, impose a factor symmetry,
or choose a reduced structure on a set.

\begin{theorem}[The divisorial boundary]
\label{thm:principal-boundary-v151}
The multiplication morphism
\[
 \Pj(S_g)\times\Gr(r,S_h)\longrightarrow Y_g,
                         \qquad ([f],J)\longmapsto fJ
\]
is finite and is the normalization.  The locus \(Y_g\) is integral,
of dimension \(s_g-1+r(s_h-r)\).  At a geometric relation \(K\)
it is normal exactly when \(F_K=\gcd(K)\) has one degree-\(g\)
divisor \([f]\) and \(\gcd(f,F_K/f)=1\).
At the preimage \(([f],K/f)\), if this last gcd has degree \(k\),
the fibre tangent dimension is \(s_k-1\).
\end{theorem}

The theorem is proved in Section~\ref{sec:canonical-boundary-v151}.
Lemma~\ref{lem:fibre-equations-v151} supplies equations for the
entire normalization fibre, not only its tangent space.  For example,
\(K=\langle x^2y,x^2z\rangle\) has a unique linear divisor, but the
normalization fibre is \(\Spec\C[a,b]/(a,b)^2\), of length three.
A different relation has two reduced divisor lifts; yet another has
a gcd-degree jump at a normal point.  The exact criterion distinguishes
these behaviours.  Corollary~\ref{cor:algebra-boundary-v151} carries
this intrinsic boundary to the first-relation algebra stack.  It is
not a classification of all finite-flat smoothings or of all linear
orbits in that boundary.

A direct consequence on the pencil side is proved separately in
Section~\ref{sec:closed-pencil-v151}.  Every quadratic pencil admits
an explicit congruence degeneration to
\(\langle x_1^2,x_1x_2\rangle\); its orbit is the closed flag
variety \(\operatorname{Fl}(1,2;V)\), and it is the unique closed
orbit in the full pencil Grassmannian.  The construction is uniform
for regular and singular pencils and lifts to a finite-flat family
of their failure algebras with exact common factor
\(\delta^{n-1}\).  This proves a particular universal orbit-closure
adjacency by equations, rather than by copying a prescribed invariant
locus through the effective-stack equivalence.  Its elementary
highest-weight mechanism is distinct from the normalization theorem,
which deals with a larger boundary of arbitrary first relations.

\subsection{Framed families and spectral consequences}
Once tensor coordinates are retained, the parameter map is simpler.
The first relation spaces give a closed immersion of
\(\Gr(2,W)\) into a Grassmannian, by
Theorem~\ref{thm:finite-parameter-immersion}.  Its proof factors through
the classical Pluecker embedding, tensoring a line with a fixed space,
and a fixed linear inclusion.  This is the framed encoding consequence
of the coefficient calculation, not a substitute for the unmarked
inverse.  Theorem~\ref{thm:universal-finite-neighbourhood} glues the
actual universal order-\(d\) neighbourhood.  It is finite locally
free over the relative socle of rank
\(\binom{n^2+d}{d}-\binom N2\), and projective and flat over the
pencil parameter scheme.  It commutes with arbitrary complex base
change with that socle retained.  Over a nonreduced base, the specified
relative socle ideal need not be the absolute nilradical.  These marked statements alone do not identify unmarked moduli
stacks or deformation groupoids.  The common-divisor descent and
ineffective-inertia analysis in
Theorem~\ref{thm:pencil-stack-equivalence} supply that separate step.

For regular pencils, the spectral sheaf and its Fitting incidence
schemes are read from the reconstructed pencil.  We retain these
consequences because they exhibit information genuinely lost by
coarser invariants: fixed discriminants and fixed reduced rank loci
can conceal different elementary divisors.  In
Theorem~\ref{thm:fixed-spectral-classification}, every classical
dominance specialization is transported to a flat family of the
finite neighbourhoods, with every strict transition detected at
order \(d\).  The spectral orbit classification itself is classical,
not a second independent Torelli theorem.

\subsection{Prior work and the publication object}
The Weierstrass--Segre classification is used in its root-labelled
form; compare \cite[Theorem~1.1 and Corollary~2.1]{FevolaMandelshtamSturmfels}.
The orthogonal nilpotent input is isolated in
Lemma~\ref{lem:orthogonal-orbits-v145}, with precise references to
\cite[Proposition~1 and Theorem~1]{Ohta1986}, an independent dimension
calculation, and an explicit distinction between full orthogonal
orbits and their special-orthogonal components.

Ballico's relative multiplication construction
\cite[Theorem~0.2; \S1, Remark~1.2 and equation~(6)]{Ballico96}
is a substantive antecedent for studying quadratic-normality failure
on finite linear sections of an embedded surface.  Here the defining
algebra varies, and the inverse starts from an unmarked infinitesimal
failure scheme.  The opening page of the distinct earlier paper \cite[p.~5]{Ballico93}
concerns higher-order embedding properties and restriction maps to
finite subschemes.  Its remaining theorem and proof text has not
been inspected; that first page does not establish the absence of
an inverse statement.  No anticipation or nonanticipation conclusion
relative to that article is asserted.

The Artin-local aspect is compared at theorem level with
\cite{EliasRossiShort,EliasRossiCompressed} in
Section~\ref{sec:first-relations-v147}, including the distinction
between canonical grading of a nonhomogeneous local algebra and
orbit rigidity of our already homogeneous first relation.  The
linear-preserver input is classical \cite{MarcusMoyls}; its precise
hypotheses and a proof in the required case are included in
Lemma~\ref{lem:linear-preserver-v147}.  The algebraic-group
quotient step is supplied by the homomorphism theorem
\cite[Theorem~5.39 and Remark~5.42]{Milne2017}; its use in the
exact stabilizer calculation distinguishes a scheme-theoretic
quotient from a bare bijection of complex points.

The present article consists of intrinsic reconstruction, the canonical
common-divisor boundary, relative rigidification, explicit pencil
specialization, and spectral readout.  The reciprocal-likelihood theorem and projective critical
schemes are preserved in a separate applications manuscript.  Its
finite-flat critical divisor requires supplied split semisimple data
and simple disjoint projective poles; its totally real cover requires
a supplied real definite realization and separated data.  None of
those hypotheses enters the all-pencil theorem here.  Independent
web, exterior-operator, K3, polarized, and boundary results are
preserved in an explicitly non-submitted research archive.  They are
not inputs to any proof in this article and are not part of its
significance claim.  Every proof required here is included here.

% BEGIN PRESERVED PART 35-coefficient-symmetries-v148.tex
\section{Recognition and the exact ambiguity of coefficient systems}
\label{sec:coefficient-symmetries-v148}

The support inequality in the preceding section is not merely a
sufficient test in the class under consideration.  We characterize
that class intrinsically and determine exactly when its orientation
is unique.  We then compute the full linear stabilizer, including
the cases in which an orientation is lost.  These are statements
about unmarked first relations, not about a supplied matrix
presentation.

Throughout this section \(n\ge2\), \(q\ge1\), and \(D=n+q\).
Let \(E\) have dimension \(n^2\), put \(S=\Sym E\), and let
\(0\ne K\subset S_D\).  The affine cone is in \(E^*\).
Suppose that the radical of the homogeneous ideal \((K)\) is a
reduced hypersurface linearly equivalent to the determinant
hypersurface in square \(n\)-by-\(n\) matrices.  We call this the
\emph{determinant-radical condition}.  The ideal \((K)\) is taken
in the full symmetric algebra, not in its Artin truncation.
Its reduced equation gives a line \(\C\delta\subset S_n\),
and unique factorization gives the intrinsic residual subspace
\begin{equation}\label{eq:residual-recognition-v148}
                   J=K/\delta\ \subset S_q.
\end{equation}
Changing the generator \(\delta\) does not change this subspace.
The reduced singular loci of the recovered cone determine its
unordered tensor rulings.  An ordering identifies
\(E=V^*\otimes U\), with \(\dim U=\dim V=n\), up to
separate factor changes.  The subgroup induced by
\(\operatorname{GL}(U)\) on the second factor is independent of
those changes.  An ordering is called \emph{admissible} when
\(J\) is invariant under this full subgroup.  Admissibility is
therefore an intrinsic condition on an ordered ruling, and does
not ask for a chosen basis of either factor.

\begin{theorem}[Intrinsic recognition and orientation dichotomy]
\label{thm:recognition-dichotomy-v148}
Under the determinant-radical condition, an ordering is admissible
if and only if its Cauchy decomposition has
\begin{equation}\label{eq:recognized-cauchy-v148}
 J=\bigoplus_{\lambda\vdash q,\,\ell(\lambda)\le n}
         L_\lambda\otimes\mathbb S_\lambda U,
 \qquad L_\lambda\subset\mathbb S_\lambda V^*.
\end{equation}
The coefficient spaces are unique and are recovered by projection
and contraction.  Suppose at least one ordering is admissible and
put \(r_\lambda=\dim L_\lambda\) and
\(m_\lambda=\dim\mathbb S_\lambda\C^n\).  Then exactly one
of the following alternatives occurs.

\smallskip
\noindent (i) Some \(0<r_\lambda<m_\lambda\).  Precisely one
of the two orderings is admissible, and every linear isomorphism
between such relation spaces preserves that ordering.

\smallskip
\noindent (ii) Every \(r_\lambda\) is either zero or
\(m_\lambda\).  Both orderings are admissible.  The residual
space is a sum of full Cauchy summands, and transposition preserves
\(K\) after identifications of the two matrix factors.

\smallskip
Conversely, every nonzero system in
\eqref{eq:recognized-cauchy-v148} satisfies the determinant-radical
condition after multiplication by \(\delta\).  Thus the theorem
characterizes the image of the coefficient construction, without
assuming a coefficient presentation as part of the input.
\end{theorem}
\begin{proof}
The determinant polynomial is irreducible for \(n\ge2\).
Every element of \(K\) belongs to the radical \((\delta)\),
so \(K\subset\delta S_q\) and the cancellation in
\eqref{eq:residual-recognition-v148} is legitimate.  The intrinsic
recovery of the unordered factors follows from
Lemma~\ref{lem:linear-preserver-v147} and the rank-one rulings.
Changing either factor frame conjugates its full general linear
group into itself; scalar changes act as scalars on \(S_q\).
Consequently the definition of admissibility is independent of
all these choices.

For an ordering, the Cauchy formula is
\[
 S_q=\bigoplus_{\lambda\vdash q,\,\ell(\lambda)\le n}
      \mathbb S_\lambda V^*\otimes\mathbb S_\lambda U.
\]
As representations of \(\operatorname{GL}(U)\), the nonzero
\(\mathbb S_\lambda U\) are irreducible and pairwise
nonisomorphic.  Finite-dimensional rational representations of
this group over \(\C\) are completely reducible.  Hence an
invariant subspace splits across its isotypic components.  In the
\(\lambda\)-component, evaluation identifies that subspace with
\[
 \Hom_{\operatorname{GL}(U)}(\mathbb S_\lambda U,J)
                              \otimes\mathbb S_\lambda U.
\]
Schur's lemma embeds the displayed multiplicity space uniquely in
\(\mathbb S_\lambda V^*\); this is \(L_\lambda\).
This proves necessity in \eqref{eq:recognized-cauchy-v148},
including the absence of unrecorded cross-summand graphs.
Sufficiency follows since the second factor is full.  Projection
onto a Cauchy summand followed by contraction with the dual of its
second factor recovers \(L_\lambda\).

If the reverse ordering is also admissible, \(J\) is invariant
under both full factor groups.  Their product acts irreducibly on
each Cauchy summand, and these product representations are pairwise
nonisomorphic.  Complete reducibility now makes \(J\) a sum of
whole summands.  Equivalently, all \(r_\lambda\) are zero or
full.  Conversely a sum of whole summands is invariant under both
groups and under the factor exchange.  This proves the dichotomy.
One may also see uniqueness directly from the contraction ranks
\((r_\lambda,m_\lambda)\): a factor exchange reverses them,
which is impossible at a nonzero proper support.  Every linear
isomorphism of relation spaces preserves the determinant cone;
Lemma~\ref{lem:linear-preserver-v147} leaves just the two
factor alternatives.  It must therefore preserve the unique
admissible ordering in case (i).

Finally choose a nonzero \(L_\lambda\).  At an invertible
matrix \(T:U\to V\), the map \(\mathbb S_\lambda T\) is
invertible.  A nonzero covector in \(L_\lambda\) consequently
has a nonzero matrix coefficient at \(T\).  Thus the residual
coefficients do not vanish simultaneously at any invertible
matrix.  All elements of \(\delta J\) vanish on the determinant
hypersurface.  Their common zero set is exactly that hypersurface;
the Nullstellensatz and its reduced irreducible equation give
\(\rad(\delta J)=(\delta)\), as asserted.
\end{proof}

The complete reducibility and Cauchy statements here are classical
representation theory; see \cite{ChengWang}.  Their role is to
supply an intrinsic image criterion \emph{after} the matrix factors
have been recovered.  For a relation space that fails the
admissibility criterion, the theorem does not assert a one-sided
coefficient description.  In particular it is not a classification
of arbitrary homogeneous ideals or arbitrary degeneracy schemes.

\subsection{Stabilizers and the full transpose ambiguity}
Use the normal action \(T\mapsto gTh^{-1}\).  Its action on
polynomial functions is inverse pullback, so it gives an honest
representation \(\rho\) on \(E=V^*\otimes U\).
The kernel of \(\rho:\operatorname{GL}(V)\times
\operatorname{GL}(U)\to\operatorname{GL}(E)\) is the diagonal
scalar subgroup.  For a recognized system define
\[
 G_L=\{g\in\operatorname{GL}(V):
          (\mathbb S_\lambda g^{-1})^*L_\lambda=L_\lambda
                                      \text{ for every }\lambda\},
 \qquad H_K=\operatorname{Stab}_{\operatorname{GL}(E)}(K).
\]
These are closed subgroup schemes, with the stabilizer scheme
structures from their actions on Grassmannians.

\begin{theorem}[Exact coefficient stabilizer]
\label{thm:exact-coefficient-stabilizer-v148}
In case (i) of Theorem~\ref{thm:recognition-dichotomy-v148},
\begin{equation}\label{eq:proper-stabilizer-v148}
                 H_K\simeq
        (G_L\times\operatorname{GL}(U))/\mathbb G_m.
\end{equation}
In case (ii), after choosing bases for the two factors,
\begin{equation}\label{eq:full-stabilizer-v148}
 H_K\simeq
 \bigl((\operatorname{GL}_n\times\operatorname{GL}_n)
                         /\mathbb G_m\bigr)\rtimes C_2,
\end{equation}
where \(C_2\) is generated by transposition.  These are
isomorphisms of affine algebraic groups over \(\C\), and hence
of their represented functors on commutative complex algebras.

For the ungraded algebra
\(A(K)=S/((K)+\mathfrak m^{D+1})\), they give the split exact
sequence
\begin{equation}\label{eq:complete-local-aut-v148}
 1\longrightarrow U_K\longrightarrow\operatorname{Aut}(A(K))
                      \longrightarrow H_K\longrightarrow1.
\end{equation}
Here \(U_K\) is the universal higher-order substitution group
from Proposition~\ref{prop:first-relation-group-v147}, of
dimension \(n^2(\length A(K)-1-n^2)\).
\end{theorem}
\begin{proof}
Every element of \(H_K(\C)\) preserves the radical of \((K)\),
and therefore the determinant cone.  In case (i), the unique
admissible orientation excludes its transposed component.  The
remaining map is \(\rho(g,h)\).  Cancellation of the determinant
line and projection onto each Cauchy summand show that it preserves
\(J\) if and only if \(g\in G_L\); the arbitrary right map
acts on the full factor and introduces no additional condition.
The kernel is the diagonal \(\mathbb G_m\).  In case (ii),
every separate factor transformation and transposition preserves
each active full summand.  These exhaust the determinant preservers,
so they exhaust \(H_K(\C)\).  The transposition has order two
and lies outside the preserving component for \(n\ge2\), giving
the asserted semidirect product.  Its action on the left-right
quotient is conjugation by the matrix-transpose map; no particular
formula for its action on chosen factor lifts is needed.

We justify the scheme statement as well.  The homomorphism theorem
for algebraic groups identifies the quotient by the kernel with
the closed image \cite[Theorem~5.39 and Remark~5.42]{Milne2017}.
Every finite-type group scheme over a characteristic-zero field is
smooth \cite[Lemma~39.8.2]{StacksGroups}.  Thus the stabilizer and
the image are reduced closed subgroup schemes.  The equalities on
complex points just proved imply equality of their defining radical
ideals.  This gives \eqref{eq:proper-stabilizer-v148} and
\eqref{eq:full-stabilizer-v148} as algebraic-group isomorphisms,
not only as bijections on complex points.  Finally
Proposition~\ref{prop:first-relation-group-v147} applies in degree
\(D\) and gives \eqref{eq:complete-local-aut-v148} with the
stated kernel, splitting and dimension.
\end{proof}

For completeness, the omitted zero system has \(K=0\).
Then \(A(K)\) is the free truncated symmetric algebra,
\(H_K=\operatorname{GL}(E)\), and its linear automorphisms
need not preserve any matrix factors.  The determinant-radical
condition deliberately excludes this third possibility.  A full
nonzero coefficient system and a zero system therefore lose
information in different ways: the former retains unordered
matrix geometry and has exactly a transpose ambiguity, while the
latter retains no intrinsic matrix geometry at all.

\subsection{The zero/full case over a projective reduction}
The local transpose ambiguity also has an exact global form.  Fix
a nonempty set of active partitions and take the corresponding
\(\mathcal L_\lambda\) to be full, with all other coefficient
bundles zero.  Use the same construction as in
\eqref{eq:coefficient-thickening-v147}, without its strict-support
hypothesis.

\begin{proposition}[Global ambiguity for full coefficient systems]
\label{prop:full-global-ambiguity-v148}
Let \(Y_i\) be two such schemes, for the same \(n,q\) and
active partitions, over smooth connected projective bases \(B_i\)
with source bundles \(\mathcal U_i\).  They are isomorphic as
abstract complex schemes if and only if there are an isomorphism
\(\sigma:B_1\to B_2\) and an actual bundle isomorphism
\(\mathcal U_1\simeq\sigma^*\mathcal U_2\).
An isomorphism exchanging the two ruling families over a fixed
\(\sigma\) exists precisely when, in addition, the common source
bundle is projectively trivial.  Equivalently,
\[
 \mathcal U_1\simeq\mathcal M^{\oplus n},\qquad
 \sigma^*\mathcal U_2\simeq\mathcal M^{\oplus n}
\]
for a line bundle \(\mathcal M\) on \(B_1\).
Thus even the full systems recover the actual source-bundle class;
the extra ambiguity is an exchange, not an undetected line twist.
\end{proposition}
\begin{proof}
An abstract isomorphism recovers the reduction, the actual normal
bundle map, the first relation, and the unordered determinant
rulings by the proof of
Theorem~\ref{thm:general-moving-coefficients-v147}, which uses
strict support only to orient the rulings.  Over the connected
base the choice to preserve or exchange the two components is
constant.  In the preserving case, projectivity makes the left
projective map constant.  A constant lift and
Lemma~\ref{lem:actual-factor-descent-v147} give an actual
isomorphism of the source bundles.

In the exchanging case each source projective bundle is identified
with the constant left projective bundle on the other side.  It
is therefore trivial.  A vector bundle with a trivialized
projectivization is a line twist of a constant vector space:
locally lift the given projective trivialization to linear frames;
on intersections their differences are scalar matrices, whose
scalar cocycle is the required line bundle.  After pullback by
\(\sigma\), write
\(\mathcal U_i^*=\mathcal M_i^*\otimes W_i\), with \(W_i\)
constant of dimension \(n\).  The normal bundles become
\(\mathcal M_i^*\otimes(V_i\otimes W_i)\).
Their actual isomorphism induces a morphism from \(B_1\) to
the projective-linear isomorphisms between two constant
\(n^2\)-spaces.  This target is affine; the morphism is constant.
Choose its constant linear lift.  On removing it, the actual
bundle map is scalar on each projective fibre.  In local line
frames its off-diagonal entries vanish and all diagonal entries
agree.  Reducedness makes these equalities hold as bundle maps.
Their common entry is a nowhere vanishing section of
\(\mathcal Hom(\mathcal M_1^*,\mathcal M_2^*)\), and hence
\(\mathcal M_1\simeq\mathcal M_2\).  This proves necessity
of the common actual line twist, and thus of the source-bundle
isomorphism.

Conversely any actual source-bundle isomorphism, together with a
constant left identification, preserves all active full Cauchy
summands, so it gives an isomorphism of the schemes.  If the
source bundles have the common line-twist presentation above,
choose identifications of the constant factor spaces and
transpose their tensor factors, leaving the line factor intact.
This is a global actual normal map.  It preserves the determinant
line, every active full summand, and the truncation, and supplies
an exchanging isomorphism.  These constructions prove all the
claims.
\end{proof}

\subsection{The intrinsic orbit space}
Fix a rank vector \(\boldsymbol r\) with at least one proper
nonzero entry and form the product of Grassmannians
\[
 X_{\boldsymbol r}=
   \prod_{\lambda\vdash q,\,\ell(\lambda)\le n}
       \Gr(r_\lambda,\mathbb S_\lambda V^*).
\]
Zero and full factors in this product are points.  Let
\(\mathcal Z_{\boldsymbol r}\) denote the set of subspaces of
\(\Sym^D\C^{n^2}\) satisfying the determinant-radical condition
and the intrinsic admissibility test, with ranks
\(\boldsymbol r\) in the unique admissible orientation.

\begin{corollary}[Geometric orbit classification]
\label{cor:intrinsic-orbits-v148}
The coefficient construction induces bijections
\begin{equation}\label{eq:intrinsic-orbits-v148}
 X_{\boldsymbol r}(\C)/\PGL(V)
  \ \simeq\ \mathcal Z_{\boldsymbol r}/\operatorname{GL}_{n^2}(\C)
  \ \simeq\ \{A(K):K\in\mathcal Z_{\boldsymbol r}\}/\simeq.
\end{equation}
Their stabilizers, including the ineffective right factor and the
nonlinear kernel of the last passage, are precisely those in
Theorem~\ref{thm:exact-coefficient-stabilizer-v148}.
For a nonempty zero/full rank pattern there is just one relation
orbit for that pattern, with the stabilizer in
\eqref{eq:full-stabilizer-v148}.
\end{corollary}
\begin{proof}
The recognition theorem gives surjectivity of the first map.
A linear isomorphism of two such relations is orientation-preserving;
its left factor carries every \(L_\lambda\) to the corresponding
coefficient space, and its right factor is unrestricted.  Scalars
act trivially on the Grassmannians.  This proves injectivity and
the first bijection.  The first-relation orbit proposition gives
the second.  For a zero/full pattern each coefficient Grassmannian
is a point, and the full-stabilizer theorem supplies the last claim.
\end{proof}

Equation~\eqref{eq:intrinsic-orbits-v148} is a statement about
geometric orbit sets with the explicitly computed affine
stabilizers.  It does not identify nonreduced moduli stacks or
assert that every infinitesimal deformation remains in this class.
For smooth projective reduced bases, its uniquely oriented case
combines with actual factor descent in
Theorem~\ref{thm:general-moving-coefficients-v147}.  The following
pencil-native application shows that this global step separates an
entire covering problem that fibrewise classification cannot see.

% BEGIN PRESERVED PART 29-curve-reconstruction-v145.tex
\section{Reconstruction on a projective curve}
\label{sec:curve-recognition-v145}

The global part of the inverse argument has a more economical geometric
realization.  It uses a proper base to make the projective left frame
constant, and a nontrivial right projective bundle to distinguish the
rulings.  Neither step requires the whole socle Grassmannian.  We now
replace that Grassmannian by a curve.  In particular, a neighbourhood
supported on one Schubert line already retains every pencil.

Let \(C\) be a smooth connected projective complex curve, and let
\(\mathcal U\) be a rank-\(n\) bundle on \(C\), where \(n\ge3\).
Assume that \(\Pj_C(\mathcal U^*)\) is not projectively trivial.
A sufficient numerical condition is
\begin{equation}\label{eq:curve-degree-obstruction-v145}
                 \deg\det\mathcal U\not\equiv0\pmod n.
\end{equation}
Indeed a projectively trivial rank-\(n\) bundle is a line-bundle twist
of a trivial bundle, so its determinant degree is divisible by \(n\).
Fix \(V\) of dimension \(n\), put \(E=\Hom(\mathcal U,V)\),
and let \(\mathfrak a\) be the zero-section ideal in
\(\Sym E^*\).  For a pencil \(R\subset\Sym^2V\), the quotient
\(\gamma_R:\Sym^2V\to S_R\) defines the homogeneous ideal
\begin{equation}\label{eq:curve-failure-ideal-v145}
 \mathcal I_{R,C}=(\det T)I_p(\gamma_R\Sym^2T),
 \qquad p=\binom{n+1}{2}-2,
\end{equation}
where \(T:\mathcal U\to V\) is the universal map on \(\Tot E\).
The expression is independent of local frames: both determinant
multiplication and the maximal-minor construction are bundle
constructions.  Define
\[
 Y_{R,C}^{[k]}=\Spec_C\bigl(\Sym E^*/
             (\mathcal I_{R,C}+\mathfrak a^{k+1})\bigr).
\]
Only its abstract complex scheme is retained, not its projection to
\(C\), its embedding in \(\Tot E\), or either tensor factor.

\begin{theorem}[Curve-supported reconstruction]
\label{thm:curve-reconstruction-v145}
For fixed \((C,\mathcal U,V)\) as above and
\(d=n^2+2n-4\), every pair of quadratic pencils satisfies
\begin{equation}\label{eq:curve-torelli-v145}
 Y_{R,C}^{[d]}\simeq Y_{R',C}^{[d]}
       \quad\Longleftrightarrow\quad R'=gR
              \quad\text{for some }g\in\PGL(V).
\end{equation}
All lower-order neighbourhoods are independent of \(R\).
The order-\(d\) schemes are finite locally free over \(C\) of rank
\(\binom{n^2+d}{d}-\binom{N}{2}\), where \(N=\binom{n+1}{2}\).
No regularity condition is imposed on the pencil.
\end{theorem}
\begin{proof}
On any frame chart, the ideal in
\eqref{eq:curve-failure-ideal-v145} is generated in degree \(d\).
Its degree-\(d\) piece \(K_R\) is a subbundle of rank
\(\binom N2\), by the fixed coefficient inclusion
\eqref{eq:finite-coefficient-factorization}; its quotient is locally
free.  The quotient algebra has its full symmetric powers in degrees
less than \(d\), has \(\Sym^dE^*/K_R\) in degree \(d\),
and is zero above that degree.  This proves the rank and the assertion
about lower orders.

Its reduction is \(C\), and its nilradical \(\mathfrak n\)
recovers \(L=\mathfrak n/\mathfrak n^2\) and
\[
       K_R=\ker[\Sym^dL\longrightarrow\mathfrak n^d].
\]
Thus the abstract finite-order scheme reconstructs the entire graded
algebra \(\Sym L/(K_R)\), by the same argument as
Lemma~\ref{lem:first-relation-cone}.  The ideal here means the
homogeneous ideal generated by the recovered kernel, not its radical.
No splitting of the reduction or supplied projection is used.

The reduction of this reconstructed cone is the determinant cone.
Indeed its ideal has the determinant factor, and when \(T\) is
invertible the quadratic quotient has full rank.  Its rank-one Fano
scheme has two ruling components,
\(\Pj(V)\times C\) and \(\Pj_C(\mathcal U^*)\).
The hypothesis distinguishes the first one intrinsically.  An abstract
isomorphism of neighbourhoods may induce a nonidentity automorphism
of \(C\); pull back the second bundle along that automorphism.
The oriented ruling map is nevertheless a morphism from \(C\) to
\(\PGL_n\).  It is constant, since \(C\) is connected projective
and \(H^0(C,\OO_C)=\C\).  Lemma~\ref{lem:oriented-segre-descent}
therefore supplies one constant projective left frame.

Cancel the intrinsic determinant ideal and apply the coefficient
construction of Theorem~\ref{thm:structural-coefficient-reconstruction}.
For pencils, Lemma~\ref{lem:pencil-irreducibility} leaves just the
single line \(\C\beta_R\).  Exterior duality
\eqref{eq:pencil-exterior-duality} recovers \(R\) with respect to
that same projective frame.  This proves the forward implication in
\eqref{eq:curve-torelli-v145}.  Conversely a linear lift of \(g\)
induces the bundle map on \(E\) and the quotient isomorphism on
\(S_R\), and carries all the displayed ideals and their truncations
to the corresponding ideals for \(R'\).  This proves the converse.
\end{proof}

\begin{corollary}[One Schubert line suffices]
\label{cor:schubert-line-reconstruction-v145}
Let \(\ell\simeq\Pj^1\) be any Schubert line in
\(B_R=\Gr(n,S_R)\), and use the canonical graph projection
\(Z_R^{[d]}\to B_R\) to form its restriction \(Z_{R,\ell}^{[d]}\).
Its abstract isomorphism class is independent of the choice of
\(\ell\), and
\[
 Z_{R,\ell}^{[d]}\simeq Z_{R',\ell'}^{[d]}
       \quad\Longleftrightarrow\quad R'=gR
                       \text{ for some }g\in\PGL(V).
\]
The smallest uniform reconstruction order for these curve-supported
neighbourhoods is again \(d\).  Their reductions all have dimension
one, instead of \(n(p-n)\).
\end{corollary}
\begin{proof}
The tautological bundle on a Schubert line is
\(\OO_{\Pj^1}(-1)\oplus\OO_{\Pj^1}^{n-1}\).
Pulling back the homogeneous ideal of the graph neighbourhood gives
exactly \eqref{eq:curve-failure-ideal-v145} with that bundle.  Its
isomorphism type depends on this bundle and \(R\), not on the
chosen inclusion of the line in the socle.  The determinant degree
is \(-1\), so Theorem~\ref{thm:curve-reconstruction-v145} applies.
Its lower-order schemes are all independent of the pencil, while
\(\langle x_1^2,x_1x_2\rangle\) and
\(\langle x_1^2,x_2^2\rangle\) are inequivalent.  This proves
the uniform minimality assertion.
\end{proof}

This restriction is made in the construction; it is not a claim that
an abstract neighbourhood canonically chooses a Schubert line.  No
line is supplied to an isomorphism in the corollary: the reduced curve
is recovered as the reduction of its source.  Nor is the corollary a
statement about one unmarked Artinian normal slice.  The projective
curve and the nontrivial ruling remain essential to this argument.
The same theorem applies, for example, to
\(\mathcal U=\OO_C(-P)\oplus\OO_C^{n-1}\) on every smooth
projective connected curve with a point \(P\).  It isolates the
properness and ruling obstruction used in reconstruction, independently
of the original Grassmannian realization.

% BEGIN PRESERVED PART 31-moving-pencils-v146.tex
\section{Reconstruction of moving pencils and their isomorphisms}
\label{sec:moving-pencils-v146}

The inverse need not be applied separately to constant pencils.  A single
unmarked finite scheme can retain a pencil varying over its entire
projective reduction.  We prove this assertion, including the recovery
of the normal bundle and the ambiguity in lifting the recovered data to
an isomorphism.  The coefficient lines in this section are subbundles;
constancy of the pencil is not assumed.

\subsection{The family and its first relation bundle}
Fix \(n\ge3\), and write
\[
 N=\binom{n+1}{2},\qquad p=N-2,\qquad
 d=n+2p=n^2+2n-4,\qquad M=\binom N2.
\]
An admissible datum consists of a smooth connected projective complex
variety \(B\), a rank-\(n\) vector bundle \(\mathcal U\) on \(B\), a
constant \(n\)-dimensional vector space \(V\), and a rank-two locally
direct-summand subbundle
\[
 \mathcal R\ \subset\ \Sym^2V\otimes\OO_B.
\]
We require that \(\Pj_B(\mathcal U^*)\) is not a trivial projective
bundle.  Put \(\mathcal S=(\Sym^2V\otimes\OO_B)/\mathcal R\) and
\(\mathcal E=\Hom(\mathcal U,V\otimes\OO_B)\).  On its total space,
let \(T\) be the universal map and let \(\gamma\) denote the quotient
onto \(\mathcal S\).  Define the homogeneous ideal and its truncations by
\begin{equation}\label{eq:moving-construction-v146}
 \begin{split}
 \mathcal I_{\mathcal R}&=(\det T)I_p(\gamma\Sym^2T),\\
 \mathcal A_{\mathcal R}^{[k]}
   &=\Sym(\mathcal E^*)/
       (\mathcal I_{\mathcal R}+\mathfrak a^{k+1}),\qquad
 Y_{\mathcal R}^{[k]}=\Spec_B\mathcal A_{\mathcal R}^{[k]},
 \end{split}
\end{equation}
where \(\mathfrak a\) is the positive-degree ideal.  Determinant-line
changes of frame multiply local generators by units, so this is an
ideal sheaf, independent of all frames.  This construction is the
pullback of the coefficient construction in
Theorem~\ref{thm:universal-finite-neighbourhood}, with an arbitrary
rank-\(n\) source bundle.  It is not necessary to choose an embedding
of \(B\) in a socle Grassmannian.

\begin{lemma}[A varying coefficient line]\label{lem:moving-coefficients-v146}
The ideal \(\mathcal I_{\mathcal R}\) has no component of degree less
than \(d\) and is generated by its degree-\(d\) component
\(\mathcal K_{\mathcal R}\).  This component is a subbundle of
\(\Sym^d\mathcal E^*\) of rank \(M\).  After cancellation of the
determinant line, its degree-\(2p\) coefficient module determines
\(\mathcal R\) as a subbundle, once the two tensor factors of
\(\mathcal E\) are supplied.  In particular,
\(Y_{\mathcal R}^{[d]}\to B\) is finite locally free of rank
\[
                   \binom{n^2+d}{d}-M.
\]
\end{lemma}
\begin{proof}
Put \(F(H)=\bigwedge^p\Sym^2H\) for rank-\(n\) vector spaces or
bundles \(H\).  This is an irreducible polynomial representation by
Lemma~\ref{lem:pencil-irreducibility}; its dimension is \(M\).  Write
\(\mathcal L_{\mathcal R}=\bigwedge^p\mathcal S^*\), viewed, by
\(\bigwedge^p\gamma^*\), as a line subbundle of
\(F(V)^*\otimes\OO_B\).  The residual minors are precisely the image of
\begin{equation}\label{eq:moving-coefficient-map-v146}
 \mathcal L_{\mathcal R}\otimes F(\mathcal U)
  \longrightarrow \Sym^{2p}(V^*\otimes\mathcal U),\qquad
 \beta\otimes\xi\longmapsto[T\mapsto\beta(F(T)\xi)].
\end{equation}
The full map \(F(V)^*\otimes F(U)\to
\Sym^{2p}(V^*\otimes U)\) is the nonzero inclusion of its unique
Cauchy summand: evaluation at \(T=1\) is nonzero, as in
\eqref{eq:general-coefficient-map}.  Restricting to a nonzero line on
the left therefore has rank \(M\) in every fibre.  The coefficients
vary algebraically.  The constant-rank criterion makes the image a
subbundle; equivalently one may choose a nonvanishing maximal minor
of this coefficient map locally on the Grassmannian of pencils and
pull it back to \(B\).  Multiplication by the determinant, a nonzero
polynomial in every fibre, again has constant rank \(M\).

All defining generators have degree \(d\), proving the assertion
about the ideal.  For the recovery assertion, the residual image,
inside the indicated Cauchy summand, is
\(\mathcal L_{\mathcal R}\otimes F(\mathcal U)\).  Contraction with
\(F(\mathcal U)^*\) has image exactly \(\mathcal L_{\mathcal R}\)
as a line subbundle of \(F(V)^*\otimes\OO_B\).  This statement is an
equality of images of bundle maps, not just equality at closed points.
Exterior duality identifies it with the Pluecker line
\(\bigwedge^2\mathcal R\), up to the fixed determinant line of
\(\Sym^2V\).  The Grassmannian Pluecker embedding consequently
recovers the morphism \(B\to\Gr(2,\Sym^2V)\), and its pullback
of the tautological bundle recovers \(\mathcal R\).
Finally, the truncated algebra is the sum of its free graded pieces
below \(d\) and the locally free degree-\(d\) quotient by
\(\mathcal K_{\mathcal R}\).  Its rank is the asserted binomial sum
minus \(M\).
\end{proof}

\subsection{The unmarked inverse}
\begin{theorem}[Reconstruction of a moving pencil]
\label{thm:moving-reconstruction-v146}
Let \((B_i,\mathcal U_i,V_i,\mathcal R_i)\), \(i=1,2\), be admissible
data of the same rank \(n\).  Then
\begin{equation}\label{eq:moving-inverse-v146}
 Y_{\mathcal R_1}^{[d]}\simeq Y_{\mathcal R_2}^{[d]}
\end{equation}
as abstract complex schemes if and only if there exist an isomorphism
\(\sigma:B_1\to B_2\), a constant linear isomorphism
\(g:V_1\to V_2\), and a bundle isomorphism
\(h:\mathcal U_1\to\sigma^*\mathcal U_2\) such that
\begin{equation}\label{eq:moving-data-v146}
                  \sigma^*\mathcal R_2=(\Sym^2g)\mathcal R_1.
\end{equation}
The reduced base, normal tensor orientation, isomorphism class of the
actual source bundle, and the varying pencil are recovered together.
Neither the projection onto the base nor the ambient bundle is
supplied in \eqref{eq:moving-inverse-v146}.  The constant left map
cannot be replaced by a separately chosen change of coordinates in
each fibre.  A pair \((g,h)\) defining the same normal-bundle map is
unique up to simultaneous multiplication by a nonzero constant.
\end{theorem}
\begin{proof}
Let \(\mathfrak n_i\) be the nilradical of
\(\OO_{Y_{\mathcal R_i}^{[d]}}\).  Since the positive-degree ideal
is nilpotent and \(B_i\) is reduced, the reduction is exactly \(B_i\),
and this ideal is exactly \(\mathfrak n_i\).  There are no relations
below degree \(d\), so multiplication intrinsically recovers
\begin{equation}\label{eq:moving-intrinsic-kernel-v146}
 \mathcal E_i^*=\mathfrak n_i/\mathfrak n_i^2,
 \qquad
 \mathcal K_i=\ker\left[
 \Sym^d(\mathfrak n_i/\mathfrak n_i^2)
       \longrightarrow\mathfrak n_i^d\right].
\end{equation}
The map is well defined since \(\mathfrak n_i^{d+1}=0\).
An abstract isomorphism gives \(\sigma\), an isomorphism of these
normal bundles, and the matching relation subbundles.  In particular
it gives an isomorphism of the untruncated homogeneous cones
\[
               \Spec_{B_i}\bigl(\Sym\mathcal E_i^*/(\mathcal K_i)\bigr).
\]
This cone is reconstructed from the first relation module; it is not
the associated graded algebra of the finite scheme, whose degrees
above \(d\) vanish.

The reduction of this cone is the determinant cone.  Indeed the
residual quadratic map is surjective wherever \(T\) is invertible,
whereas the determinant factor annihilates the ideal wherever
\(T\) is singular.  On a frame chart the determinant hypersurface
is reduced over the smooth base, so the radical ideal is precisely
the determinant ideal.  Successive relative reduced singular loci
give the lower matrix-rank cones.  The rank-one cone has the two
ruling parameter schemes
\[
             \Pj(V_i)\times B_i,\qquad \Pj_{B_i}(\mathcal U_i^*),
\]
by Lemma~\ref{lem:rank-one-fano-scheme}.  Only the first is
projectively trivial.  Thus an isomorphism cannot interchange them.
The induced map on the trivial ruling is a morphism
\(B_1\to\operatorname{Isom}(\Pj(V_1),\Pj(V_2))\).
After choosing one identification this target is the affine variety
\(\PGL_n\).  Since \(B_1\) is connected and projective, every such
morphism is constant.  Choose one constant linear lift \(g\).

We check that no line-bundle ambiguity remains in the other factor.
Locally trivialize both source bundles.  The linear rank-one
preserver with its left projective map fixed to \([g]\) is of the
form
\begin{equation}\label{eq:moving-normal-map-v146}
                           T\longmapsto gTh^{-1}
\end{equation}
by Lemma~\ref{lem:oriented-segre-descent}.  Once this particular
\(g\) is fixed, the local linear map \(h\) is unique: if two right
maps give the same value for every \(T\), they are equal.  The local
maps therefore agree on overlaps and glue to an actual bundle
isomorphism \(\mathcal U_1\to\sigma^*\mathcal U_2\).  Alternatively,
after removing \(g\) the normal-bundle map commutes with
\(\End(V_1)\); the commutant is the right-factor homomorphism bundle.
Rescaling the constant lift \(g\) rescales \(h\) by the same constant.
This proves both existence and the asserted uniqueness for a given
normal-bundle map.

Cancel the intrinsic determinant ideal in the recovered homogeneous
ideal.  On a frame chart the determinant polynomial is a
non-zero-divisor and
\(( (\det T)\mathcal J:(\det T))=\mathcal J\); the identities glue
without selecting a global determinant generator.  Apply
Lemma~\ref{lem:moving-coefficients-v146} to the degree-\(2p\) part.
The right map \(h\) acts on the full factor \(F(\mathcal U)\), so
contraction against its dual removes that change of frame.  The
remaining line subbundle transforms by the one constant left map
\(g\).  Exterior duality yields exactly
\eqref{eq:moving-data-v146}.  This uses the subbundle version proved
above, not the constant-coefficient hypothesis in
Theorem~\ref{thm:structural-coefficient-reconstruction}.

Conversely \(\sigma,g,h\) satisfying \eqref{eq:moving-data-v146}
identify the quotient bundles \(\mathcal S_i\), the quadratic maps,
and their maximal-minor ideals.  The determinant changes by a unit
in local frames.  The map \eqref{eq:moving-normal-map-v146} consequently
identifies both the homogeneous ideals and every truncation,
proving \eqref{eq:moving-inverse-v146}.
\end{proof}

\begin{corollary}[Uniform sharpness over an admissible base]
\label{cor:moving-sharpness-v146}
For fixed \(B,\mathcal U,V\), all schemes
\(Y_{\mathcal R}^{[k]}\) with \(k<d\) are independent of
\(\mathcal R\).  Order \(d\) is the smallest uniform order at which
the moving pencil is recoverable in the sense of
Theorem~\ref{thm:moving-reconstruction-v146}.
\end{corollary}
\begin{proof}
Below \(d\) the ideal \(\mathcal I_{\mathcal R}\) contributes no
relations, and the algebra is
\(\bigoplus_{j=0}^k\Sym^j\mathcal E^*\).  Constant pencils are
among the admissible subbundles.  The pencils
\(\langle x_1^2,x_1x_2\rangle\) and
\(\langle x_1^2,x_2^2\rangle\) are inequivalent: all members of
the first have a common linear factor and those of the second do
not.  A base automorphism does not change either constant pencil.
The theorem separates their order-\(d\) schemes, whereas every
lower truncation agrees.  This proves uniform, not pairwise, minimality.
\end{proof}

\subsection{The part of an isomorphism that the inverse does not record}
The recovered data classify isomorphism classes, but an isomorphism
of finite schemes need not itself be linear in the nilpotent
directions.  The following statement identifies this distinction
without replacing it by a claim about moduli stacks.

Fix an admissible datum, put \(Y=Y_{\mathcal R}^{[d]}\), and let
\(\mathscr L_{\mathcal R}\) be the group of triples
\((\sigma,g,h)\) satisfying \eqref{eq:moving-data-v146}, modulo
\((g,h)\sim(cg,ch)\) for \(c\in\C^*\).  Its group law is
composition of their induced normal-bundle maps.  Write
\(\operatorname{Aut}_{\C}(Y)\) for the group of complex scheme
automorphisms.  For \(j\ge1\), define \(G_j\) to consist of the
automorphisms inducing the identity on the reduction and satisfying
\begin{equation}\label{eq:automorphism-filtration-v146}
 (\alpha^*-1)(\mathfrak n^k)\subseteq\mathfrak n^{k+j}
 \quad\text{for every }k\ge0,\qquad \mathfrak n^0=\OO_Y.
\end{equation}
They act on sheaves on the same underlying topological space.

\begin{theorem}[Linear data and the nilpotent kernel]
\label{thm:automorphism-kernel-v146}
There is an exact sequence of groups
\begin{equation}\label{eq:automorphism-sequence-v146}
 1\longrightarrow G_1\longrightarrow
 \operatorname{Aut}_{\C}(Y)\longrightarrow
 \mathscr L_{\mathcal R}\longrightarrow1.
\end{equation}
The displayed homogeneous presentation of \(Y\) supplies a splitting.
Moreover \(G_{d+1}=1\),
\([G_i,G_j]\subseteq G_{i+j}\), and there are natural injections
\begin{equation}\label{eq:automorphism-symbol-v146}
 G_j/G_{j+1}\ \lhook\joinrel\longrightarrow\
 H^0\!\left(B,\operatorname{Der}_{\C}
       (\gr_{\mathfrak n}\OO_Y)_j\right)
\end{equation}
into additive groups of degree-\(j\) derivations.  Thus the kernel is
nilpotent and records precisely the automorphisms invisible on the
associated graded scheme.  No surjectivity onto the derivation spaces
is asserted.
\end{theorem}
\begin{proof}
An automorphism induces a base map and a degree-one bundle map.  The
proof of Theorem~\ref{thm:moving-reconstruction-v146} identifies these
with a unique class of triples in \(\mathscr L_{\mathcal R}\).
The graded algebra is generated in degree one over degree zero, with
relation subbundle \(\mathcal K_{\mathcal R}\).  Hence its entire
graded automorphism is already determined by that triple.  The kernel
is the identity on every graded piece, which is exactly condition
\eqref{eq:automorphism-filtration-v146} for \(j=1\).
Conversely the homogeneous maps \(T\mapsto gTh^{-1}\) give actual
automorphisms of \(Y\), respect composition, and induce the prescribed
triples.  This proves exactness and the splitting.  A choice of the
homogeneous presentation supplies that section; the assertion does
not endow an arbitrary unmarked presentation with a preferred section.

For \(\alpha\in G_j\), put \(D_\alpha=\alpha^*-1\).  It raises
filtration by \(j\), and is therefore nilpotent as an additive
operator.  The inverse of \(1+D_\alpha\) is the finite geometric
series.  Composition of filtration-raising operators shows that a
commutator of elements of \(G_i,G_j\) raises filtration by at least
\(i+j\).  Also \(G_{d+1}=1\) follows by taking \(k=0\), since
\(\mathfrak n^{d+1}=0\).

The leading term of \(D_\alpha\) defines maps
\(\mathfrak n^k/\mathfrak n^{k+1}\to
\mathfrak n^{k+j}/\mathfrak n^{k+j+1}\).
They are independent of lifts.  The identity
\[
 D_\alpha(ab)=aD_\alpha(b)+D_\alpha(a)b+
                           D_\alpha(a)D_\alpha(b)
\]
reduces to the Leibniz identity in the indicated graded degree, since
\(2j\ge j+1\).  Thus the leading term is a sheaf derivation over
\(\C\), including its action on degree zero; it need not be an
\(\OO_B\)-linear derivation.  Leading terms add under composition,
and their kernel is exactly \(G_{j+1}\).  They glue globally,
proving \eqref{eq:automorphism-symbol-v146} and all assertions.
\end{proof}

\begin{example}[A nontrivial invisible automorphism]
\label{ex:invisible-shear-v146}
Let \(B=\Pj^1\) and
\(\mathcal U=\OO(-1)\oplus\OO^{n-1}\).
The bundle \(\mathcal E^*=V^*\otimes\mathcal U\) has at least two
trivial direct summands.  Choose their degree-one generators \(x,y\).
The substitution
\[
                          x\longmapsto x+y^2
\]
fixing \(y\), the other summands and \(\OO_B\), defines a nonidentity
automorphism of \(Y_{\mathcal R}^{[d]}\) in \(G_1\), for every
admissible \(\mathcal R\).  Indeed each relation has degree \(d\);
its change under this substitution has degree at least \(d+1\) and
vanishes in the truncation.  The inverse is \(x\mapsto x-y^2\).
The automorphism is nonidentity because there are no degree-two
relations.  Therefore identifying all finite-scheme automorphisms
with the stabilizer of the pencil would be false, even in this simple
case.  The inverse theorem and the exact sequence retain, rather than
suppress, this higher-order freedom.
\end{example}

\subsection{A nonconstant family across spectral collisions}
\begin{example}\label{ex:nonconstant-pencil-v146}
Take \(n=4\), \(B=\Pj^1_{[s:t]}\),
\(\mathcal U=\OO(-1)\oplus\OO^3\), and \(V=\C^4\).  In diagonal
quadratic coordinates set
\[
 A=\operatorname{diag}(1,1,1,1),\qquad
 B_{s,t}=\operatorname{diag}(0,s,t,2s+3t).
\]
The map \(\OO\oplus\OO(-1)\to\Sym^2V\otimes\OO\) with columns
\(A,B_{s,t}\) has rank two everywhere.  The first diagonal entry
shows that any scalar relation involving \(A\) has zero coefficient
on \(A\); the next two entries then force \(s=t=0\), which is
impossible on \(\Pj^1\).  Its image \(\mathcal R\) is therefore a
locally direct-summand pencil bundle.  Since it contains the
nondegenerate form \(A\), it is regular at every point, including
all spectral collisions.

On the open set of four distinct eigenvalues, an ordered cross-ratio
of \(0,s,t,2s+3t\) is
\[
 \chi=\frac{t(s+3t)}{(2s+3t)(t-s)}.
\]
The rational function
\(J=(\chi^2-\chi+1)^3/(\chi^2(1-\chi)^2)\)
is invariant under the six cross-ratio transformations arising from
permutations of four points.  It is nonconstant: at \([s:t]=[1:2]\)
and \([1:3]\), the values of \(\chi\) are \(7/4\) and \(15/11\),
and the corresponding values of \(J\) differ.  Thus the pencil
family is not projectively isotrivial, even allowing a change of
parameter within each pencil.  The discriminant of
\(\det(zA-B_{s,t})\) is
\begin{equation}\label{eq:moving-discriminant-v146}
 4s^2t^2(2s+3t)^2(s-t)^2(s+3t)^2(s+t)^2.
\end{equation}
No factor is inverted in forming \(Y_{\mathcal R}^{[20]}\).
This finite scheme is locally free over its entire projective
reduction of rank \(\binom{36}{20}-45\).
By Theorem~\ref{thm:moving-reconstruction-v146}, its unmarked
isomorphism class recovers the whole morphism into the pencil
Grassmannian, including the six collision points and the pencil
there, up to one constant transformation of \(V\) and an automorphism
of the base.  This is a consequence about a genuinely moving family,
not just the pointwise Pluecker encoding.  The cross-ratio and
discriminant calculations are classical elementary readouts; their
role here is to exhibit a family to which the unmarked global inverse
applies.
\end{example}

\begin{remark}[Scope of the family statement]
\label{rem:moving-scope-v146}
The reduction in this section is smooth, projective and reduced.  The
source projective bundle is assumed nontrivial, which orients the two
rulings.  The theorem allows all fibre pencils, singular or regular;
regularity is used only in the spectral description of the last
example.  Theorem~\ref{thm:universal-finite-neighbourhood} separately
allows arbitrary complex base change with its specified relative
socle ideal.  The present intrinsic inverse does not replace that
ideal by an absolute nilradical over a nonreduced base.  Nor does
\eqref{eq:automorphism-sequence-v146} assert an equivalence of
unmarked moduli stacks or a description of every infinitesimal
deformation.  It is an exact classification of geometric
isomorphisms, their graded data, and the residual automorphism kernel
for the stated admissible objects.
\end{remark}

% BEGIN PRESERVED PART 38-pencil-deformation-equivalence-v149.tex
\section{The effective moduli stack of pencil failure algebras}
\label{sec:pencil-deformations}

We apply the common-divisor inverse to the actual first relations of
quadratic pencils.  Their common divisor has a higher multiplicity
than the single determinant factor used to recover their support.
Removing that multiplicity places every pencil, without a regularity
hypothesis, in the smooth stratum of the preceding section.

Keep \(W=\Sym^2V\), \(N=\binom{n+1}{2}\), \(p=N-2\), and put
\begin{equation}\label{eq:primitive-pencils}
 \begin{split}
 M&=\binom N2,\qquad s=n-1,\qquad h=3n-4,\\
 \mu&=(\underbrace{3,\ldots,3}_{n-2},2,0),\qquad
 A=\mathbb S_\mu V^*,\qquad F=\mathbb S_\mu U.
 \end{split}
\end{equation}
Both \(A\) and \(F\) have dimension \(M\).  The Cauchy inclusion
\(A\otimes F\subset S_h(E)\) will be fixed; its subspace and all
its projections are equivariant, independent of a choice of bases.

\begin{lemma}[Exact common multiplicity]\label{lem:primitive-pencil}
For every \(R\in\Gr(2,W)\), including every singular pencil, the
first relation space of its order-\(d\) failure algebra is
\begin{equation}\label{eq:primitive-factorization}
 K_R=\delta^{\,n-1}\widetilde J_R,\qquad
 \widetilde J_R=L_R^0\otimes F\subset S_h(E),
                   \qquad \gcd(\widetilde J_R)=1.
\end{equation}
Here \(L_R^0\subset A\) is the Pluecker line of \(R\), under the
canonical projective identification
\(\Pj(A)=\Pj(\bigwedge^2W)\).
The common divisor of \(K_R\) is therefore exactly
\(\delta^{n-1}\), up to scalar.  Its degree remains
\(d=sn+h=n^2+2n-4\).
\end{lemma}
\begin{proof}
Exterior duality and
\(\bigwedge^2\Sym^2H=\mathbb S_{(3,1)}H\) give
\[
 \bigwedge^{N-2}\Sym^2H
   =\mathbb S_{\lambda}H,\qquad
 \lambda=(\underbrace{n+1,\ldots,n+1}_{n-2},n,n-2)
   =\mu+(n-2)^n.
\]
Thus this representation is
\((\det H)^{n-2}\otimes\mathbb S_\mu H\).
Applied to the universal matrix, its coefficient polynomials factor
as \(\delta^{n-2}\) times the matrix coefficients of
\(\mathbb S_\mu T\).  The additional determinant in the failure
ideal gives \eqref{eq:primitive-factorization}.  On the left,
exterior duality also yields
\[
                A\simeq\bigwedge^2W\otimes(\det V)^{-3},
\]
which identifies the residual line with the Pluecker line.

We prove a slightly stronger assertion: the matrix coefficient
space \(\beta\otimes\mathbb S_\mu U\) has greatest common divisor
\(1\) for every nonzero \(\beta\in\mathbb S_\mu V^*\).
Its common divisor is taken to a scalar multiple of itself by the
full right group.  Its zero divisor is therefore invariant under
right multiplication.  A nonconstant invariant divisor cannot meet
the open orbit of invertible matrices.  Since the complement of
that orbit is the irreducible determinant hypersurface, any such
common divisor is a power of \(\delta\).

It remains to exclude \(\delta\).  Choose a rank-\(n-1\) map
\(T_0\).  The representation \(\mathbb S_\mu\) has length \(n-1\),
so \(\mathbb S_\mu T_0\ne0\).  Choose a nonzero vector in its image.
The span of its left \(\operatorname{GL}(V)\)-orbit is the entire
irreducible representation.  Consequently some left translate
\(gT_0\) has a coefficient not annihilated by \(\beta\).
Not all residual coefficients vanish on the rank-\(n-1\) locus.
They have no determinant factor, proving the assertion.  The same
argument at an invertible matrix shows that the coefficient vector
is nonzero everywhere on the invertible locus.
\end{proof}

Let \(\mathscr P=\Gr(2,W)\).  The morphism
\begin{equation}\label{eq:primitive-pencil-embedding}
 \iota:\mathscr P\longrightarrow X_{h,M}(E),
                 \qquad R\longmapsto L_R^0\otimes F
\end{equation}
is a closed immersion.  Indeed, the Pluecker embedding is followed
by the closed immersion \(\Pj(A)\to\Gr(M,A\otimes F)\),
\(L\mapsto L\otimes F\), and then by a linear Grassmannian
inclusion.  The inverse to the middle map is contraction with
\(F^*\).  This factorization proves the assertion on schemes,
including nonreduced test schemes, rather than only on points.
Lemma~\ref{lem:primitive-pencil} shows that its image lies in the
open set \(X_{h,M}(E)\).

Define the \emph{pencil relation stratum} \(Z_{\rm pen}\) inside
\(Z^{\det}_{n-1,3n-4,M}\) by the following intrinsic scheme
conditions.  Recover the common-divisor line, its determinant root,
and its unordered tensor rulings by
Theorem~\ref{thm:determinant-stratum}.  In an ordering, the primitive
residual relation is required to lie in the indicated Cauchy
summand, to have the form \(L\otimes F\), and to satisfy the
Pluecker equations on \(L\) under
\(A\simeq\bigwedge^2W\otimes(\det V)^{-3}\).
These are equations on subbundles; a test on the fibres alone is
not their definition.  Use both orderings and descend.  The two
ordered images are disjoint: their support ranks are respectively
\((1,M)\) and \((M,1)\), with \(M>1\).  Each is closed in the
residual Grassmannian.  Thus this prescription defines a locally
closed scheme, independent of the tensor frames.  It is an
intrinsic image criterion for a natural family of first relations,
not a definition by an algebra isomorphism to a preselected pencil.

\begin{theorem}[Effective all-pencil moduli and deformations]
\label{thm:pencil-stack-equivalence}
The stack of first relations in \(Z_{\rm pen}\) is
\begin{equation}\label{eq:pencil-stack-before-rigidification}
 [Z_{\rm pen}/\operatorname{GL}(E)]\simeq
                         [\mathscr P/G^\circ].
\end{equation}
There is an exact sequence of groups
\begin{equation}\label{eq:pencil-ineffective-group}
 1\longrightarrow\operatorname{GL}(U)\longrightarrow G^\circ
                 \longrightarrow\PGL(V)\longrightarrow1,
\end{equation}
and its kernel acts trivially on \(\mathscr P\).
Quotient the morphisms in the algebra stack first by the
nonlinear tangent-identity kernel of
Proposition~\ref{prop:relative-first-relation}, and then by this
ineffective right group, and fppf stackify.  The resulting
\emph{effective pencil failure stack} satisfies
\begin{equation}\label{eq:effective-pencil-stack}
       \mathscr F_{\rm pen}^{\rm eff}\simeq
                [\Gr(2,\Sym^2V)/\PGL(V)].
\end{equation}
This equivalence is valid over all complex parameter schemes.
It identifies automorphisms after the specified quotients,
infinitesimal deformation groupoids, and families of degenerating
pencils, including singular pencils.

The stack is smooth of dimension \(n-3\).  At \(R\), its tangent
complex is
\begin{equation}\label{eq:pencil-tangent-complex}
 [\,\mathfrak{pgl}(V)\xrightarrow{\ d_R\ }\Hom(R,W/R)\,],
 \qquad d_R(X)(r)=(X\otimes1+1\otimes X)r\bmod R.
\end{equation}
In particular, the dimension of its space of first-order
isomorphism classes is \(n-3+\dim\operatorname{Stab}_{\PGL(V)}(R)\).
Within this stratum, every infinitesimal pencil deformation and
its effective failure-algebra deformation lift together across a
small extension of local Artinian complex algebras.
\end{theorem}
\begin{proof}
In the determinant frame of
Theorem~\ref{thm:determinant-stratum}, the residual pencil locus is
\(\iota(\mathscr P)\amalg\tau\iota(\mathscr P)\), where
\(\tau\) exchanges the two factors.  The disjointness just proved
and the equivariance of the Cauchy and Pluecker maps give, as
schemes with group action,
\[
 \iota(\mathscr P)\amalg\tau\iota(\mathscr P)
                         \simeq G\times^{G^\circ}\mathscr P.
\]
Quotienting by \(G\) gives
\eqref{eq:pencil-stack-before-rigidification}.  All maps used here
are morphisms and all inverse operations are contractions,
common-divisor descent, or the inverse on the Pluecker image.
Consequently this calculation remains valid on a nonreduced base.
It is not an inference from equality of complex-point stabilizers.

The map in \eqref{eq:pencil-ineffective-group} sends the class of
\((g,h)\) to \([g]\).  If \(g=aI\), multiplication by the diagonal
scalar \(a^{-1}\) writes its class uniquely as \((I,a^{-1}h)\).
This proves the asserted kernel.  The action on the Pluecker line
is the projective action of \(g\); the full right factor has no
effect on \(R\).  Every \(\PGL(V)\)-valued isomorphism lifts fppf
locally to \(G^\circ\), and two lifts differ by that kernel.
The quotient of morphisms and stackification therefore gives
\([\mathscr P/\PGL(V)]\).  Combining this with
Proposition~\ref{prop:relative-first-relation} proves
\eqref{eq:effective-pencil-stack}.  No splitting of
\eqref{eq:pencil-ineffective-group} over the base is required.

A Grassmannian chart writes a nearby pencil as the graph of a map
\(R\to W/R\).  Infinitesimal coordinate changes act by \(d_R\),
giving \eqref{eq:pencil-tangent-complex}.  Grassmannian charts and
\(\PGL(V)\) are smooth, so the quotient is smooth and the graph
coefficients lift over every small extension.  A coordinate change
can also be lifted in the smooth group, and its action determines
a compatible lift of the corresponding object.  Torsors over a local
Artinian complex algebra can be trivialized, since smooth torsors
lift a point of their algebraically closed residue field.
This proves the lifting assertion for deformation groupoids.
Finally,
\[
 \dim\mathscr P-\dim\PGL(V)=2(N-2)-(n^2-1)=n-3.
\]
Taking kernel and cokernel of \(d_R\) gives the stated first-order
dimension, including stabilizer jumps at singular pencils.
\end{proof}

\begin{corollary}[Native specialization loci]\label{cor:pencil-specialization-stack}
Every \(\PGL(V)\)-invariant locally closed subscheme
\(Q\subset\Gr(2,\Sym^2V)\), with its specified scheme structure,
has a corresponding intrinsic effective failure stratum
isomorphic to \([Q/\PGL(V)]\).  Its infinitesimal equations and its
specializations are retained, not just its closed points.
In particular, this applies to rank-incidence and spectral Fitting
conditions on the universal pencil, including their intersections.
A family in such a locus specializes to a pencil in its closure
exactly when its effective failure family has the corresponding
specialization in the associated closure stratum.
\end{corollary}
\begin{proof}
Restrict the equivariant closed immersion
\eqref{eq:primitive-pencil-embedding} and its inverse to \(Q\).
The construction of \(Z_{\rm pen}\) and its two quotient steps
then restrict to give the stated isomorphism.  Equations for a
rank or a Fitting condition are equivariant equations in the
universal pencil, so they are included in this argument with their
scheme structures.  Apply the same construction to the closure of
\(Q\) to obtain the specialization assertion.
\end{proof}

\begin{remark}\label{rem:moduli-scope}
Theorem~\ref{thm:pencil-stack-equivalence} does not say that the
unrigidified algebra stack is the pencil stack.  At a complex
pencil its linear stabilizer remains
\((G_R\times\operatorname{GL}(U))/\mathbb G_m\), and its full
algebra automorphism group also has the nonlinear kernel in
Proposition~\ref{prop:first-relation-group-v147}.  Both are real
automorphisms, not redundant notation.  They are explicitly removed
only in the effective comparison.  Nor is the quotient asserted to
be proper, although its framed Grassmannian is projective.

The parameter base in this theorem is distinct from the intrinsic
projective reduction in Theorem~\ref{thm:unrestricted-moving-v146}.
Here projective coordinate changes may vary over the parameter
scheme.  In that moving theorem, properness of the intrinsic
reduction forces one constant left map and recovers an actual
source bundle.  Neither assertion is substituted for the other.
The present equivalence upgrades the geometric orbit statement to
a family statement without changing the hypotheses or conclusions
of the moving inverse.
\end{remark}

% BEGIN PRESERVED PART 41-rigidification-v151.tex
\section{Universal properties of the two rigidifications}
\label{sec:rigidification-v151}

The effective stacks used above have a universal characterization.
This specifies what is forgotten from isomorphisms, as opposed to what
is retained from objects.  We use rigidification in its standard sense
of removing a flat normal subgroup of inertia; see
\cite[Tag~04V2]{StacksRigidV151}.  The argument below gives the
universal property directly in the present situation, including the
nonconstant subgroup obtained by descent.

Let \(H_0\subset\Gr(r,S_D(E))\) be a
\(\operatorname{GL}(E)\)-invariant locally closed subscheme.  Let
\(\mathscr A_{H_0}\) be the stack of algebras satisfying the trace,
filtration and Hilbert-function conditions of
Proposition~\ref{prop:relative-first-relation}, whose intrinsic
kernel is in \(H_0\).  Write
\(\mathscr R_{H_0}=[H_0/\operatorname{GL}(E)]\).

\begin{theorem}[First-relation rigidification]
\label{thm:rigidification-v151}
Taking the first relation defines an fppf gerbe
\[
                      q:\mathscr A_{H_0}\longrightarrow\mathscr R_{H_0}.
\]
Its relative inertia \(\mathscr U\) consists exactly of the
substitutions which induce the identity on \(\mathcal N/\mathcal N^2\).
It is a smooth affine normal subgroup of the inertia.  In a homogeneous
presentation its underlying scheme is
\(\Hom(\mathcal E,\mathcal N^2)\), with the substitution group law.
For every stack \(\mathscr Z\) in groupoids on complex schemes,
precomposition with \(q\) identifies
\begin{equation}\label{eq:rigidification-universal-v151}
 \Hom(\mathscr R_{H_0},\mathscr Z)
 \simeq
 \{F\in\Hom(\mathscr A_{H_0},\mathscr Z):F(\mathscr U)=1\}.
\end{equation}
The right side is the full subgroupoid, including all natural
isomorphisms, of functors killing the indicated relative inertia.
The construction is compatible with base change on the relation stack.
\end{theorem}
\begin{proof}
A relation bundle has its homogeneous truncated algebra as a lift.
Every algebra over that relation bundle is locally isomorphic to this
lift.  The first-relation proposition shows that an isomorphism of
relation bundles lifts fppf locally, and any two lifts differ by a
unique element of the indicated kernel.  These assertions are exactly
the gerbe and relative-inertia assertions.  The kernel is intrinsic,
so it is stable under conjugation.  Its local affine-space description
shows smoothness and flatness; no additive group law is being asserted.

For clarity, we prove the descent implicit in
\eqref{eq:rigidification-universal-v151}.  Given \(F\) on the right
and an object \(y\) of \(\mathscr R_{H_0}\), choose local lifts
\(a_i\) over an fppf cover.  Locally choose comparison isomorphisms
on overlaps above the identity of \(y\).  Two choices differ by
\(\mathscr U\), hence have the same image under \(F\).  On triple
overlaps the failure of the chosen lifts to satisfy the cocycle
condition is again in \(\mathscr U\).  Their images therefore give
a genuine descent datum for \(F(a_i)\), which is effective because
\(\mathscr Z\) is a stack.  Refinements and different lift choices
give canonically isomorphic descended objects.  An arrow between two
objects of \(\mathscr R_{H_0}\) is treated in the same way, using
local lifts of that arrow.  This constructs the descended functor.
Natural isomorphisms descend by the same argument, proving full
faithfulness as well as essential surjectivity.  Conversely every
functor pulled back by \(q\) kills its relative inertia.  All the
constructions use covers and isomorphism sheaves, so commute with
base change.
\end{proof}

\begin{corollary}[The effective pencil stack is a successive rigidification]
\label{cor:two-rigidifications-v151}
On the pencil relation stratum the second morphism
\[
 [\mathscr P/G^\circ]\longrightarrow[\mathscr P/\PGL(V)],
                     \qquad \mathscr P=\Gr(2,\Sym^2V),
\]
is the rigidification by the normal relative inertia induced from
\(\ker(G^\circ\to\PGL(V))=\operatorname{GL}(U)\).
Consequently Theorem~\ref{thm:pencil-stack-equivalence} has the
universal property of first killing \(\mathscr U\) and then this
right-factor inertia.  It does not kill the stabilizer of the pencil
in \(\PGL(V)\).
\end{corollary}
\begin{proof}
The right factor acts trivially on \(\mathscr P\), and the quotient
map of groups is fppf surjective.  Objects and isomorphisms of the
second quotient therefore lift locally to the first, with differences
in that kernel.  Apply the descent proof of
Theorem~\ref{thm:rigidification-v151}.  Over an object represented by
a \(G^\circ\)-torsor the subgroup is its conjugation-associated
\(\operatorname{GL}(U)\)-bundle; it need not be a constant group
scheme on the base.  This describes the actual subgroup stack and
shows that no global splitting is assumed.  Composing with the first
rigidification gives the asserted successive universal property.
All remaining pencil automorphisms survive in the target.
\end{proof}

% BEGIN PRESERVED PART 42-pencil-specialization-v151.tex
\section{A common closed specialization of all quadratic pencils}
\label{sec:closed-pencil-v151}

We give an orbit-closure consequence directly on the pencil side.
Unlike restricting an equivalence to a previously specified locus, the
following argument constructs the specialization.  It works for every
pencil, including pencils with singular generic member.  Its mechanism
is elementary projective representation theory; the additional failure
algebra statement uses the all-pencil, exact-multiplicity construction.

\begin{theorem}[Explicit common closed orbit]
\label{thm:closed-pencil-v151}
Let \(\dim V=n\geq3\).  The action of \(\PGL(V)\) on
\(\mathscr P=\Gr(2,\Sym^2V)\) has a unique closed orbit, namely
\begin{equation}\label{eq:flag-orbit-v151}
 \mathscr C=\{\ell H:\ell\subset H\subset V,\
                      \dim\ell=1,\ \dim H=2\}.
\end{equation}
It is the flag variety \(\operatorname{Fl}(1,2;V)\), of dimension
\(2n-3\).  More precisely, after a fixed change of coordinates every
pencil has a degeneration along the one-parameter subgroup
\[
                  \lambda(t)=\diag(1,t,t^2,\ldots,t^2)
\]
to \(\langle x_1^2,x_1x_2\rangle\).  The pencil family is a
subbundle on the whole affine line.
\end{theorem}
\begin{proof}
Take independent forms \(q_0,q_1\).  The rational function
\(q_1/q_0\) on \(\Pj(V^*)\) is nonconstant.  In characteristic
zero its differential is not identically zero: on an affine coordinate
chart the common kernel of the coordinate derivations of the rational
function field is \(\C\).  Choose a point \([v]\) where \(q_0(v)\ne0\)
and this differential is nonzero.  Replace \(q_1\) by
\(q_1-(q_1(v)/q_0(v))q_0\).  Then \(q_1(v)=0\), whereas its
differential has a nonzero value on some vector \(w\) independent of
\(v\).  Euler's identity shows that its value on the radial vector
\(v\) is zero.  Extending \(v,w\) to a basis of \(V^*\) gives dual
coordinates in which
\[
 [x_1^2]q_0=a\ne0,\qquad [x_1^2]q_1=0,
                                  \qquad [x_1x_2]q_1=b\ne0.
\]
Here brackets denote a coefficient.

Apply \(\lambda(t)\) to the variables.  The forms
\[
 Q_0(t)=q_0(x_1,tx_2,t^2x_3,\ldots,t^2x_n),\qquad
 Q_1(t)=t^{-1}q_1(x_1,tx_2,t^2x_3,\ldots,t^2x_n)
\]
have polynomial coefficients: the only possible weight-zero term in
\(q_1\) was removed.  Their limits are \(a x_1^2\) and
\(b x_1x_2\).  For \(t\ne0\) they span a coordinate transform of
\(R\); at zero they are independent.  Thus their coefficient matrix
has rank two everywhere, defines a subbundle, and supplies the asserted
morphism from \(\A^1\) to \(\mathscr P\).

The orbit of \(\langle x_1^2,x_1x_2\rangle\) is exactly
\(\mathscr C\).  A pencil of the form \(\ell H\) recovers \(\ell\)
as its gcd and \(H\) by division.  These are scheme-theoretic inverse
operations on the exact-gcd-one-degree stratum by
Corollary~\ref{cor:canonical-exact-v151}, applied to \(g=h=1,r=2\).
Under that inverse the condition \(\ell\subset H\) is the usual
closed flag incidence.  The flag-to-pencil map is therefore a locally
closed immersion; since the flag variety is projective it is a closed
immersion.  It identifies \(\mathscr C\) with the flag variety, and
the dimension is \((n-1)+(n-2)\).  Every pencil orbit closure contains
\(\mathscr C\), by the constructed family and equivariance.  Any
closed orbit must consequently be \(\mathscr C\), proving uniqueness.
\end{proof}

\begin{corollary}[Flat closed specialization of failure algebras]
\label{cor:closed-failure-v151}
For every quadratic pencil \(R\) there is a finite locally free
family over \(\A^1\), of rank
\(\binom{n^2+d}{d}-\binom N2\), whose nonzero fibres are isomorphic
to \(\mathfrak A_R^{[d]}\) and whose zero fibre is
\(\mathfrak A_{\langle x_1^2,x_1x_2\rangle}^{[d]}\).
The intrinsic first relations throughout the family have exact gcd
\(\delta^{n-1}\).  In a fixed determinant frame, the two ordered
copies of the flag orbit form the unique closed \(G\)-orbit in the
pencil relation locus.  Equivalently the effective pencil failure
stack has the common closed specialization described in
Theorem~\ref{thm:closed-pencil-v151}.
\end{corollary}
\begin{proof}
Let \(R_t\) be the polynomial subbundle just constructed.  Its
Pluecker line gives, under the fixed exterior-duality identification,
a line bundle \(L_t^0\subset\mathbb S_\mu V^*\otimes\OO_{\A^1}\).
Set
\[
 \mathcal K_t=\delta^{n-1}(L_t^0\otimes\mathbb S_\mu U),\qquad
 \mathcal A_t=\Sym(E\otimes\OO_{\A^1})/
                       ((\mathcal K_t)+\mathfrak m^{d+1}).
\]
The first is a rank-\(M\) subbundle.  The second is finite locally
free by its graded direct-sum presentation, with the claimed rank.
Lemma~\ref{lem:primitive-pencil} gives the exact gcd at every fibre,
including zero; in fact this family lies scheme-theoretically in the
fixed-divisor exact stratum.  Equivariance and the local inverse
identify its nonzero fibres as asserted.  In a determinant frame the
pencil locus is \(G\times^{G^\circ}\mathscr P\).  The previous
closed-orbit calculation therefore gives the assertion for \(G\).
This is a statement inside that locus; it does not assert closedness
under arbitrary degenerations of the determinant form in the full
relation Grassmannian.
\end{proof}

The closed specialization is very coarse compared with the unmarked
local inverse: many pairwise nonisomorphic finite algebras have this
same limiting fibre.  The full first relation, not its common limit,
retains the original pencil.  The explicit family also differs from
transport of spectral rank loci: it crosses from arbitrary pencils to
a singular common-factor pencil without a regularity hypothesis.
Theorem~\ref{thm:boundary-normalization-v151}, on the other hand,
describes divisor collisions in the larger natural first-relation
boundary, outside the exact pencil image.  These are two distinct
geometric statements, and neither is a classification of every
congruence-orbit adjacency of quadratic pencils.

% BEGIN PRESERVED PART 39-transverse-relations-v149.tex
\section{Transverse deformations with unchanged determinant support}
\label{sec:transverse-relations}

The larger determinant-divisor stratum has deformations which are
not pencil deformations and need not satisfy either factor-invariance
condition.  We now identify their tangent directions and construct
complete families in which the homogeneous radical remains exactly
the determinant ideal.  Thus the intrinsic stratum theorem has content
outside the coefficient class, not merely a different description
of the same symmetry hypothesis.

Let \(J=L\otimes F\) be the primitive relation of a pencil
\(R\in\Gr(2,W)\).  In degree \(h=3n-4\), write the equivariant
Cauchy decomposition as
\[
                    S_h(E)=(A\otimes F)\oplus H_{\rm other},
\]
where \(A,F,L\) are as in \eqref{eq:primitive-pencils}.
The summand \(H_{\rm other}\) contains every partition different
from \(\mu\); this notation denotes its full vector space, not a
quotient chosen separately at each pencil.

\begin{theorem}[Normal directions to the pencil stratum]
\label{thm:transverse-pencil-directions}
The normal space at \(J\) of the pencil locus in \(X_{h,M}(E)\)
has the following equivariant direct-sum description:
\begin{equation}\label{eq:three-transverse-modes}
 \begin{split}
 N_J\simeq{}&
 \Hom(\bigwedge^2R,\bigwedge^2(W/R))\\
 &\oplus\bigl(\Hom(L,A/L)\otimes\mathfrak{sl}(F)\bigr)\\
 &\oplus\Hom(L\otimes F,H_{\rm other}).
 \end{split}
\end{equation}
The summands respectively record departure from the Pluecker locus,
departure from a full right multiplicity factor inside the same
Cauchy summand, and departure into other Cauchy summands.  They are
all transverse directions in the effective first-relation stack:
the infinitesimal \(G^\circ\)-orbit lies in the pencil tangent
space, so no vector in \(N_J\) is removed by that orbit.  Every
such direction is the first-order term of a local algebraic family
in the determinant-divisor stratum.  This last assertion refers to
the common-divisor stratum; preservation of the entire homogeneous
radical for complete families is supplied below.
\end{theorem}
\begin{proof}
The tangent space of a Grassmannian and the Cauchy splitting give
\begin{equation}\label{eq:ambient-cauchy-tangent}
 \Hom(J,S_h(E)/J)=
 \bigl(\Hom(L,A/L)\otimes\End(F)\bigr)
       \oplus\Hom(L\otimes F,H_{\rm other}).
\end{equation}
Under \(L\mapsto L\otimes F\), the tangent to \(\Pj(A)\) is the
first factor tensored with the identity of \(F\).
In characteristic zero the trace splits
\(\End(F)=\C I\oplus\mathfrak{sl}(F)\).
It remains to identify the normal space of the Pluecker embedding.
The natural map
\(\bigwedge^2W\to\bigwedge^2(W/R)\) has kernel \(R\wedge W\).
At the line \(\bigwedge^2R\), the tangent image is
\(\Hom(\bigwedge^2R,(R\wedge W)/\bigwedge^2R)\), which is the
image of \(\Hom(R,W/R)\).  Its quotient in the projective tangent
space is
\(\Hom(\bigwedge^2R,\bigwedge^2(W/R))\).
The common determinant twist in \(A\) cancels in this Hom space.
Substitution in \eqref{eq:ambient-cauchy-tangent} proves
\eqref{eq:three-transverse-modes} as a normal-space isomorphism.
No splitting of the tangent sequence of the Pluecker embedding is
being asserted.

The left group preserves the Pluecker locus, and the right group
fixes \(L\otimes F\); hence the differential of its orbit is in
the pencil tangent.  Passing to the quotient tangent complexes in
\eqref{eq:ambient-tangent-complex} does not change this normal
quotient.  Finally \(X_{h,M}(E)\) is open in a Grassmannian.  A
tangent vector is integrated by a graph with coefficients linear
in one parameter, after restricting to a neighbourhood of zero
which stays in this open set.  Multiplication by \(\delta^{n-1}\)
and Proposition~\ref{prop:relative-first-relation} give the claimed
family of finite flat algebras.
\end{proof}

\begin{theorem}[Complete fixed-support families through every pencil]
\label{thm:fixed-support-grassmannians}
Put \(e=n^2\).  For every pencil \(R\) there is an \(e\)-dimensional
subspace \(B\subset J=\widetilde J_R\) such that
\(\gcd(B)=1\) and the forms in \(B\) have no common zero on the
invertible matrix locus.  The projective Grassmannian
\begin{equation}\label{eq:fixed-support-grassmannian}
 \Sigma_B=\{J'\in\Gr(M,S_h(E)):B\subset J'\}
          \simeq\Gr(M-e,S_h(E)/B)
\end{equation}
contains \(J\), is contained in \(X_{h,M}(E)\), and satisfies
\begin{equation}\label{eq:fixed-radical}
              \sqrt{(\delta^{n-1}J')}=(\delta)
                         \quad\text{for every geometric }J'\in\Sigma_B.
\end{equation}
Its universal first relation gives a finite flat family of the
prescribed Hilbert function.  It contains a projective line through
\(J\) whose every other point is invariant under neither full
factor group.  Consequently determinant-radical relations through
each pencil admit positive-dimensional complete families outside
the recognized coefficient class, and their intrinsic deformation
problem is described by Theorem~\ref{thm:determinant-stratum}.
\end{theorem}
\begin{proof}
By Lemma~\ref{lem:primitive-pencil}, the residual coefficient vector
is nonzero on every invertible matrix.  Evaluation gives a morphism
\[
 \operatorname{GL}_n/\mathbb G_m\longrightarrow\Pj(J^*).
\]
Its projective image closure \(Y\) has dimension at most \(e-1\).
Since \(M>e\) for \(n\ge3\), a general projective linear subspace
of codimension \(e\) in \(\Pj(J^*)\) misses \(Y\).  For such a
subspace \(\Pj(B^\perp)\), the associated \(e\)-plane \(B\subset J\)
has no common zero on the invertible locus.  To justify the generic
avoidance directly, the incidence of pairs consisting of a point
of \(Y\) and an \(e\)-plane annihilated by it has dimension strictly
smaller than \(\Gr(e,J)\).  Its projection is closed because
\(Y\) is projective.

The condition \(\gcd(B)=1\) is a second nonempty open condition
on \(\Gr(e,J)\).  To see nonemptiness, take any nonzero \(b_1\in J\).
It has only finitely many irreducible factors.  Since \(\gcd(J)=1\),
a general \(b_2\in J\) avoids divisibility by each such factor;
then \(\gcd(b_1,b_2)=1\).  Extend their span to an \(e\)-plane.
Openness follows from Lemma~\ref{lem:universal-gcd}.  The two
nonempty opens intersect, as \(\Gr(e,J)\) is irreducible.

Every \(J'\) containing this \(B\) has greatest common divisor
\(1\), and its forms cannot vanish simultaneously on an invertible
matrix.  The common zero set of \(\delta^{n-1}J'\) is thus exactly
\(V(\delta)\).  The Nullstellensatz over an algebraically closed
field proves \eqref{eq:fixed-radical}.  This is a statement about
geometric homogeneous radicals.  The scheme-theoretic family used
here is the universal subbundle \(\delta^{n-1}\mathcal J'\), not
an operation of taking radicals on a nonreduced parameter ring.
Multiplication by this fixed nonzero form gives a subbundle;
its quotient in degree \(d\) and all the lower symmetric powers
are locally free.  The truncated algebras are therefore finite flat.

For the final assertion, choose \(j_0\in J\setminus B\), a
complement \(C\) to \(\C j_0\) in \(J\) containing \(B\), and
a nonzero decomposable vector
\[
 \eta\in\Sym^h V^*\otimes\Sym^h U\subset H_{\rm other}.
\]
The inequality \(M>e\) allows the choice of \(C\).
The family
\[
      J_{[a:b]}=C\oplus\C(aj_0+b\eta),\qquad[a:b]\in\Pj^1,
\]
is a projective line in \(\Sigma_B\) through \(J=J_{[1:0]}\).
For \(b\ne0\) its projection to the \((h)\)-Cauchy summand is
exactly the line \(\C\eta\).  A decomposable line in
\(\Sym^h V^*\otimes\Sym^h U\) is invariant under neither full
factor group: each of the two irreducible factors has dimension
greater than one.  Equivariance of the Cauchy projection therefore
excludes both possible admissible orientations of \(J_{[a:b]}\).
The determinant-preserver theorem says that these are the only
orientations after any change of the recovered matrix coordinates.
This proves the assertion intrinsically, not only in the chosen
matrix frame.
\end{proof}

\begin{remark}\label{rem:scope-transverse}
The three normal modes and the complete fixed-support families
separate two questions.  The effective pencil stratum has exactly
the deformation theory of quadratic pencils, including their
singular boundary.  The larger determinant-divisor stratum contains
additional directions, explicitly described here, which cease to
be pencil or coefficient relations.  Retaining the determinant
radical does not suppress these directions.  Conversely, no
classification of all ideals with a prescribed radical is needed
to describe this larger locally closed stratum and its family
inverse.  These assertions enlarge the reconstruction problem
without replacing the all-pencil theorem by a generic statement.
\end{remark}

% BEGIN PRESERVED PART 36-covering-moduli-v148.tex
\section{Covering maps inside one fibre orbit of quadratic pencils}
\label{sec:covering-moduli-v148}

We give a consequence native to the pencil construction.  It embeds
the equivalence problem for maps \(\Pj^1\to\Pj^1\) into
unmarked finite failure schemes, while keeping the fibre algebra
and every nilradical graded vector bundle fixed.  Unlike a varying
spectral invariant, the phenomenon occurs within a single orbit
of singular pencils.

Fix \(n\ge3\), a two-dimensional subspace
\(H=\langle x,y\rangle\subset V\), and the trivial source
bundle \(\mathcal U=\OO_{\Pj^1}^n\).  Let
\(f=[a:b]:\Pj^1\to\Pj(H)\) be a morphism of degree
\(m\ge1\), where \(a,b\) are homogeneous of degree \(m\)
with no common zero.  Its tautological line subbundle is
\[
 \mathcal L_f=\OO(-m)
       \xrightarrow{\ (a,b)\ }H\otimes\OO_{\Pj^1}.
\]
Multiplication in \(\Sym V\) defines a rank-two subbundle
\begin{equation}\label{eq:cover-pencil-v148}
 \mathcal R_f=\mathcal L_f H
       \subset\Sym^2V\otimes\OO_{\Pj^1},
 \qquad
 \mathcal R_f|_{[s:t]}=
   \langle (ax+by)x,(ax+by)y\rangle.
\end{equation}
Let \(Y_f=Y_{\mathcal R_f}^{[d]}\) be its order-\(d\)
failure thickening, with
\[
 N=\binom{n+1}{2},\qquad p=N-2,\qquad
 d=n^2+2n-4,\qquad M=\binom N2,\qquad e=n^2.
\]
All schemes below are compared without projection, grading,
embedding, matrix factors, or coefficient markings.

\begin{theorem}[Unmarked realization of the covering equivalence problem]
\label{thm:covering-pencil-moduli-v148}
For fixed \(n\ge3\) and \(m\ge1\), the schemes \(Y_f\)
have the following properties.

\smallskip
\noindent (i) They are finite locally free over their reduction
\(\Pj^1\), of rank \(\binom{e+d}{d}-M\), and are
Zariski locally trivial with the same Artin algebra as fibre.
Every fibre pencil is equivalent to \(\langle x^2,xy\rangle\).

\smallskip
\noindent (ii) Writing \(\mathfrak n_f\) for their nilradicals,
their graded vector bundles are independent of \(f\):
\begin{equation}\label{eq:cover-graded-bundles-v148}
 \begin{split}
 \mathfrak n_f^j/\mathfrak n_f^{j+1}
   &\simeq\OO^{\binom{e+j-1}{j}}\qquad(1\le j<d),\\
 \mathfrak n_f^d
   &\simeq\OO^{\binom{e+d-1}{d}-3M}
                              \oplus\OO(m)^{\oplus2M}.
 \end{split}
\end{equation}
Their abstract first-relation bundles are all
\(\OO(-2m)^{\oplus M}\), and their pencil and quotient bundles
are respectively
\(\OO(-m)^2\) and \(\OO(2m)\oplus\OO^{N-3}\).

\smallskip
\noindent (iii) For two degree-\(m\) maps,
\begin{equation}\label{eq:cover-equivalence-v148}
 Y_f\simeq Y_{f'}
 \quad\Longleftrightarrow\quad
 f'=\alpha\circ f\circ\sigma
 \quad\text{for some }\alpha,\sigma\in\PGL_2(\C).
\end{equation}
Thus the double-orbit set of degree-\(m\) morphisms
\(\Pj^1\to\Pj^1\) is realized faithfully by unmarked
quadratic-pencil failure schemes with all the additive data in
(ii) fixed.  This is a statement about isomorphism classes, not an
equivalence of their automorphism groupoids.
\end{theorem}
\begin{proof}
\emph{The pencil and local triviality.}
Multiplication by a nonzero linear form is injective on \(H\).
Since \((a,b)\) never vanishes, the map
\(\OO(-m)\otimes H\to\Sym^2H\otimes\OO\) has rank
two everywhere.  Its image is a subbundle with locally free
quotient.  In the basis \(x^2,xy,y^2\) its two columns are
\begin{equation}\label{eq:cover-columns-v148}
                 \begin{pmatrix}a&0\\ b&a\\0&b\end{pmatrix},
 \qquad\text{with Pluecker coordinates }(a^2,ab,b^2).
\end{equation}
It follows that \(\mathcal R_f\simeq\OO(-m)^2\).
The quotient of \(\Sym^2H\otimes\OO\) is a line bundle
of degree \(2m\).  A constant complement to \(\Sym^2H\)
in \(\Sym^2V\) gives
\(\mathcal S_f\simeq\OO(2m)\oplus\OO^{N-3}\).
On an open set where \(\mathcal L_f\) is trivialized, complete
its nonvanishing generator to a frame of \(H\otimes\OO\).
The inverse frame map takes \(\mathcal R_f\) to
\(\langle x^2,xy\rangle\).  Extend it by the identity on a
fixed complement of \(H\) in \(V\).  With the source bundle
trivial, it identifies the defining minors and the truncated
algebra with a product on that open set.  This proves local
triviality, not merely equality of fibre isomorphism classes.
Lemma~\ref{lem:moving-coefficients-v146} gives finite local
freeness and the asserted rank.

\emph{The entire graded bundle calculation.}
The first relations occur in degree \(d\), so all lower
positive graded pieces are the full constant symmetric powers of
\(E=V^*\otimes\C^n\).  Put
\(F(V)=\bigwedge^p\Sym^2V\), of dimension \(M\).
The coefficient line
\(\mathcal L_{\mathcal R_f}=\bigwedge^p\mathcal S_f^*\)
is \(\OO(-2m)\).  Under exterior duality
\[
 F(V)^*\simeq\bigwedge^2(\Sym^2V)
                              \otimes\det(\Sym^2V)^*,
\]
its inclusion is the Pluecker line of \(\mathcal R_f\).
The constant determinant twist does not affect the bundle
calculation.  Thus it lies in a fixed three-dimensional subspace
of \(F(V)^*\), with column \((a^2,ab,b^2)\).
There is an exact sequence
\begin{equation}\label{eq:conic-quotient-v148}
 0\longrightarrow\OO(-2m)
 \xrightarrow{\ (a^2,ab,b^2)^{\mathsf t}\ }\OO^3
 \xrightarrow{\left(\begin{smallmatrix}-b&a&0\\0&-b&a\end{smallmatrix}\right)}
                  \OO(m)^2\longrightarrow0.
\end{equation}
Indeed the composite is zero; the maximal minors of the second
map are \(b^2,-ab,a^2\), which have no common zero.  It is
therefore everywhere surjective.  Its rank-one kernel is the
image of the displayed nowhere vanishing column.  Taking a
constant complement proves
\begin{equation}\label{eq:coefficient-quotient-v148}
 (F(V)^*\otimes\OO)/\mathcal L_{\mathcal R_f}
                   \simeq\OO^{M-3}\oplus\OO(m)^2.
\end{equation}

The fixed Cauchy inclusion, followed by multiplication by the
fixed determinant, embeds
\(F(V)^*\otimes F(\C^n)\) as an \(M^2\)-dimensional
constant subspace of \(\Sym^d E\).
The first relation bundle is precisely
\(\mathcal L_{\mathcal R_f}\otimes F(\C^n)\) in that
subspace.  Choose a constant complement to the whole subspace,
not to its varying relation subbundle.  The top graded quotient is
therefore
\[
 \OO^{\binom{e+d-1}{d}-M^2}
 \oplus\left((F(V)^*\otimes\OO)/
                         \mathcal L_{\mathcal R_f}\right)^{\oplus M}.
\]
Equation~\eqref{eq:coefficient-quotient-v148} gives exactly
\eqref{eq:cover-graded-bundles-v148}.  It also shows that the
first relation bundle is \(\OO(-2m)^{\oplus M}\).
All assertions in (ii) concern vector bundles; their
multiplication maps have not been identified.

\emph{The global inverse is the covering equivalence.}
Suppose \(Y_f\simeq Y_{f'}\).  Their reductions give an
automorphism of \(\Pj^1\).  By
Theorem~\ref{thm:unrestricted-moving-v146}, or the actual
descent theorem with the pencil coefficient line, the varying
pencil subbundles are identified by that base map and a single
constant \(g\in\operatorname{GL}(V)\).
For a nonzero \(\ell\in H\), the common polynomial divisor
of the two-dimensional space \(\ell H\subset\Sym^2V\)
is precisely \(\ell\), up to scalar: \(x\) and \(y\)
have no common nonconstant divisor in \(\Sym V\).
Dividing the pencil space by this common factor recovers \(H\).
These operations are equivariant under \(\operatorname{GL}(V)\).
Consequently an equality between \(g(\ell H)\) and a pencil
of the form \(\ell'H\) forces \(gH=H\) and
\([g\ell]=[\ell']\).  The constant restriction \(g|_H\)
therefore supplies a single \(\alpha\in\PGL(H)\) for the
whole map.  Replacing the recovered base map by its inverse if
necessary gives the right-hand side of
\eqref{eq:cover-equivalence-v148}.

Conversely choose linear lifts of \(\alpha\) on \(H\) and
extend by the identity on a complement in \(V\).  Together with
the base map and the identity identification of the trivial source
bundles, this gives equivalent pencil data.  Their homogeneous
normal maps induce an isomorphism of the finite schemes.  This
proves (iii) and the theorem.
\end{proof}

\begin{corollary}[A pencil-native separating pair]
\label{cor:pencil-native-pair-v148}
For every \(n\ge3\), the degree-three maps
\[
          f_0=[s^3:t^3],\qquad f_1=[s^3+st^2:t^3]
\]
give nonisomorphic unmarked order-\(d\) failure schemes of
quadratic pencils, although they have the same Artin fibres, the
same conormal and first-relation bundles, the same pencil and
quotient bundles, and isomorphic nilradical graded vector bundles
in every degree.  For \(n=3\) their common finite rank is
\(167945\), and their top graded bundle is
\[
                  \OO^{75537}\oplus\OO(3)^{\oplus30}.
\]
\end{corollary}
\begin{proof}
The pairs of cubics have no common zero.  In the affine coordinate
\(z=s/t\) the maps are \(z^3\) and \(z^3+z\).
The first has two ramification points, both of index three.
The second has the two distinct zeros of \(3z^2+1\), each of
index two, and infinity of index three.  Independent projective
changes in source and target preserve this multiset of indices,
so the two maps are inequivalent.  The theorem proves all the
geometric assertions.  For \(n=3\),
\(e=9,d=11,M=15\), whence
\(\binom{20}{11}-15=167945\) and
\(\binom{19}{11}-45=75537\).
\end{proof}

This consequence uses the global inverse essentially.  Pointwise
local reconstruction returns the same pencil orbit at every
point for every \(f\); it cannot recover the map \(f\).
Even supplementing it with all the abstract bundles displayed in
the theorem leaves the covering equivalence problem unresolved.
It is the multiplication in the total unmarked thickening,
followed by descent to one constant left transformation, that
recovers that problem.  The double-orbit description is not an
assertion that the scheme has the automorphism group of the
cover: the unused complement of \(H\), the right factor, and
the higher-order kernel remain present, as recorded by the
previous automorphism theorems.

% BEGIN PRESERVED PART 21-spectral-specialization.tex
\section{Spectral torsion and rank-preserving specialization}
\label{sec:spectral-specialization}

The finite-neighbourhood theorem recovers more than the discriminant
or the corank of each member of a regular pencil.  We identify its
classical spectral invariant and then give a flat specialization on
which the entire discriminant and every reduced rank locus remain
constant, while the recovered invariant changes.  The two specializations
are separated at exactly the order of Theorem~\ref{thm:sharp-finite-pencil}.
The parameter space throughout is the usual Grassmannian of pencils;
the failure map is the original multiplication map, not an added
coefficient-evaluation block.

\subsection{The spectral torsion sheaf}

Let \(R\subset\Sym^2V\) be a regular pencil, meaning that it
contains a nonsingular quadratic form.  On its parameter line
\(L_R=\Pj(R)\), the tautological quadratic form gives an exact sequence
\begin{equation}\label{eq:spectral-torsion-sequence}
 0\longrightarrow V^*\otimes\OO_{L_R}(-1)
 \xrightarrow{\ q_R\ }V\otimes\OO_{L_R}
 \longrightarrow\mathcal T_R\longrightarrow0.
\end{equation}
The first arrow is injective because its determinant is nonzero at the
generic point of the smooth integral curve \(L_R\).  Thus
\(\mathcal T_R\) is a torsion sheaf of length \(n\), whose
zeroth Fitting divisor is the full discriminant.  The pair
\((L_R,\mathcal T_R)\) is considered up to isomorphism of the
parameter line.  Its definition involves no choice of a basis of \(R\).
The relation with elementary divisors is classical; see
\cite[\S2, Corollary~2.1 and equation~(5)]{FevolaMandelshtamSturmfels}.
We record the argument to make precise which datum is recovered.

\begin{lemma}[Sheaves on the affine spectral line]
\label{lem:spectral-sheaf-similarity}
Let \(A,B\) be symmetric complex matrices with \(A\) invertible.
On the affine line with coordinate \(z\), the cokernel of
\(zA-B\) is the finite-dimensional \(\C[z]\)-module on
\(\C^n\) for which multiplication by \(z\) is
\(C=A^{-1}B\).  An isomorphism of such sheaves, in this fixed
coordinate, is equivalent to similarity of the corresponding matrices.
\end{lemma}
\begin{proof}
Multiplication on the target by \(A^{-1}\) identifies the
presentation with \(zI-C\).  Reduction of a polynomial vector
modulo these relations sends \(\sum_jz^jv_j\) to
\(\sum_jC^jv_j\); constant vectors give the inverse identification.
The equivalence between coherent sheaves on an affine scheme and
finitely generated modules identifies a finite-length sheaf supported
in this affine line with this module.  The support is finite, so its
underlying complex vector space is finite dimensional.  A module
isomorphism is a linear isomorphism \(G\) satisfying
\(GC=C'G\), which is exactly similarity.  Extension by zero across
a point outside the support loses no sheaf data.
\end{proof}

\begin{lemma}[Polynomial square roots of invertible operators]
\label{lem:artinian-polynomial-square-root}
For every invertible complex matrix \(K\) there is
\(h\in\C[x]\) with \(h(K)^2=K\).  If \(K\) is
self-adjoint for a nondegenerate symmetric form and commutes with
\(C\), then \(h(K)\) has both properties.
\end{lemma}
\begin{proof}
Write the minimal polynomial as
\(\prod_\lambda(x-\lambda)^{a_\lambda}\), where
\(\lambda\ne0\).  In its local Artinian factor choose a square
root \(\mu_\lambda^2=\lambda\), and put
\[
 h_\lambda(x)=\mu_\lambda
 \sum_{j=0}^{a_\lambda-1}\binom{1/2}{j}
          \left(\frac{x-\lambda}{\lambda}\right)^j.
\]
The formal binomial identity, truncated modulo
\((x-\lambda)^{a_\lambda}\), gives \(h_\lambda(x)^2=x\).
The Chinese remainder theorem joins these polynomials to a single
class modulo the minimal polynomial.  Evaluation at \(K\) proves
the assertion even when \(K\) is not semisimple.  Polynomials in
one self-adjoint operator remain self-adjoint, and every polynomial
in \(K\) commutes with every operator commuting with \(K\).
\end{proof}

\begin{theorem}[Spectral reconstruction at the first relation]
\label{thm:spectral-finite-torelli}
Let \(n\ge3\), \(d=n^2+2n-4\), and let \(R,R'\) be
regular pencils.  The following are equivalent:
\begin{enumerate}
\item the unmarked schemes \(Z_R^{[d]}\) and \(Z_{R'}^{[d]}\)
are isomorphic;
\item there is an isomorphism \(f:L_R\to L_{R'}\) with
\(\mathcal T_R\simeq f^*\mathcal T_{R'}\);
\item the pencils \(R,R'\) are projectively equivalent.
\end{enumerate}
In particular the finite scheme determines the full root-labelled
spectral torsion sheaf, not only its Fitting divisor or its support.
This is an equivalence of isomorphism classes of objects; it is not
an assertion identifying their automorphism groups or moduli stacks.
\end{theorem}
\begin{proof}
The equivalence of the first and third conditions is
Theorem~\ref{thm:sharp-finite-pencil}.  A congruence of pencils
induces an isomorphism of \eqref{eq:spectral-torsion-sequence},
so the third implies the second.

For the converse choose a common affine coordinate whose point at
infinity avoids the two spectral supports, using \(f\) to identify
the lines.  Write the corresponding symmetric pairs as \((A,B)\)
and \((A',B')\), with \(A,A'\) invertible, and set
\(C=A^{-1}B\), \(C'=A'^{-1}B'\).  Choose the presentation \(zA-B\) for the affine
coordinate on both lines.  Lemma~\ref{lem:spectral-sheaf-similarity}
identifies the two sheaves with the modules of \(C,C'\), so their
isomorphism gives a similarity.  After applying
that change of coordinates, assume \(C'=C\).  Both \(A\) and
\(A'\) are nondegenerate symmetric forms for which \(C\) is
self-adjoint.  Then \(K=A^{-1}A'\) commutes with \(C\)
and is self-adjoint for \(A\).  Lemma~\ref{lem:artinian-polynomial-square-root}
provides a polynomial \(h\) with \(h(K)^2=K\).  Put \(H=h(K)\).  It commutes with \(C\), is
self-adjoint for \(A\), and satisfies
\[
 H^tAH=A',\qquad H^tBH=H^tACH=A'C=B'.
\]
Thus the pairs, and hence the pencils, are congruent.  This also shows
why the spectral sheaf alone suffices here; no additional pairing on
that sheaf is required over \(\C\).
\end{proof}

For \(x\in\Supp\mathcal T_R\), put
\begin{equation}\label{eq:spectral-length-profile}
 \ell_a(x)=\length\bigl(\mathcal T_{R,x}/
                  \mathfrak m_x^a\mathcal T_{R,x}\bigr),
 \qquad a\ge0.
\end{equation}
If the elementary-divisor exponents at \(x\) are
\(e_1,\ldots,e_b\), the structure theorem over the discrete
valuation ring \(\OO_{L_R,x}\) gives
\begin{equation}\label{eq:spectral-partition-inversion}
 \ell_a(x)=\sum_{j=1}^b\min(a,e_j),\qquad
 \ell_a(x)-\ell_{a-1}(x)=\#\{j:e_j\ge a\}.
\end{equation}
Thus these intrinsic local lengths recover every exponent, including
repeated roots.  The corank of the pencil at \(x\) gives only
\(\ell_1(x)\), and the multiplicity of the discriminant gives
only the eventual value \(\ell_a(x)=\sum_j e_j\).  The
intermediate values are genuinely additional information.

\subsection{A specialization with constant discriminant and coranks}

The following family keeps both ends of this length profile fixed.
Fix \(i\in\C\) with \(i^2=-1\), and set
\begin{equation}\label{eq:rank-preserving-four-block}
 N_0=\begin{pmatrix}1&i\\i&-1\end{pmatrix},\qquad
 P_t=\frac t2N_0,\qquad
 A_0=\begin{pmatrix}0&I_2\\I_2&0\end{pmatrix},\qquad
 C_t=\begin{pmatrix}P_t&I_2\\0&P_t\end{pmatrix}.
\end{equation}
Put
\begin{equation}\label{eq:symmetric-specialization-pair}
 B_t=A_0C_t=\begin{pmatrix}0&P_t\\P_t&I_2\end{pmatrix},
 \qquad R_t=\langle A_0,B_t\rangle\subset\Sym^2\C^4.
\end{equation}
Both matrices are symmetric and \(A_0\) is invertible.  They
are independent for every \(t\).  For \(n>4\), append
\(I_{n-4}\) to \(A_0\) and append
\(\diag(\mu_5,\ldots,\mu_n)\) to \(B_t\), where the
\(\mu_j\) are fixed distinct nonzero numbers.  Denote the resulting
matrices by \(A^{(n)},B_t^{(n)}\) and their pencil again by \(R_t\).

\begin{theorem}[Separation inside fixed reduced rank data]
\label{thm:rank-preserving-specialization}
For every \(n\ge4\), the family \(R_t\), \(t\in\A^1\),
has the following properties.
\begin{enumerate}
\item Its full discriminant, with multiplicities, is independent of \(t\):
\begin{equation}\label{eq:fixed-discriminant-specialization}
 \det(sA^{(n)}+uB_t^{(n)})
       =s^4\prod_{j=5}^n(s+\mu_j u).
\end{equation}
For every minor size, the reduced determinantal locus on the marked
parameter line is also independent of \(t\).
\item At the root \(s=0\), the spectral module has exponents
\((3,1)\) for \(t\ne0\) and \((2,2)\) for \(t=0\).
Consequently \(Z_{R_t}^{[d]}\not\simeq Z_{R_0}^{[d]}\)
for every \(t\ne0\), whereas all the punctured fibres are isomorphic.
\item The schemes \(Z_{R_t}^{[d]}\) form a flat projective family
over \(\A^1\), with fixed reduction.  After one identification
of the socle quotient bundles, every truncation of order \(k<d\)
is a constant family.  Thus the first separating order remains
exactly \(d=n^2+2n-4\) even under the constraints in part~(1).
\end{enumerate}
\end{theorem}
\begin{proof}
The identities \(N_0^2=0\) and \(\rank N_0=1\) give
\begin{equation}\label{eq:rank-preserving-power-identities}
 C_t^2=\begin{pmatrix}0&tN_0\\0&0\end{pmatrix},\qquad
 C_t^3=0,\qquad
 C_t=\binom{I_2}{P_t}\begin{pmatrix}P_t&I_2\end{pmatrix}.
\end{equation}
The right factor in the last expression is surjective and the left
factor is injective.  Hence \(\rank C_t=2\) for every \(t\),
whereas \(\rank C_t^2=1\) precisely when \(t\ne0\).
The nilpotent partitions are therefore \((3,1)\) and \((2,2)\).
Also \(\det A_0=1\) and \(\det(sI_4+uC_t)=s^4\),
proving \eqref{eq:fixed-discriminant-specialization}.  At \(s=0\)
the full pencil has corank two for all \(t\); at each remaining
root it has corank one, and elsewhere corank zero.  Equality of these
functions on \(\Pj^1(\C)\) implies equality of the reduced
closed determinantal subschemes for each minor size.

At the unique multiplicity-four root, formula
\eqref{eq:spectral-partition-inversion} yields respectively
\[
 (\ell_1,\ell_2,\ell_3,\ldots)=(2,3,4,\ldots),
 \qquad (2,4,4,\ldots).
\]
No automorphism of the parameter line can exchange this root with a
simple one.  Theorem~\ref{thm:spectral-finite-torelli} separates
the special fibre from the punctured fibres and identifies the latter
with one another.  This is a statement about geometric fibres; no
choice of a simultaneous trivializing congruence on \(\A^1\setminus0\)
is needed.

For flatness we use the relative gluing of
Theorem~\ref{thm:universal-finite-neighbourhood}; the following
coordinates identify its pullback in the present family.  First note that \(t\mapsto R_t\) is a morphism
to \(\Gr(2,\Sym^2V)\), not merely a collection of pencils.
In symmetric matrix coordinates the two fixed linear functionals
\[
 a(Q)=Q_{13}+iQ_{14},\qquad b(Q)=Q_{33}
\]
satisfy \(a(A^{(n)})=b(B_t^{(n)})=1\) and
\(a(B_t^{(n)})=b(A^{(n)})=0\).  Their common kernel is a
fixed complement to every \(R_t\).  It trivializes
\(S_{R_t}=\Sym^2V/R_t\), and hence identifies the relative
socle Grassmannian with \(B\times\A^1\).

On a frame chart of its tautological bundle, the graph description
\eqref{eq:finite-homogeneous-ideal} identifies the truncated algebra
with
\begin{equation}\label{eq:relative-spectral-finite-algebra}
 \bigoplus_{j=0}^{d-1}\Sym^j E^*
       \ \oplus\ \bigl(\Sym^d E^*/K_{R_t}\bigr).
\end{equation}
By Theorem~\ref{thm:finite-parameter-immersion}, \(K_{R_t}\)
is a subbundle of constant rank \(\binom N2\), with locally
free quotient.  This statement commutes with pullback to the present
base and with right changes of frame, so the local descriptions glue.
Thus \eqref{eq:relative-spectral-finite-algebra} is locally free
on \(B\times\A^1\), of rank
\(\binom{n^2+d}{d}-\binom N2\).  The graph projection is a
finite flat morphism to that base.  Since \(B\) is projective,
its composition with projection to \(\A^1\) is flat and
projective.  Below degree \(d\) there is no relation, and the
fixed complement identifies the entire lower truncated algebra,
including multiplication, with a constant family.  These arguments
prove flatness of the finite neighbourhoods; they make no assertion
about flatness of the unrestricted failure schemes.
\end{proof}

The qualifier ``reduced'' in part~(1) is necessary.  At \(s=0\),
in a local parameter \(z=s/u\), the gcd of the three-row minors of
the four-dimensional block is a unit times \(z\) for \(t\ne0\)
and a unit times \(z^2\) for \(t=0\).  Both radicals are
\((z)\).  Thus the family retains exactly the higher Fitting
information discarded by reduction; it does not contradict the
classical classification by the full determinantal ideals.

\subsection{An independently visible change of reciprocal geometry}

For a regular pencil let its reciprocal curve be the projective closure
of the inverses of its nonsingular members.  These curves and their
strata are studied independently of multiplication-failure schemes in
\cite[\S3]{FevolaMandelshtamSturmfels}.

\begin{corollary}[A conic specializes to a line]
\label{cor:reciprocal-conic-line}
For the four-dimensional family \eqref{eq:symmetric-specialization-pair},
the reciprocal curve is a smooth plane conic when \(t\ne0\)
and a projective line when \(t=0\).  The unmarked twentieth
neighbourhood distinguishes these two cases, while all earlier
neighbourhoods, the full discriminant, and all reduced rank loci agree.
\end{corollary}
\begin{proof}
For \(s\ne0\), multiplication on the right by the fixed
invertible matrix \(A_0^{-1}\) identifies the projective inverse
with
\[
 (sI_4+uC_t)^{-1}
   =s^{-1}I_4-us^{-2}C_t+u^2s^{-3}C_t^2.
\]
If \(t\ne0\), the matrices \(I_4,C_t,C_t^2\) are
linearly independent: their span has dimension equal to the degree
three of the minimal polynomial.  The closure is therefore the
quadratic Veronese image \([s:u]\mapsto[s^2:-su:u^2]\),
a smooth conic in its spanning plane.  At \(t=0\) the last term
vanishes, and \(I_4,C_0\) are independent, giving a line.
The remaining assertions follow from
Theorem~\ref{thm:rank-preserving-specialization}, with \(d=20\).
The degree drop does not assert that the reciprocal curves themselves
form a flat family; the flat family proved above consists of the
finite failure neighbourhoods.
\end{proof}

% BEGIN PRESERVED PART 22-relative-spectral-strata.tex
\section{The universal finite neighbourhood and spectral strata}
\label{sec:relative-spectral-strata}

The finite-order construction is a family of actual neighbourhoods in
multiplication-failure schemes.  This fact is needed when passing from
individual pencils to degenerations and infinitesimal parameter spaces.
We establish it before discussing the spectral incidence schemes.

Put \(P=\Gr(2,\Sym^2V)\), and write \(\mathcal R\) and
\(\mathcal S=(\Sym^2V\otimes\OO_P)/\mathcal R\) for its
universal relation and quotient bundles.  Let
\[
 \pi:\mathcal B=\Gr_P(n,\mathcal S)\longrightarrow P,
 \qquad \mathcal E=\Hom(\mathcal U,V\otimes\OO_{\mathcal B}),
\]
where \(\mathcal U\) is tautological of rank \(n\).  As before,
\(N=n(n+1)/2\), \(p=N-2\), \(d=n+2p\), and
\(M=\binom N2\).

\begin{theorem}[Universal neighbourhood and arbitrary base change]
\label{thm:universal-finite-neighbourhood}
The universal multiplication-failure scheme has a canonical order-\(d\)
neighbourhood \(\mathcal Z^{[d]}\) of \(\mathcal B\).  There is a
rank-\(M\) vector subbundle
\(\mathcal K\subset\Sym^d\mathcal E^*\) such that
\begin{equation}\label{eq:universal-truncated-algebra}
 \mathcal Z^{[d]}=
 \Spec_{\mathcal B}\left(
 \Sym\mathcal E^*/
 (\mathcal K,\Sym^{d+1}\mathcal E^*)\right).
\end{equation}
This is finite locally free over \(\mathcal B\), of rank
\begin{equation}\label{eq:universal-finite-rank}
                   \binom{n^2+d}{d}-\binom N2,
\end{equation}
and is projective and flat over \(P\).  Its construction, its
ideal-adic filtration, and its identification with the actual
multiplication-failure neighbourhood commute with every morphism of
complex schemes \(T\to P\), including nonreduced ones.
For \(k<d\), its order-\(k\) truncation is
\(\Spec_{\mathcal B}\bigoplus_{j=0}^k\Sym^j\mathcal E^*\),
with truncated multiplication.  After trivializing \(\mathcal S\)
on the parameter base, these lower-order families are constant.
\end{theorem}
\begin{proof}
The locus in \(\Gr_P(n,V\oplus\mathcal S)\) on which projection
to \(\mathcal S\) is injective is canonically \(\Tot\mathcal E\).
It is an open neighbourhood of \(\mathcal B\), not a choice of
formal coordinates.  On a frame chart of \(\mathcal U\), let
\(T:\mathcal U\to V\) be the universal graph map and let
\(\gamma:\Sym^2V\to\mathcal S\) be the quotient.  The same
block elimination that proves
\eqref{eq:finite-homogeneous-ideal} works over its coordinate ring:
the ideal of the multiplication-failure scheme is
\[
                   (\det T)I_p(\gamma\Sym^2T).
\]
Every generator is homogeneous of degree \(d\).

Here is also a coordinate-free description of its first relation
bundle.  The coefficient map, followed by determinant multiplication,
is a morphism
\begin{equation}\label{eq:universal-relation-bundle-map}
 \det V^*\otimes\det\mathcal U\otimes\det\mathcal S^*
       \otimes\bigwedge^p\Sym^2\mathcal U
 \longrightarrow \Sym^d\mathcal E^*.
\end{equation}
The inclusion used in this expression is
\(\det\mathcal S^*\hookrightarrow(\bigwedge^p\Sym^2V)^*\).
The determinant factors evaluate \(\det T\), and the remaining
factors evaluate \(\bigwedge^p\Sym^2T\) followed by the quotient.
Thus the image is exactly the span of the displayed maximal minors,
not a separately inserted coefficient module.

On a frame chart this map is the fixed injective linear coefficient
map \(\nu\) of Theorem~\ref{thm:finite-parameter-immersion},
restricted to a line subbundle tensored with the fixed space
\(\bigwedge^p\Sym^2\C^n\).  Its image is a subbundle of rank
\(M\): locally split the line in the source, and split the fixed
complex linear inclusion \(\nu\).  Both splittings remain
splittings over any complex algebra.  Denote its image by
\(\mathcal K\).  A right change of frame of \(\mathcal U\)
changes the graph coordinates and this coefficient map
contragrediently.  The determinant character in
\eqref{eq:universal-relation-bundle-map} is exactly the one arising
from \(\det T\).  Hence these local ideals and their degree-\(d\)
subbundles agree on overlaps.  In particular, they glue as ideals
of the given graph neighbourhood.

Let \(\mathfrak b\) be the ideal of the zero section inside the
multiplication-failure scheme.  In graph coordinates its preimage
is the positive-degree ideal.  Quotienting by
\(\mathfrak b^{d+1}\) therefore gives precisely
\eqref{eq:universal-truncated-algebra}.  This identifies the gluing
with the actual ideal-adic neighbourhood, not merely with algebras
having the correct geometric fibres.  Its underlying module is
\[
 \bigoplus_{j=0}^{d-1}\Sym^j\mathcal E^*
       \ \oplus\ (\Sym^d\mathcal E^*/\mathcal K).
\]
All terms are locally free.  Counting monomials in \(n^2\)
variables gives \eqref{eq:universal-finite-rank}.  The morphism
\(\mathcal B\to P\) is smooth and projective, so the asserted
flatness and projectivity follow.  Finally the graph description,
the split coefficient inclusion, and the homogeneous quotient all
commute with base change.  Below degree \(d\) the ideal contributes
no relation.  This proves the final assertion as well.
\end{proof}

\begin{remark}\label{rem:relative-reduction-scope}
Over a nonreduced base, the theorem concerns the specified relative
socle subscheme \(\mathcal B_T\) and its ideal.  It does not identify
that ideal with the absolute nilradical of \(\mathcal Z_T^{[d]}\),
which can also contain nilpotents from the base.  On geometric fibres
\(\mathcal B_t\) is exactly the reduction.  Nor is this theorem an
equivalence of unmarked deformation groupoids: extra automorphisms
of truncated algebras are not being identified with pencil
congruences.
\end{remark}

\subsection{Spectral incidence schemes from the first relation space}

On the regular-pencil locus, pass locally to ordered generators
\((A,B)\) with \(A\) invertible, and put \(C=A^{-1}B\).
For a marked point with affine coordinate \(z\), set \(X=C-zI\).
The next elementary presentation works over an arbitrary coefficient
ring; no diagonalization or division by eigenvalue differences is
involved.

\begin{lemma}[Truncated spectral presentation]
\label{lem:truncated-spectral-presentation}
Let \(D\) be a commutative ring, \(X\in\operatorname{Mat}_n(D)\),
and \(a\ge1\).  There is a natural isomorphism
\begin{equation}\label{eq:truncated-module-power}
 \coker(wI-X)\otimes_{D[w]}D[w]/(w^a)
                  \simeq \coker(X^a:D^n\to D^n).
\end{equation}
Equivalently, the \(an\)-square block presentation with diagonal
blocks \(X\) and subdiagonal blocks \(-I\) is equivalent, by
invertible row and column operations over \(D\), to
\(I_{n(a-1)}\oplus X^a\).
\end{lemma}
\begin{proof}
Map the class of \(\sum_{j=0}^{a-1}w^jv_j\) to the class of
\(\sum_{j=0}^{a-1}X^jv_j\) modulo \(X^aD^n\).
The relation \(wv=Xv\) and the relation \(w^av=0\) show that
this map is well-defined.  Sending a constant vector to its class
is inverse.  Eliminating successively the blocks with coefficient
\(-I\) gives the stated presentation equivalence, using only unit
pivots.  This also proves naturality under ring homomorphisms.
\end{proof}

\begin{theorem}[Scheme-theoretic spectral readout]
\label{thm:relative-spectral-readout}
In the fixed tensor presentation of
Theorem~\ref{thm:finite-parameter-immersion}, the first relation
subspace determines, on the regular-pencil locus, every closed
spectral incidence scheme
\begin{equation}\label{eq:spectral-incidence-minors}
 \mathfrak S_{a,h}
       =V\bigl(I_{n-h+1}((C-zI)^a)\bigr),
       \qquad 1\le a,h\le n.
\end{equation}
This determination commutes with arbitrary complex base change.
At a geometric point, membership is equivalent to
\(\length(\mathcal T_z/\mathfrak m_z^a\mathcal T_z)\ge h\).
Consequently the schemes of specified spectral length profiles, and
not only their reduced supports, can be read from the order-\(d\)
relation subspace.
\end{theorem}
\begin{proof}
The closed immersion \(\Phi\) of
Theorem~\ref{thm:finite-parameter-immersion} has an inverse on its
scheme-theoretic image.  Thus a family of relation subspaces in that
image recovers its morphism to \(P\), over any base, and hence its
universal pencil.  In the displayed chart its spectral sheaf has
presentation \((z+w)I-C=wI-X\).  Lemma~\ref{lem:truncated-spectral-presentation}
gives the presentation \(X^a\) for restriction to the order-\(a\)
fat point centred at \(z\).  Its \((h-1)\)-st Fitting ideal is
exactly the ideal in \eqref{eq:spectral-incidence-minors}; Fitting
ideals of finitely presented modules commute with base change.
On a field the cokernel has dimension at least \(h\) exactly when
these minors vanish.

These descriptions are independent of the chosen frames.  A change
of basis in \(V\) conjugates \(X\); a change of affine coordinate
replaces the parameter ideal at the marked point by a unit multiple
locally.  Alternatively, restriction of the spectral sheaf to the
fat diagonal on its parameter line defines the same finite module
without a coordinate.  Thus the Fitting ideals glue.  Rank equality
is imposed by the corresponding rank inequality together with the
open complement of the next inequality.  Intersecting these locally
closed conditions gives each specified length profile.  Formula
\eqref{eq:spectral-partition-inversion} then identifies its geometric
fibres with the prescribed elementary-divisor partition.
\end{proof}

\subsection{Realizing the complete Segre stratification}

Let \(P_\sigma\) denote the reduced closure in \(P\) of a
regular Segre stratum, and let \(J_\sigma=\Phi(P_\sigma)\).
Because \(\Phi\) is a closed immersion, it gives isomorphisms
\begin{equation}\label{eq:segre-scheme-realization}
 P_\sigma\simeq J_\sigma,
 \qquad
 P_{\sigma_1}\cap_P\cdots\cap_P P_{\sigma_s}
 \simeq
 J_{\sigma_1}\cap_{\Phi(P)}\cdots\cap_{\Phi(P)}J_{\sigma_s},
\end{equation}
where the intersections carry their scheme structures.  In
particular their local rings and formal completions are preserved.
Pulling back Theorem~\ref{thm:universal-finite-neighbourhood} realizes
these parameter schemes by flat projective families of the actual
finite failure neighbourhoods.  The classical inclusions among
Segre-stratum closures \cite[Theorem~5.1]{FevolaMandelshtamSturmfels}
are therefore all realized, rather than just the one degeneration
of Theorem~\ref{thm:rank-preserving-specialization}.

There are two distinct assertions here.  The Segre closure order is
classical, and transport of closed subschemes through a closed
immersion is formal.  The geometric input proved here is the
base-change-compatible realization by actual multiplication-failure
neighbourhoods and the recovery of their spectral Fitting schemes.
The ambient parameter space in \eqref{eq:segre-scheme-realization}
is the image of \(\Phi\), not a moduli space of every possible
unmarked deformation of a finite scheme.

% BEGIN PRESERVED PART 25-fixed-spectral-strata-v145.tex
\section{Fixed-discriminant strata and all admissible specializations}
\label{sec:fixed-spectral-strata}

The universal family of Section~\ref{sec:relative-spectral-strata}
and the example of Section~\ref{sec:spectral-specialization}
lead to a classification of the entire fixed-discriminant,
fixed-reduced-rank locus. The spectral orbit order is classical.
Every allowed specialization is realized here by a flat family of
the original finite multiplication-failure neighbourhoods, and the
same first relation detects every strict transition.

Fix distinct \(\alpha_1,\ldots,\alpha_s\in\C\), positive
\(m_i\) with \(\sum m_i=n\ge3\), and \(1\le c_i\le m_i\).
Fix a nonsingular symmetric \(A\), and let \(\mathfrak p_A\)
be its self-adjoint endomorphisms. Consider the reduced locally closed
variety
\begin{equation}\label{eq:fixed-spectral-locus}
 X_{\boldsymbol m,\boldsymbol c}(A)=
 \left\{C\in\mathfrak p_A:
 \det(zI-C)=\prod_i(z-\alpha_i)^{m_i},\quad
 \dim\ker(C-\alpha_i I)=c_i\right\}.
\end{equation}
Exclude \(s=1,c_1=n\), for which \(C\) is scalar and
\(\langle A,AC\rangle\) is not a pencil. Every other datum
is realized. Put \(R_C=\langle A,AC\rangle\),
\(d=n^2+2n-4\), and \(Z_C=Z_{R_C}^{[d]}\).
The full binary determinant is
\(\det(A)\prod_i(u+\alpha_i v)^{m_i}\). The coranks and
labelled roots specify all reduced determinantal loci on the line.

Write \(\mathcal P(m,c)\) for partitions into exactly \(c\)
positive parts. For equal-size partitions, \(\mu\preceq\lambda\)
means \(\sum_{j\le a}\mu_j\le\sum_{j\le a}\lambda_j\)
for every \(a\); write \(\lambda'\) for the conjugate partition.

\subsection{The classical orthogonal input}
\label{sec:orthogonal-input-v145}

We isolate the orbit theorem used below, with its precise group and
source.  Let \((H,A)\) be a nondegenerate complex symmetric space of
dimension \(m\).  The acting group is the full orthogonal group
\(O(A)\), and the operators are \(A\)-self-adjoint, not elements
of the skew-adjoint Lie algebra \(\mathfrak o(A)\).

\begin{lemma}[Self-adjoint nilpotent orbits]
\label{lem:orthogonal-orbits-v145}
The full \(O(A)\)-orbits of self-adjoint nilpotent operators are
indexed by all partitions \(\lambda\) of \(m\).  If
\(\mathcal O_\lambda\) denotes the orbit, then
\[
 \dim\mathcal O_\lambda=
       \frac12\left(m^2-\sum_a(\lambda'_a)^2\right),\qquad
 \mathcal O_\mu\subset\overline{\mathcal O_\lambda}
       \ \Longleftrightarrow\ \mu\preceq\lambda.
\]
An \(O(A)\)-orbit is one \(SO(A)\)-orbit when \(\lambda\)
has an odd part.  When every part is even, it is the disjoint union
of two \(SO(A)\)-orbits, interchanged by an isometry of determinant
\(-1\).  The closure statement above is for the full orthogonal
orbit and does not identify the closures of its two components.
\end{lemma}
\begin{proof}
The classification and closure assertions are
\cite[Proposition~1, p.~444; Theorem~1, p.~447]{Ohta1986}, respectively,
with \(\varepsilon=1\) and \(G(V)\) as defined in
\cite[\S1.1, p.~443]{Ohta1986}.  Remark~1 on p.~444 specifies that
all partitions occur; the order in \S1.4 on p.~446 is precisely the
row-dominance order used here.  The proof of Theorem~1 occupies
pp.~447--451.  These references concern the full group, not its
identity component.

For the dimension formula there is a short independent calculation.
The trace pairing on \(\End H\) is nondegenerate on both its
self-adjoint and skew-adjoint summands.  For a self-adjoint nilpotent
\(N\), the commutator with \(N\) exchanges the summands, and the
two resulting maps are adjoint up to sign for that pairing.  They
have equal rank.  Consequently
\(\dim\operatorname{GL}(H)N=2\dim O(A)N\), also the identity
in \cite[\S2.4, Remark~8(i), p.~456]{Ohta1986}.
The general linear centralizer has dimension
\(\sum_{i,j}\min(\lambda_i,\lambda_j)
=\sum_a(\lambda'_a)^2\), as follows by writing homomorphisms
between the cyclic \(\C[t]\)-modules of the Jordan decomposition.
This proves the stated dimension.

We give the component argument explicitly.  The good bases of
\cite[Lemma~1, p.~444]{Ohta1986} decompose the symmetric space
orthogonally into cyclic Jordan blocks, each with a reversed-identity
Gram matrix.  If one block has odd size, its negative identity,
extended by the identity on the other blocks, centralizes \(N\)
and has determinant \(-1\).  The centralizer then meets both
components of \(O(A)\), so the orbit is a single \(SO(A)\)-orbit.

For the converse, write
\[
 E_k=\ker N^k/(\ker N^{k-1}+N\ker N^{k+1}).
\]
The form induced by \(A(u,N^{k-1}v)\) on \(E_k\) is
nondegenerate and symmetric; this is immediate on the same cyclic
blocks.  A commuting isometry induces \(g_k\in O(E_k)\).
A filtration of the Jordan module by its cyclic lengths and powers
of \(t\) gives
\[
                        \det g=\prod_k(\det g_k)^k.
\]
Indeed the multiplicity-space action of a length-\(k\) block occurs
\(k\) times on the associated graded; the strictly triangular
terms of a module automorphism do not change its determinant.
If all lengths are even, this product is one.  The centralizer lies
in \(SO(A)\), and the index-two quotient of \(O(A)\) therefore
gives exactly two \(SO(A)\)-orbits.  Since \(SO(A)\) is connected,
these are the connected components of the full orbit.  This does not
preclude their closures from meeting.
\end{proof}


\begin{theorem}[Classification and specialization at the first relation]
\label{thm:fixed-spectral-classification}
The following statements hold.
\begin{enumerate}
\item The \(O(A)\)-orbits in \(X_{\boldsymbol m,\boldsymbol c}(A)\)
are indexed by \(\boldsymbol\lambda\in\prod_i\mathcal P(m_i,c_i)\),
with dimensions
\begin{equation}\label{eq:fixed-spectral-orbit-dimension}
 \dim\mathcal O_{\boldsymbol\lambda}
 =\frac12\left(n^2-\sum_i\sum_a(\lambda'_{i,a})^2\right).
\end{equation}
Unmarked isomorphism classes of \(Z_C\) are these tuples modulo
the permutation group induced on the roots by projective transformations
preserving their multiplicities and coranks.
\item A specialization from \(\boldsymbol\lambda\) to
\(\boldsymbol\mu\) inside \eqref{eq:fixed-spectral-locus}
exists if and only if
\begin{equation}\label{eq:fixed-spectral-dominance}
                 \mu_i\preceq\lambda_i\quad\hbox{for every }i.
\end{equation}
Every such pair is realized over a smooth pointed algebraic curve,
after shrinking at the special point, with \(A\) fixed and with
constant full determinant and reduced rank loci.
\item Every such pencil family induces a flat projective family of
\(Z_C\). After shrinking the curve, every truncation of order
\(k<d\) is constant. If at least one dominance inequality is
strict, the generic and special \(Z_C\) are nonisomorphic;
their first distinguishing order is \(d\).
\item The recovered spectral module at \(\alpha_i\) has length
profile
\begin{equation}\label{eq:fixed-spectral-lengths}
                 L_{i,a}=\sum_j\min(a,\lambda_{i,j}).
\end{equation}
Every \(L_{i,a}\) can only increase in the stated specialization.
All these inequalities together are equivalent to
\eqref{eq:fixed-spectral-dominance}, giving an intrinsic numerical
criterion for every admissible transition.
\end{enumerate}
\end{theorem}
\begin{proof}
We use Lemma~\ref{lem:orthogonal-orbits-v145}.  Its full-orthogonal
classification and dominance order have exactly the group and
self-adjoint convention needed here.  The comparison for Grassmannian
Segre strata is \cite[Theorem~5.1]{FevolaMandelshtamSturmfels}.

The primary spaces \(V_i=\ker(C-\alpha_iI)^{m_i}\) are an
orthogonal direct sum: the Chinese remainder idempotents are
polynomials in the self-adjoint \(C\), hence self-adjoint projections.
Each restricted form is nondegenerate, and
\(C|_{V_i}=\alpha_iI+N_i\), with \(N_i\) self-adjoint
nilpotent. The prescribed corank is its number of Jordan blocks.
The classical theorem gives the indexing. Conversely, a Jordan block
is self-adjoint for the reversed-identity Gram matrix; an isometry
with the chosen nondegenerate form realizes each partition.

The space of orthogonal decompositions has dimension
\((n^2-\sum_i m_i^2)/2\). Adding the nilpotent orbit dimensions
proves \eqref{eq:fixed-spectral-orbit-dimension}. These are dimensions
in the fixed-\(A\) matrix space, not in an unmarked moduli stack.
Full orthogonal orbits need not be connected; no irreducibility
assertion is needed. Theorem~\ref{thm:spectral-finite-torelli}
identifies the unmarked neighbourhood with the pencil, and its spectral
module determines the partition at each root. A projective change of
the line can therefore do exactly the stated root permutations.
For one or two roots the projective stabilizer may be infinite;
only its finite permutation image is used.

For necessity, the same polynomial idempotents decompose a family
with fixed characteristic polynomial into subbundles of ranks \(m_i\).
The rank of \(N_i^a\) cannot increase at specialization. Since
\(m_i-\rank N_i^a=\sum_j\min(a,\lambda_{i,j})\), this gives
the length inequalities. Summing the first \(a\) columns of
conjugate Young diagrams shows that they are equivalent to dominance:
conjugation reverses the dominance order.

For sufficiency, fix an orthogonal decomposition of \((V,A)\).
The closure criterion puts a chosen representative of \(\mu_i\)
in the closure of the orbit for \(\lambda_i\). Thus the special
tuple lies in the closure of the product of the generic orbits.
Choose an irreducible component of that closure containing the point
and meeting the generic orbit product. Successive general hyperplane
sections through the point, avoiding the closed complement, give an
integral algebraic curve with generic point in that open orbit product.
Normalize it and choose a point over the prescribed point. A normal
complex curve is smooth. Remove the other points outside the generic
orbit product. A constant curve suffices in the equality case.
The resulting block operator \(\bigoplus_i(\alpha_iI+N_i(t))\)
is self-adjoint for the same \(A\). Its determinant is unchanged,
and its coranks agree generically and specially, hence throughout the
shrunk curve. This gives the required specialization. It does not
assert constancy of the full higher Fitting ideals.

For part~(3), pull back Theorem~\ref{thm:universal-finite-neighbourhood}
along the curve's morphism to \(P=\Gr(2,\Sym^2V)\).
That theorem identifies the family with the actual ideal-adic
neighbourhoods and proves flatness and projectivity under arbitrary
base change. No inference from fibrewise isomorphisms is needed.
An affine Grassmannian chart around the special pencil supplies a
fixed complement to nearby \(R_C\), identifying \(S_{R_C}\),
\(B\), and \(E\). There is no relation below degree \(d\),
so all lower neighbourhoods become constant. For a strict dominance
transition, \(\sum_{i,a}(\lambda'_{i,a})^2\) strictly increases;
this sum is invariant under permutations of the decorated roots.
The generic and special tuples therefore cannot be related by such
a permutation. Part~(1) separates their unmarked degree-\(d\)
neighbourhoods. Flatness is asserted only for these finite
neighbourhoods, not unrestricted failure schemes.
\end{proof}

\begin{corollary}[Exactly when the coarse data lose a labelled class]
\label{cor:fixed-spectral-ambiguity}
The number of root-labelled congruence classes with the prescribed
data is \(\prod_i\#\mathcal P(m_i,c_i)\). It is one precisely
when, for every \(i\), either \(c_i=1\) or \(m_i-c_i\le1\).
The first possible ambiguity with fixed full determinant and all
reduced rank loci occurs in dimension four, at \((m,c)=(4,2)\),
with partitions \((3,1)\) and \((2,2)\).
\end{corollary}
\begin{proof}
Subtracting one from every part identifies \(\mathcal P(m,c)\)
with partitions of \(m-c\) into at most \(c\) parts.
There is one if \(c=1\) or \(m-c\le1\). Otherwise
\((m-c)\) and \((m-c-1,1)\) are distinct admissible partitions.
The theorem gives the count and the minimal example.
\end{proof}

\subsection{The first relation as a standard graded Betti space}

The finite invariant has a usual Hilbert-scheme interpretation.
The first Betti space means the generator space in the first free
module of a minimal graded resolution of the quotient, not the space
of relations among generators of its defining ideal. The following
statement retains the marked quotient presentation.

\begin{proposition}[A graded Hilbert locus and its first Betti space]
\label{prop:first-syzygy-hilbert}
In the fixed coordinates of Theorem~\ref{thm:finite-parameter-immersion},
let \(S=\Sym H^*\), \(N=\binom{n+1}{2}\), \(M=\binom N2\),
and \(Q_R=S/((K_R)+S_{\ge d+1})\). Then
\begin{equation}\label{eq:first-relation-tor}
                     \Tor_1^S(Q_R,\C)_d=K_R.
\end{equation}
This identity commutes with arbitrary complex base change.
The functor of graded quotients with Hilbert function
\[
 h_j=\begin{cases}\dim S_j,&0\le j<d,\\
 \dim S_d-M,&j=d,\\0,&j>d\end{cases}
\]
is represented by \(\Gr(M,S_d)\). The pencil locus is a closed
copy of \(\Gr(2,\Sym^2V)\). Its degree-\(d\) first Betti
space recovers every scheme-theoretic pencil parameter, including
parameters over nonreduced bases.
\end{proposition}
\begin{proof}
Put \(J=(K_R)+S_{\ge d+1}\) and \(S_+=\bigoplus_{j>0}S_j\).
The augmented sequence \(0\to J\to S\to Q_R\to0\) gives
\(\Tor_1^S(Q_R,\C)=J/S_+J\), since its map to \(S/S_+\)
is zero. Its degree-\(d\) part is \(K_R\). The same argument
over a base algebra applies to the degree-\(d\) subbundle with
locally free quotient and commutes with pullback.
A graded quotient with the indicated locally free components has
zero ideal below degree \(d\), a rank-\(M\) subbundle in degree
\(d\), and the full ideal in higher degrees. Conversely, every
such subbundle defines an ideal, as its positive-degree products lie
above degree \(d\). This is the Grassmannian functor. The closed
immersion and its inverse on the image are
Theorem~\ref{thm:finite-parameter-immersion}. This is not an inference
from closed-point Torelli to unmarked deformation-groupoid equivalence.
\end{proof}

% BEGIN PRESERVED PART 17-assistance.tex
\section*{Acknowledgment of assistance}
AI assistance was used in developing and drafting the arguments and
finite verification scripts.  All structural proofs are supplied for
independent mathematical scrutiny; computational receipts are not
formal proof certificates.  No AI system is listed as an author.

% BEGIN PRESERVED PART 43-assistance-specific-v151.tex
For the present revision, GPT-6 Astra Pro assisted specifically with the
incidence normalization and overlap criterion in
Section~\ref{sec:canonical-boundary-v151}, the universal-property proof
in Section~\ref{sec:rigidification-v151}, the explicit specialization
in Section~\ref{sec:closed-pencil-v151}, and the associated exact tests
and exposition.  These arguments require independent human verification.
No statement here represents a completed journal review or a formal
proof-assistant certification.

% BEGIN PRESERVED PART 47-assistance-specific-v152.tex
\paragraph{Additional disclosure for the present revision.}
GPT-6 Astra Pro assisted specifically with the conductor and collision
models, the scalar-matrix degeneration, the flag-jet calculation,
the reciprocal-polynomial fibre presentation, literature comparisons,
and exact computational checks.  The proofs above state the complete
mathematical arguments; finite computations are consistency checks,
not universal proof certificates.  No journal decision or historical
priority determination is inferred from the build and source audits.

\section*{Assistance in the present revision}
An AI assistant helped formulate and write the binary conductor stratification,
the simultaneous incidence atlas, the evaluated singular Fitting ideal, and
the minimal-relation-module argument for pure-power fibres. It also prepared
the source-preserving reorganization and finite exact consistency checks.
The written proofs, rather than those finite checks, are the mathematical
basis of the statements and are submitted to independent scrutiny.

\paragraph{Revision 154.}
AI assistance was used for the polarized moduli compatibility,
spectral divisor-fibre construction, homological calculations and
proof expansions. The fixed five-variable Betti table uses exact
computer-assisted rational elimination with its complete finite
inputs recorded. The general theorems rely on the written arguments
and the explicitly attributed classical results, not on regression
counts or source-preservation receipts.

\begin{thebibliography}{99}

\bibitem{Ballico93}
E.~Ballico,
On the failure locus of higher order properties of embeddings in
projective spaces,
\emph{Math. Nachr.} \textbf{163} (1993), 5--13.
doi:10.1002/mana.19931630102.

\bibitem{Ballico96}
E.~Ballico,
On the failure cycles for the quadratic normality of a projective
variety,
\emph{Pacific J. Math.} \textbf{172} (1996), 307--313.

\bibitem{ChengWang}
S.-J. Cheng and W. Wang, Howe duality for Lie superalgebras,
\emph{Compositio Math.} \textbf{128} (2001), 55--94;
arXiv:math/0008093v1, \S2 and Theorem 3.5.

\bibitem{FevolaMandelshtamSturmfels}
C.~Fevola, Y.~Mandelshtam, and B.~Sturmfels,
Pencils of quadrics: old and new,
\emph{Le Matematiche} \textbf{76} (2021), no.~2, 319--335.
doi:10.4418/2021.76.2.2; arXiv:2009.04334v2,
Theorem~1.1, Corollary~2.1, Remark~2.3, Conjecture~4.5, and Theorem~5.1.

\bibitem{Ohta1986}
T.~Ohta, The singularities of the closures of nilpotent orbits in certain
symmetric pairs, \emph{Tohoku Math. J. (2)} \textbf{38} (1986), 441--468.
doi:10.2748/tmj/1178228456.

\bibitem{SmithQuadrics}
I.~Smith, Floer cohomology and pencils of quadrics,
\emph{Invent. Math.} \textbf{189} (2012), 149--250.
doi:10.1007/s00222-011-0364-1; arXiv:1006.1099, \S4.2.

\bibitem{EliasRossiShort}
J.~Elias and M.~E.~Rossi,
Isomorphism classes of short Gorenstein local rings via Macaulay's
inverse system,
\emph{Trans. Amer. Math. Soc.} \textbf{364} (2012), no.~9, 4589--4604.
doi:10.1090/S0002-9947-2012-05430-4;
arXiv:0911.3565v1, Theorem~3.3 and Corollaries~3.4--3.5.

\bibitem{EliasRossiCompressed}
J.~Elias and M.~E.~Rossi,
Analytic isomorphisms of compressed local algebras,
\emph{Proc. Amer. Math. Soc.} \textbf{143} (2015), no.~3, 973--987.
doi:10.1090/S0002-9939-2014-12313-6;
arXiv:1207.6919v1, Theorem~3.1 and Examples~3.3--3.5.

\bibitem{MarcusMoyls}
M.~Marcus and B.~N.~Moyls,
Transformations on tensor product spaces,
\emph{Pacific J. Math.} \textbf{9} (1959), no.~4, 1215--1221.
doi:10.2140/pjm.1959.9.1215, Theorem~1.

\bibitem{StacksGroups}
The Stacks Project Authors,
\emph{The Stacks Project}, Section~39.8, Lemma~39.8.2,
\url{https://stacks.math.columbia.edu/tag/0BF6}, accessed September 24, 2026.
\bibitem{Milne2017}
J.~S. Milne, \emph{Algebraic Groups: The Theory of Group Schemes of
Finite Type over a Field}, Cambridge Studies in Advanced Mathematics
\textbf{170}, Cambridge University Press, Cambridge, 2017.

\bibitem{StacksClosedImmersion}
The Stacks Project Authors, \emph{The Stacks Project},
Lemma~41.7.2 (Tag 04XV),
\url{https://stacks.math.columbia.edu/tag/04XV}, accessed September 24, 2026.

\bibitem{StacksQuotient}
The Stacks Project Authors, \emph{The Stacks Project},
Theorem~97.17.2 and Section~97.18,
\url{https://stacks.math.columbia.edu/tag/06PI}, accessed September 24, 2026.

\bibitem{StacksSmoothQuotient}
The Stacks Project Authors, \emph{The Stacks Project},
Lemma~101.33.7 (Tag 0DLS),
\url{https://stacks.math.columbia.edu/tag/0DLS}, accessed September 24, 2026.

\bibitem{StacksFiniteV151}
The Stacks Project Authors, \emph{The Stacks Project},
Lemma~37.44.1 (Tag 02LS),
\url{https://stacks.math.columbia.edu/tag/02LS}, accessed September 24, 2026.

\bibitem{StacksImageV151}
The Stacks Project Authors, \emph{The Stacks Project},
Section~29.6, Scheme theoretic image (Tag 01R5),
\url{https://stacks.math.columbia.edu/tag/01R5}, accessed September 24, 2026.

\bibitem{StacksNormalV151}
The Stacks Project Authors, \emph{The Stacks Project},
Section~29.55, Normalization (Tag 035E),
\url{https://stacks.math.columbia.edu/tag/035E}, accessed September 24, 2026.

\bibitem{StacksRigidV151}
The Stacks Project Authors, \emph{The Stacks Project},
Subsection~112.5.7, Rigidification (Tag 04V2), including the discussion
of noncentral normal inertia subgroups,
\url{https://stacks.math.columbia.edu/tag/04V2}, accessed September 24, 2026.


\bibitem{Chipalkatti2003}
J.~V. Chipalkatti, On equations defining coincident root loci,
\emph{J. Algebra} \textbf{267} (2003), 246--271.
doi:10.1016/S0021-8693(03)00336-3;
arXiv:math/0110224v1, Proposition~5.1, Theorem~5.4,
and Corollary~5.8.

\bibitem{Kurmann2012}
S.~Kurmann, Some remarks on equations defining coincident root loci,
\emph{J. Algebra} \textbf{352} (2012), 223--231.
doi:10.1016/j.jalgebra.2011.10.045;
arXiv:1108.4532v1, Proposition~1.3, Theorem~1.5,
and Proposition~2.1.

\bibitem{Ferrand2003}
D.~Ferrand, Conducteur, descente et pincement,
\emph{Bull. Soc. Math. France} \textbf{131} (2003), no.~4, 553--585.
doi:10.24033/bsmf.2455.

\bibitem{AOV2008}
D.~Abramovich, M.~Olsson, and A.~Vistoli, Tame stacks in positive
characteristic, \emph{Ann. Inst. Fourier (Grenoble)}
\textbf{58} (2008), no.~4, 1057--1091.
doi:10.5802/aif.2378; Appendix~A, Theorem~A.1.

\bibitem{HuLinShao2007}
Y.~Hu, J.~Lin, and Y.~Shao, A compactification of the space of
algebraic maps from $\Pj^1$ to $\Pj^n$,
arXiv:math/0701255v1 (2007), Theorem~1.1,
\S2.2, and Corollary~3.6.


\bibitem{Grinberg2019}
D.~Grinberg, \emph{A basis for a quotient of symmetric polynomials},
arXiv:1910.00207, version of September 24, 2021, Theorem~2.7.

\bibitem{MFK1994}
D.~Mumford, J.~Fogarty, and F.~Kirwan,
\emph{Geometric Invariant Theory}, third edition,
Ergebnisse der Mathematik und ihrer Grenzgebiete (2), vol.~34,
Springer-Verlag, Berlin, 1994.

\bibitem{MillerSturmfels2005}
E.~Miller and B.~Sturmfels, \emph{Combinatorial Commutative Algebra},
Graduate Texts in Mathematics, vol.~227, Springer, New York, 2005.
Chapter~5 (Stanley--Reisner rings and Hochster's formula).

\bibitem{ShafieiSymmetric}
S.~Shafiei, \emph{Apolarity for determinants and permanents of generic
symmetric matrices}, arXiv:1303.1860v3 (2021), Proposition~3.1,
Theorem~3.10, and Corollary~2.7.

\bibitem{Krishna2026}
A.~Krishna, \emph{A universal discriminant formula for pencils of
quadrics}, arXiv:2607.05340v2 (2026), Theorem~1.1 and Lemma~3.1.

\bibitem{StacksLocalV154}
The Stacks Project Authors, \emph{The Stacks Project},
Tags \href{https://stacks.math.columbia.edu/tag/00R4}{00R4}
(miracle flatness),
\href{https://stacks.math.columbia.edu/tag/07QV}{07QV}
(quasi-excellent rings are Nagata), and
\href{https://stacks.math.columbia.edu/tag/0BJ0}{0BJ0}
(geometrically reduced formal fibres), accessed September 25, 2026.

\end{thebibliography}

\end{document}
